\documentclass[11pt,a4paper]{article}
\usepackage{fontspec}
\setmainfont{Latin Modern Roman}
\usepackage{amsmath,amssymb,amsthm,mathrsfs,bm,mathtools}
\usepackage[a4paper,top=23mm,bottom=23mm,left=23mm,right=23mm]{geometry}
\usepackage{graphicx,booktabs,longtable,array,multirow,xcolor,xurl}
\usepackage[unicode,colorlinks=true,linkcolor=black,citecolor=black,urlcolor=blue!50!black]{hyperref}
\usepackage{bookmark,fancyhdr,enumitem}
\setlength{\headheight}{16pt}
\pagestyle{fancy}\fancyhf{}
\fancyhead[L]{\small Nonconvex exclusion for simple pentagons}
\fancyhead[R]{\small Analytic proof}
\fancyfoot[C]{\thepage}
\renewcommand{\headrulewidth}{0.2pt}
\setlength{\parindent}{1.5em}
\setlength{\parskip}{3pt}
\linespread{1.08}
\setlength{\emergencystretch}{3em}
\setcounter{tocdepth}{1}
\allowdisplaybreaks[1]
\newtheoremstyle{proofmanuscript}{7pt}{5pt}{\normalfont}{}{\bfseries}{.}{\newline}{\thmname{#1}\thmnumber{ #2}\thmnote{ (#3)}}
\theoremstyle{proofmanuscript}
\newtheorem{theorem}{Theorem}[section]
\newtheorem{lemma}[theorem]{Lemma}
\newtheorem{definition}[theorem]{Definition}
\newtheorem{proposition}[theorem]{Proposition}
\newtheorem{corollary}[theorem]{Corollary}
\providecommand{\tightlist}{\setlength{\itemsep}{0pt}\setlength{\parskip}{0pt}}
\title{\textbf{Nonconvex Exclusion for Simple Pentagons of Fixed Area}\\[6pt]\large Fourth English analytic edition}
\author{}
\date{12 October 2026}
\begin{document}
\maketitle
\begin{abstract}
Every nonconvex simple polygon of fixed area with at most five active sides admits a strictly better same-area competitor for the first Dirichlet eigenvalue. Continuous Steiner deformation and ear reflection reduce the problem to candidates with their actual Euclidean section geometry. Small heights are excluded by a triangular container. Middle heights are treated by head and tail responses at the same physical interface and energy. A Volterra equation constructs the regular singular-tail solutions, while direct differential comparisons prove the tail and cotangent response bounds. In this edition, one convex lower bound handles the complete short-tail-dominant branch, and one completed square proves positivity of the tail denominator without an order split. A differential argument for a positive rescaling of the head gap replaces the first terminal curve's degree-seven squared-remainder argument; its original payment and subsequent curvature estimate are preserved. At large heights, a positive comparison function glued across the physical interface proves both noncrossing spectral bounds directly on the full form domain. The crossing bound remains uniform for every real order between two and three. The appendix connects the result to existence of minimizers. Its final regular-pentagon corollary explicitly assumes a separate convex-pentagon extremal theorem.
\end{abstract}
\clearpage
\tableofcontents
\clearpage


\section{The main theorem and its scope}

\begin{theorem}[Nonconvex exclusion for simple pentagons of fixed area]\label{thm-main-nonconvex-pentagon}

Every nonconvex simple polygon \(P\in\mathscr P_{\le5}(A)\) of area \(A>0\) admits a simple polygon \(Q\) of the same area with at most five active sides such that \(\lambda_1(Q)<\lambda_1(P)\). If the admissible class requires exactly five active sides, \(Q\) can also be chosen with exactly five active sides. Consequently, any global minimizer in either fixed-area class, if it exists, is convex. The admissible classes and spectral functional are those of \ref{def-domain}.

\end{theorem}


\section{Admissible domains and critical variations}

\begin{definition}[Admissible domains, criticality, and the spectral functional]\label{def-domain}

Let \(\mathscr P_{\le5}(A)\) denote the bounded open domains of area \(A>0\) whose boundaries are simple polygonal Jordan curves with at most five active sides. Active sides are obtained by merging consecutive collinear sides; self-intersections, retraced edges, and cracks are excluded. A reflex angle has interior angle in \((\pi,2\pi)\), and a strictly convex angle has interior angle in \((0,\pi)\). A domain is called area-critical if every infinitesimal vertex displacement with zero first variation of area also has zero first variation of the first eigenvalue. Set

\[
\lambda_1(P)=\inf_{u\in H_0^1(P)\setminus\{0\}}\frac{\int_P|\nabla u|^2}{\int_Pu^2},\qquad F(P)=|P|\lambda_1(P).
\]

A strict competitor is a simple polygon \(Q\) with at most five sides, the same area, and \(\lambda_1(Q)<\lambda_1(P)\). An ear \(ABC\) has boundary sides \(AB,BC\), an open diagonal \(AC\) contained in the domain, and triangular interior contained in \(P\).

\end{definition}

\begin{lemma}[A pentagon strictly below the square threshold]\label{lem-square}

In the notation of \ref{def-domain}, for every \(A>0\) there is a convex polygon \(Q\) with exactly five active sides, \(|Q|=A\), and \(F(Q)<2\pi^2\).

\end{lemma}

\begin{proof}[Proof of \ref{lem-square}]

Use the spectral functional in \ref{def-domain}. Let \(S=(0,1)^2\), choose \(0<\varepsilon\le1/32\), and remove \(\{x,y\ge0:x+y\le\varepsilon\}\) from the corner at the origin. The resulting open domain \(Q_\varepsilon\) is a convex pentagon, and \[
|Q_\varepsilon|=1-\varepsilon^2/2.
\] The positive normalized ground state of the square is \(U=2\sin(\pi x)\sin(\pi y)\). Define \[
\eta_\varepsilon(r)=\min\{1,\max\{0,r/\varepsilon-1\}\},\qquad
r=(x^2+y^2)^{1/2}.
\] This function vanishes on the cut, so \(\eta_\varepsilon U\in H_0^1(Q_\varepsilon)\). Testing the weak eigenvalue equation on the square with \(\eta_\varepsilon^2U\) gives the exact identity \[
\int_S|\nabla(\eta_\varepsilon U)|^2
=2\pi^2\int_S\eta_\varepsilon^2U^2
+\int_SU^2|\nabla\eta_\varepsilon|^2.
\] In the first quadrant, \(0\le U\le2\pi^2xy\le\pi^2r^2\). Estimate over the quarter disk of radius \(2\varepsilon\) and put \(C_0=16\pi^5/3\). Then \[
1-\int_S\eta_\varepsilon^2U^2\le C_0\varepsilon^6,
\qquad
\int_SU^2|\nabla\eta_\varepsilon|^2\le C_0\varepsilon^4.
\] These estimates use only \(|\nabla\eta_\varepsilon|\le\varepsilon^{-1}\) and \(\int_0^{\pi/2}\int_0^{2\varepsilon}\pi^4r^5\,dr\,d\theta =C_0\varepsilon^6\). Since \(\pi<4\) and \(\varepsilon\le1/32\), \(C_0\varepsilon^6<1/2\), and \[
\frac{C_0\varepsilon^4}{1-C_0\varepsilon^6}
<2C_0\varepsilon^4
\le\frac{\pi^5}{96}\varepsilon^2
<\frac{2\pi^2}{3}\varepsilon^2.
\] The Rayleigh principle therefore gives \[
F(Q_\varepsilon)
<\left(1-\frac{\varepsilon^2}{2}\right)
 \left(2\pi^2+\frac{2\pi^2}{3}\varepsilon^2\right)
=2\pi^2-\frac{\pi^2}{3}\varepsilon^2
 -\frac{\pi^2}{3}\varepsilon^4
<2\pi^2.
\] For example, one may fix \(\varepsilon=1/32\). Finally, take \(t=(A/|Q_\varepsilon|)^{1/2}\). The planar scaling law preserves \(F\) and gives \(|tQ_\varepsilon|=A\).

\end{proof}

\begin{lemma}[Flux moments and the energy tensor at an area-critical polygon]\label{lem-variations}

Let \(P\) be an area-critical simple polygon as in \ref{def-domain}, with no flat angles. Let \(U\) be its positive \(L^2\)-normalized ground state and \(\lambda=\lambda_1(P)\). On each side \(x_i+t(x_{i+1}-x_i)\), put \(g_i(t)=(\partial_\nu U)^2\). Then

\[
\int_0^1g_i(t)dt=\lambda/|P|,\quad \int_0^1(t-1/2)g_i(t)dt=0,\quad \int_P\nabla U\otimes\nabla U=(\lambda/2)I.\tag{G.2.1}
\]

All the flux integrals are finite.

\end{lemma}

\begin{proof}[Proof of \ref{lem-variations}]

Use the admissible class, area-criticality, normalized Dirichlet ground state, and spectral functional of \ref{def-domain}.

The corner Hadamard formula \ref{lem-hadamard-corners} applies to simple polygons without flat angles and to boundary fields obtained by affine interpolation of vertex velocities. To construct such an admissible field \(V\), take a finite nondegenerate triangulation in a large rectangle containing \(P\), compatible with \(\partial P\). Choose a generic direction so that all vertex projections are distinct. Join the unique leftmost and rightmost vertices to the corresponding sides of the rectangle. Except at their initial endpoints, these two disjoint segments lie outside \(P\) and split the surrounding annular region into two simple polygons. Triangulate these polygons and \(P\), using identical subdivisions on shared sides. Assign the prescribed velocity \(V_i\) to each original vertex; at any additional mesh point on an old side, use the linear interpolation of its endpoint velocities. Set all other mesh velocities, and the field outside the rectangle, to zero. The resulting \(V\) is continuous, piecewise affine, and globally Lipschitz. Its restriction to every old side is affine, so \(I+tV\) creates no additional boundary vertices. If \(|t|\operatorname{Lip}(V)<1\), the distance estimate and the contraction mapping construction of inverse images give a global bi-Lipschitz map. The image domain remains simple. Positive edge lengths and nonflat angles persist under sufficiently small perturbations, so the number of active sides is unchanged.

The vertex velocity space is \(\mathbb R^{2m}\). The Hadamard formula makes \(D\lambda\) linear on this space, and the vertex area formula gives linearity of the area differential as well. Uniform dilation gives \(D|P|=2|P|>0\), so the area differential is nonzero. Criticality means \(\ker D|P|\subset\ker D\lambda\). Choose \(W\) with \(D|P|[W]=1\). For every \(V\), one has \(V-D|P|[V]W\in\ker D|P|\), whence \(D\lambda[V]=D\lambda[W]D|P|[V]\). Thus a single multiplier \(\kappa=D\lambda[W]\) exists.

Write \(D\lambda=\kappa D|P|\). Moving only \(x_i\), with velocity \(V_i\), gives

\[
D\lambda[V_i]=-|e_{i-1}|V_i\cdot\nu_{i-1}
\int_0^1t g_{i-1}(t)\,dt
-|e_i|V_i\cdot\nu_i\int_0^1(1-t)g_i(t)\,dt,
\]

\[
D|P|[V_i]=\tfrac12V_i\cdot
(|e_{i-1}|\nu_{i-1}+|e_i|\nu_i).
\]

The two adjacent normals are linearly independent. Hence both weighted moments equal \(-\kappa/2\). Uniform dilation gives \(-2\lambda=2\kappa|P|\), so \(\kappa=-\lambda/|P|\). Adding and subtracting the two endpoint moments proves the flux identities.

The absolute integrability of the edge fluxes near all corners and the disappearance of the artificial circular boundary terms follow from \ref{lem-hadamard-corners}.

The tensor identity follows from \ref{lem-energy-tensor}, whose proof uses only area-criticality and the Rayleigh comparison in both parameter directions.

\end{proof}

\begin{lemma}[G.2.2 Isotropy of the energy tensor]\label{lem-energy-tensor}

Let \(P\) be an area-critical simple polygon in the sense of \ref{def-domain}. Let \(U\) be its positive \(L^2\)-normalized first Dirichlet eigenfunction and \(\lambda=\lambda_1(P)\). Then \[
\int_P\nabla U\otimes\nabla U=\lambda I/2.
\]

\end{lemma}

\begin{proof}[Proof of \ref{lem-energy-tensor}]

Use area-criticality and the Rayleigh definition in \ref{def-domain}, and put \(E=\int_P\nabla U\otimes\nabla U\). For any symmetric trace-free matrix \(M\), let \(B_t=\exp(tM)\) and \(P_t=B_tP\). Since \(\det B_t=1\), this is an area-preserving deformation through simple polygons. Criticality gives \[
\lambda_1(P_t)=\lambda+o(t)\qquad(t\to0).
\] Use the same ground state as an affine trial function, \(U_t=U\circ B_t^{-1}\in H_0^1(P_t)\). Its mass remains one, and \(B_t^{-1}B_t^{-T}=\exp(-2tM)\). Hence \[
\lambda_1(P_t)
\le\int_P\nabla U^{\mathsf T}\exp(-2tM)\nabla U
=\lambda-2t(M:E)+O(t^2).
\] For \(t>0\), divide by \(t\) and let \(t\downarrow0\) to obtain \(M:E\le0\). For \(t<0\), divide by \(t\) and let \(t\uparrow0\) to obtain \(M:E\ge0\). Thus \(M:E=0\) for every symmetric trace-free \(M\), so \(E\) is a scalar matrix. Since \(\operatorname{tr}E=\int_P|\nabla U|^2=\lambda\), \[
E=\frac\lambda2 I.\tag{G.2.2}
\] This argument does not differentiate the deformed ground state or use criticality of any one-dimensional comparison operator.

\end{proof}

\begin{lemma}[Corner control and the vertex Hadamard formula at convex and reflex angles]\label{lem-hadamard-corners}

Let \(P\) be a bounded simple polygon, with all interior angles \(\omega_i\in(0,2\pi)\). Let \(U\in H_0^1(P)\) be its positive first Dirichlet eigenfunction, normalized by \(\int_PU^2=1\), and put \(\lambda=\lambda_1(P)\). Let \(V\) be a continuous piecewise affine vector field on a finite compatible triangulation, extended outside as a Lipschitz field. For small \(t\), set \(F_t=I+tV\). Then

\[
D\lambda[V]=-\int_{\partial P}(\partial_\nu U)^2V\cdot\nu\,ds.
\]

Every edge integral is absolutely convergent. The conclusion covers both strictly convex and reflex corners.

\end{lemma}

\begin{proof}[Proof of \ref{lem-hadamard-corners}]

Use the spectral definition in \ref{def-domain}.

First justify differentiation and the corner estimates. If \(\|tDV\|_\infty<1\), then \(F_t\) is bi-Lipschitz and preserves simplicity of the boundary. Pulling functions back to \(P\) gives the energy and mass forms

\[
a_t(u,u)=\int_P\det(DF_t)|DF_t^{-T}\nabla u|^2,
\qquad m_t(u,u)=\int_P\det(DF_t)u^2.
\]

The first variation follows directly from the Rayleigh principle. Write \(\lambda_t=\lambda_1(F_tP)\). All forms have the fixed domain \(H_0^1(P)\). Since \(DV\) is bounded, their coefficients admit the uniform expansions \[
a_t=a_0+t\dot a_0+O(t^2),\qquad
m_t=m_0+t\dot m_0+O(t^2).
\] Here the energy remainder, evaluated on \((u,u)\), is bounded by \(Ct^2\|\nabla u\|_2^2\), and the mass remainder by \(Ct^2\|u\|_2^2\), with \(C\) independent of \(u\) and small \(t\). Indeed, \[
\det(DF_t)=1+t\operatorname{div}V+O(t^2),
\] \[
\det(DF_t)DF_t^{-1}DF_t^{-T}
=I+t\bigl((\operatorname{div}V)I-DV-DV^{\mathsf T}\bigr)+O(t^2)
\] almost everywhere, with bounded remainders. Comparing each energy and mass form with its value at \(t=0\), uniformly over all nonzero trial functions, gives \(\lambda_t\to\lambda\).

For each small \(t\), choose the nonnegative minimizing eigenfunction \(u_t\) on the fixed domain, normalized by the same condition \(m_0(u_t,u_t)=1\). Coercivity and Rellich compactness give such a minimizer. The identity \[
a_t(u_t,u_t)=\lambda_t m_t(u_t,u_t)
\] and the coefficient comparison show that \(u_t\) is bounded in \(H_0^1(P)\) and that \(a_0(u_t,u_t)\to\lambda\). Every sequence \(t\to0\) therefore has a subsequence converging weakly in \(H_0^1(P)\) and strongly in \(L^2(P)\) to a nonnegative normalized function \(u_*\). Weak lower semicontinuity and the Rayleigh principle give \[
\lambda\le a_0(u_*,u_*)\le\liminf_{t\to0}a_0(u_t,u_t)=\lambda.
\] Thus \(u_*\) is a ground state. Positivity and simplicity on the connected domain \(P\) imply \(u_*=U\). The convergence of the energy norms then improves the weak convergence to strong convergence in \(H_0^1(P)\). Applying this argument to every sequence \(t\to0\) shows that \(u_t\to U\) strongly as \(t\to0\).

Put \(D=\dot a_0(U,U)-\lambda\dot m_0(U,U)\). The fixed trial function \(U\) gives the upper bound \[
\lambda_t-\lambda\le tD+O(t^2).
\] For the lower bound, use the minimizing function \(u_t\) and expand its quotient: \[
\lambda_t-\lambda
=\frac{(a_0-\lambda m_0)(u_t,u_t)}{m_t(u_t,u_t)}
+t\frac{(\dot a_0-\lambda\dot m_0)(u_t,u_t)}{m_t(u_t,u_t)}
+O(t^2)
\ge tD+o(|t|).
\] The first numerator is nonnegative by the Rayleigh principle at \(t=0\); strong convergence makes the second quotient tend to \(D\). Both bounds hold for positive and negative \(t\), so \(\lambda_t=\lambda+tD+o(|t|)\). Evaluating the coefficient derivatives gives

\[
D\lambda[V]=\int_P\{(\operatorname{div}V)(|\nabla U|^2-\lambda U^2)-2\nabla U\cdot DV\nabla U\}.
\]

Near each vertex, the domain is exactly a sector \(0<r<r_0,0<\theta<\omega\). Put \(\alpha=\pi/\omega>1/2\) and choose \(r_0\) so that \(\lambda r_0^2<\min(1/8,\alpha^2/2)\) and the sector meets no other vertex. Odd reflection of \(U\) across each open straight side preserves the weak Helmholtz equation because the trace is zero. Standard interior elliptic regularity therefore makes the circular trace \(F(\theta)=U(r_0,\theta)\) smooth up to both endpoints, where it vanishes. Write its sine coefficients as \(c_n\); they are bounded and satisfy \(\sum n^2|c_n|^2<\infty\). For \(\nu=n\alpha\), define

\[
B_\nu(r)=r^\nu\sum_{j\ge0}\frac{(-\lambda r^2/4)^j}{j!(\nu+1)\cdots(\nu+j)},
\]

with the empty product equal to one. The series converges absolutely and solves the radial equation term by term. On \(0\le r\le r_0\), the alternating terms decrease in magnitude, giving

\[
\frac56\le B_\nu(r)/r^\nu\le1,\qquad
|B_\nu'(r)|\le C(\nu+1)r^{\nu-1},
\]

where \(C\) is independent of \(n\). The derivative estimate follows directly by differentiating the absolutely convergent series: every denominator factor \(\nu+j\ge3/2\), and the additional factor \(\nu+2j\) is uniformly controlled by the associated exponential series. The regular solution is therefore

\[
U(r,\theta)=\sum_{n\ge1}c_n\frac{B_{n\alpha}(r)}{B_{n\alpha}(r_0)}\sin(n\alpha\theta)
\]

To identify this series with the original weak solution, note that \(\sum n^2|c_n|^2<\infty\) ensures convergence in the energy space of the whole sector, with the same boundary trace. The energy of each term is bounded by \(C(n\alpha+1)|c_n|^2\), and its radial equation holds termwise. The angular Poincaré inequality gives the zero-trace spectral bound \(\lambda_D\ge\alpha^2/r_0^2>\lambda\), so the weak solution with this boundary trace is unique. The singular radial branch \(r^{-n\alpha}\) is not in \(H^1\) and supplies no additional term. On \(r\le r_0/2\), the series and its first derivatives converge uniformly and absolutely because \(\sum_{n\ge1}(n+1)(r/r_0)^{n\alpha}\) is bounded by a constant multiple of \((r/r_0)^\alpha\). Consequently,

\[
|U|\le Cr^\alpha,\qquad |\nabla U|\le Cr^{\alpha-1}.
\]

This covers convex corners \(\alpha>1\), flat angles \(\alpha=1\), and reflex corners \(1/2<\alpha<1\); it controls the differentiated expansion as well as the function itself.

Set \(T=(|\nabla U|^2-\lambda U^2)I-2\nabla U\otimes\nabla U\). Away from the vertices, direct calculation using \(-\Delta U=\lambda U\) gives \(\operatorname{div}T=0\). Remove a sector cap of radius \(\varepsilon\) at each vertex and apply the divergence theorem to \(TV\). On every original open side, \(U=0\) and its tangential derivative vanishes, so \(T\nu=-(\partial_\nu U)^2\nu\). Each artificial circular arc contributes in absolute value at most

\[
C\|V\|_\infty\varepsilon(\varepsilon^{2\alpha-2}+\lambda\varepsilon^{2\alpha})
=O(\varepsilon^{2\alpha-1})\longrightarrow0.
\]

The original edge flux near the vertex is also integrable, since \(\int_0^{r_0}r^{2\alpha-2}dr<\infty\). The volume integral converges because \(U\in H^1\) and \(DV\in L^\infty\). Letting \(\varepsilon\to0\) proves the formula. For piecewise affine \(V\), one may integrate cell by cell; terms on shared mesh edges cancel by continuity of \(V\). Equivalently, one may apply the weak divergence theorem for Lipschitz fields.

\end{proof}


\section{Geometric reduction to one reflex angle}

\begin{lemma}[G.5.1 Triangulation from the unique reflex vertex]\label{lem-reflex-fan}

Use \ref{def-domain}. If \(ABCDE\) is a simple pentagon whose unique reflex vertex is \(C\), then it admits the decomposition \(P=\triangle ABC\cup\triangle ACE\cup\triangle CDE\) with disjoint interiors.

\end{lemma}

\begin{proof}[Proof of \ref{lem-reflex-fan}]

Use the definitions of reflex vertex and ear in \ref{def-domain}. The external geometric inputs are nondegenerate triangulation of a simple polygon and diagonal flips in a convex quadrilateral.

Suppose \(C\) is the unique reflex vertex and the boundary order is \(A B C D E\). We claim

\[
P=\triangle ABC\cup\triangle ACE\cup\triangle CDE
\quad\text{(disjoint interiors)}.
\tag{G.5.1}
\]

\textbf{The triangulation argument.} A simple pentagon admits a nondegenerate triangulation. Its two noncrossing interior diagonals must have a common endpoint. The reflex vertex \(C\) cannot be an ear vertex belonging to only one triangle.

Here are the structural details. Two pentagon diagonals without a common endpoint have alternating endpoints on the boundary. An interior diagonal is a crosscut of the Jordan domain and separates it into two sides. Alternating endpoints lie on opposite sides, so the other interior connecting arc must cross this diagonal. Therefore two noncrossing triangulation diagonals have a common endpoint, and the triangulation is a fan from that vertex to its two nonadjacent vertices. If \(C\) belonged to a single triangle, the angle of that triangle at \(C\) would fill the entire local interior angle of the polygon. A nondegenerate triangle has angle less than \(\pi\), contradicting reflexivity. In a fan centered at \(B\) or \(D\), the vertex \(C\) belongs to exactly one triangle. Both centers are therefore excluded, leaving only \(A,C,E\).

The only possible fan centers are \(A,C,E\). If the fan is centered at \(C\), the claim already holds. If it is centered at \(A\), the diagonals are \(AC,AD\). In the quadrilateral \(ACDE\), the angles at \(A,D,E\) do not exceed the corresponding strictly convex polygon angles, and the angle at \(C\) is the angle of triangle \(ACD\). Thus this quadrilateral is strictly convex. Flipping \(AD\) to \(CE\) gives (G.5.1). The case of a fan centered at \(E\) is symmetric. In particular, the argument does not assume that the convex hull is a quadrilateral.

\end{proof}

\begin{lemma}[Strict cancellation of a reflex and a convex angle at one height]\label{lem-steiner}

Suppose a simple pentagon in \ref{def-domain} has the form \(P=\{(x,y):y_-<y<y_+,a(y)<x<b(y)\}\) in some Euclidean coordinates, where \(a,b\) are continuous and piecewise affine, the interior width is strictly positive, and there are no horizontal boundary segments at interior heights. At one interior height, suppose there are exactly two nonadjacent vertices, one on each boundary chain, one strictly reflex and the other strictly convex; no other vertex has this height. Then \(P\) has a strict competitor, which may be chosen with exactly five active sides.

\end{lemma}

\begin{proof}[Proof of \ref{lem-steiner}]

Use \ref{def-domain} and set \[
\begin{gathered}
m(y)=\frac{a(y)+b(y)}2,\qquad r(y)=\frac{b(y)-a(y)}2,\\
P_q=\{(x,y):y_-<y<y_+,\ qm(y)-r(y)<x<qm(y)+r(y)\}.
\end{gathered}
\tag{G.3.3}
\] For \(0\le q\le1\), every section keeps the same width, so area is preserved.

\textbf{Cancellation parameter and side count.} Let \(y_0\) be the specified height and let \(J_a,J_b\) be the slope jumps of the endpoint functions. The counterclockwise boundary travels down the left chain and up the right chain, so the signs of the turning determinants are \(J_a\) and \(-J_b\), respectively. One convex and one reflex angle imply \(J_aJ_b>0\), and in particular \(J_a+J_b\ne0\). After deformation the jumps are exactly \[
J_{a_q}=\frac{J_a-J_b+q(J_a+J_b)}2,
\qquad
J_{b_q}=\frac{-J_a+J_b+q(J_a+J_b)}2.
\tag{G.3.4}
\] Consequently, \[
q_*=\frac{|J_a-J_b|}{|J_a+J_b|}\in[0,1)
\tag{G.3.5}
\] cancels at least one jump. The top and bottom extrema account for at least two original vertices, while \(y_0\) accounts for two more. Thus there is at most one further interior height at which a slope changes. Each interior height contributes at most two vertices to \(P_q\). A general parameter \(q\) therefore gives at most six sides, and at \(q_*\) at least one side disappears, leaving at most five. If the top or bottom is a horizontal side, it uses one additional original vertex and correspondingly reduces the number of possible interior breakpoints; the same count holds. Positive interior width and continuity of the endpoints ensure that \(P_q\) remains a simple connected polygon.

\textbf{Spectral monotonicity.} Brock \cite{nc:brock2000} proves that continuous Steiner rearrangement preserves the zero trace and does not increase Dirichlet energy for a nonnegative \(H_0^1\) function on an arbitrary open domain. Equimeasurability preserves its \(L^2\) norm. The interval formula multiplies each section center by \(e^{-t}\). Therefore, if \(0<q_2\le q_1\le1\), choosing \(t=\log(q_1/q_2)\) rearranges \(P_{q_1}\) exactly into \(P_{q_2}\). Applying the theorem and the Rayleigh principle to the positive ground state of \(P_{q_1}\) gives \[
\lambda_1(P_{q_2})\le\lambda_1(P_{q_1}).
\] The case \(q_2=0\) is complete Steiner rearrangement and follows from the complete-rearrangement endpoint of the same theorem. The function being rearranged is the ground state itself. Every domain here is a bounded open polygon, and its ground state has a zero extension in \(H^1(\mathbb R^2)\), so all function-space hypotheses hold.

\textbf{Strictness of complete rearrangement.} Apply Capriani \cite{nc:Capriani}, exchanging the coordinate names so that the rearrangement direction is our \(x\), and taking \(n=2,k=1\) and \(f(\xi)=|\xi|^2\). Condition (2.13) holds because the domain is bounded and its vertical projection is a connected interval. For (2.14), the polygon has finite perimeter. Over the interior of the projection, an open side with \(\nu_x=0\) would be horizontal, and such sides are excluded by hypothesis. The finitely many vertices have zero one-dimensional measure. Any horizontal top or bottom side lies at an endpoint of the projection, outside the interior cylinder appearing in the theorem.

It remains to check the transverse critical-set condition (2.12). The positive ground state \(U\) is real analytic in the connected domain \(P_1\): the function \(e^{\sqrt\lambda z}U(x,y)\) is weakly harmonic in three dimensions, so interior regularity and the Poisson formula give convergent local power series. The derivative \(\partial_xU\) cannot vanish identically. Otherwise \(U\) would be constant on almost every horizontal section, whereas the one-dimensional slices of its \(H_0^1\) zero extension vanish at both endpoints. This would force \(U=0\), contradicting normalization. The zero set of a nonzero real-analytic function has measure zero: at each zero, select a nonzero derivative of the lowest order and remove one differentiation. The point lies on a regular zero level set of the resulting derivative. The implicit function theorem and a countable coordinate cover place the full zero set in a countable union of smooth curves. Hence \(\{\partial_xU=0\}\) already has zero planar measure and satisfies the subset condition required in (2.12).

If \(\lambda_1(P_0)=\lambda_1(P_1)\), the inequalities \[
\lambda_1(P_0)
\le\frac{\int|\nabla U^\star|^2}{\int|U^\star|^2}
\le\frac{\int|\nabla U|^2}{\int|U|^2}
=\lambda_1(P_1)
\] force equality of the rearranged Dirichlet energies. The equality theorem then says that \(U\) and \(U^\star\) differ by a single fixed horizontal translation. Positivity of the ground states gives the same relation between their positivity sets. Continuity of the section endpoints upgrades the almost everywhere relation to every height, so \(m\) must be constant. But \[
m'(y_0+)-m'(y_0-)=\frac{J_a+J_b}{2}\ne0,
\] and therefore \(\lambda_1(P_0)<\lambda_1(P_1)\).

\textbf{Strictness at finite parameters.} For every finite real \(q\), the map \[
F_q(x,y)=(x+(q-1)m(y),y)\tag{G.3.6}
\] is an area-preserving bi-Lipschitz map from \(P_1\) to \(P_q\). Its derivative is \(\left(\begin{smallmatrix}1&(q-1)m'(y)\\0&1\end{smallmatrix}\right)\). The pulled-back energy coefficient \(DF_q^{-1}DF_q^{-T}\) is quadratic in \(q\) and is uniformly bounded and positive definite on every compact parameter interval. The common closed form domain is \(H_0^1(P_1)\), and the mass form is fixed. Rellich compactness gives compact resolvent; connectedness gives an isolated simple first eigenvalue. Kato \cite{nc:kato} therefore makes \(q\mapsto\lambda_1(P_q)\) real analytic. If \(0<q_*<1\) and the endpoint were isospectral with \(q=1\), monotonicity would force constancy on \([q_*,1]\). The identity theorem would then force constancy on \([0,1]\), contradicting the strict inequality at complete rearrangement. For \(q_*=0\), that strict inequality applies directly. Thus \(P_{q_*}\) is the required strict competitor.

If it has fewer than five active sides, apply \ref{lem-five-edge-approximation} and choose the approximation sufficiently close to preserve the strict spectral gap. This gives a competitor with exactly five sides.

\end{proof}

\begin{lemma}[All configurations with two reflex vertices]\label{lem-double}

Every simple polygon in \ref{def-domain} with five active sides and two strictly reflex vertices has a strict competitor.

\end{lemma}

\begin{proof}[Proof of \ref{lem-double}]

Use all the definitions and notation of \ref{def-domain}.

Every direction used in \ref{lem-steiner} is a direction in the original Euclidean plane. The affine coordinates below are used only to compute determinant signs. The sum of the pentagon's interior angles is \(3\pi\), so there are at most two strictly reflex vertices. If there are exactly two, neither is an extreme point of the convex hull; the remaining three vertices form its triangular convex hull. Along the three boundary chains between successive hull vertices, the reflex vertices either lie on two different chains or on one common chain.

The oblique or affine coordinates below are used only to determine directions, collinearity, and the ordering of heights. The Steiner operation is always performed in the original Euclidean metric; a general affine transformation is not treated as preserving the Dirichlet spectrum.

\textbf{G.4.1 Separated configuration}

Let the convex hull vertices be

\[
A=(0,0),\quad C=(1,0),\quad E=(0,1),
\]

and the two interior reflex vertices be \(B=(b,c_0),D=(d_0,e)\). Then

\[
b,c_0,d_0,e>0,\quad b+c_0<1,\quad d_0+e<1.
\]

The boundary order is \(A B C D E\). Strict convexity at \(C\) gives

\[
e(1-b)>c_0(1-d_0).
\tag{G.4.1}
\]

For the direction \(BE\), use \(\det(E-B,X-B)\) as transverse height. Then

\[
H(A)=b>0,\quad H(C)=b+c_0-1<0,\quad H(B)=H(E)=0,
\]

\[
H(D)=b(1-e)-d_0(1-c_0).
\]

If

\[
b(1-e)-d_0(1-c_0)<0,
\tag{G.4.2}
\]

both boundary chains are monotone in height and each section is a single interval. The equal-height pair \(B,E\) consists of a reflex and a convex vertex, so the cancellation argument G.3 applies.

For the direction \(AD\), the analogous sufficient condition is

\[
d_0c_0-eb<0
\tag{G.4.3}
\]

It gives single-interval sections and the required convex-reflex pair \(A,D\).

If both (G.4.2) and (G.4.3) fail, then

\[
e\le\frac{c_0d_0}{b},\qquad
 e\le1-\frac{d_0(1-c_0)}b.
\]

The first inequality together with (G.4.1) gives \(d_0>b\). The second together with (G.4.1) gives

\[
0<1-\frac{d_0(1-c_0)}b-\frac{c_0(1-d_0)}{1-b}
=\frac{(1-b-c_0)(b-d_0)}{b(1-b)},
\]

so \(d_0<b\), a contradiction. At least one direction is available, strictly excluding the separated configuration.

\textbf{G.4.2 Configuration on one chain}

Let the hull vertices be \(A=(0,0),D=(1,0),E=(0,1)\) and the reflex vertices be \(B=(b,c_0),C=(d_0,e)\), again in boundary order \(A B C D E\). Reflexivity means

\[
be<c_0d_0,
\qquad e(1-b)>c_0(1-d_0).
\tag{G.4.4}
\]

These inequalities give \(d_0>b\). Along the direction \(BE\),

\[
H(A)=b>0,\quad H(D)=b+c_0-1<0,
\]

\[
H(C)=b(1-e)-d_0(1-c_0)
<\frac{(1-b-c_0)(b-d_0)}{1-b}<0.
\tag{G.4.5}
\]

The sections are therefore single intervals, with equal-height reflex-convex pair \(B,E\). The cancellation argument applies again.

\textbf{Conclusion G.4.} Every simple pentagon with two reflex vertices has a same-area strict competitor with at most five sides.

For completeness, check the convex hull and section hypotheses. A strictly reflex vertex cannot lie on the boundary of the convex hull. A supporting line through such a vertex would put the whole polygon in a closed half-plane, but its local interior sector has angle greater than \(\pi\) and cannot lie in a half-plane. Thus both reflex vertices lie strictly inside the triangular hull, justifying \(b,c_0,d_0,e>0\) and the two strict inequalities for the sums of coordinates.

In the separated configuration, along \(BE\),

\[
H(D)-H(C)=(1-c_0)(1-d_0)-be>0,
\]

since \(1-c_0>b\) and \(1-d_0>e\). Under (G.4.2), the full ordering is

\[
H(C)<H(D)<H(B)=H(E)=0<H(A).
\]

The boundary chains from the lowest vertex \(C\) to the highest vertex \(A\) are \(CBA\) and \(CDEA\). Every edge of each chain strictly increases in height, and \(B,E\) lie on different chains.

For the separated configuration along \(AD\), take \(H(X)=\det(D,X)\). Then \(H(C)=-e\), \(H(B)=d_0c_0-eb\), \(H(A)=H(D)=0\), and \(H(E)=d_0\). Also \(H(B)-H(C)=d_0c_0+e(1-b)>0\). Under (G.4.3), the full ordering is

\[
H(C)<H(B)<H(A)=H(D)<H(E),
\]

with strictly increasing chains \(CBAE\) and \(CDE\). The vertices \(A,D\) lie on opposite chains.

For the configuration on one chain, along \(BE\),

\[
H(C)-H(D)=(1-c_0)(1-d_0)-be>0.
\]

Together with (G.4.5), this gives \(H(D)<H(C)<H(B)=H(E)=0<H(A)\), with strictly increasing chains \(DCBA\) and \(DEA\).

In each case, the two strictly increasing polygonal chains are continuous piecewise affine graphs of height. Simplicity makes them disjoint except at their common top and bottom endpoints. Hence their interior width is strictly positive and every interior section is a single interval. Adjacent vertices have distinct heights, so no interior horizontal boundary segment occurs. The specified equal-height vertices lie on opposite chains, and every other vertex has a different height. All section hypotheses of \ref{lem-steiner} are therefore satisfied.

\end{proof}

\begin{lemma}[Signed reduction for a unique reflex vertex]\label{lem-single}

Let \(ABCDE\) be a simple pentagon as in \ref{def-domain} whose unique reflex vertex is \(C\). It admits the triangulation \(ABC,ACE,CDE\) with disjoint interiors. If it has no strict competitor, then, after translating \(C\) to \(C=0\), it has the representation

\[
P=(u,-su-\beta v,0,-\gamma u-tv,v),\quad u,v\text{ linearly independent},\quad s,t\ge0,\quad\beta,\gamma>0,\quad\beta\gamma>st.
\]

Moreover, \(|P|=|\det(u,v)|(1+\beta+\gamma)/2\).

\end{lemma}

\begin{proof}[Proof of \ref{lem-single}]

Use the definitions and notation in \ref{def-domain}, including ears and active sides, and the strict cancellation result \ref{lem-steiner}.

By \ref{lem-reflex-fan}, the triangles \(ABC,ACE,CDE\) form a triangulation with disjoint interiors. Introduce oblique coordinates only on this fixed triangulation.

\textbf{G.5.2 Oblique coefficients need not all be positive}

Let \(C=0,A=u,E=v\), where \(u,v\) are linearly independent. By (G.5.1), there exist

\[
\beta,\gamma>0,\qquad\alpha,\delta\in\mathbb R
\]

such that

\[
B=\alpha u-\beta v,
\qquad D=-\gamma u+\delta v.
\tag{G.5.2}
\]

Only the coefficients \(\beta,\gamma\), corresponding to the ear areas, are required to be positive. The coefficients \(\alpha,\delta\) may have either sign; a pentagon with one reflex vertex may have a triangular convex hull.

If \(\delta>0\), the transverse heights for the direction \(CA\) are

\[
A,C:0,\quad B:-\beta,\quad D:\delta,\quad E:1.
\]

The boundary chains are monotone in height, and \(A,C\) form an equal-height convex-reflex pair on opposite chains. The cancellation argument G.3 applies. If \(\alpha>0\), the symmetric argument in direction \(CE\) applies.

A candidate not directly excluded by this cancellation must therefore satisfy

\[
\alpha\le0,\qquad\delta\le0.
\]

Writing

\[
s=-\alpha\ge0,\quad t=-\delta\ge0,
\]

gives the residual family

\[
\boxed{P=(u,-su-\beta v,0,-\gamma u-tv,v),}
\tag{G.5.3}
\]

\[
\boxed{s,t\ge0,\quad\beta,\gamma>0,
\quad\epsilon:=\beta\gamma-st>0.}
\tag{G.5.4}
\]

The last inequality is exactly the determinant condition that \(C\) be strictly reflex. Set \(K_0=|\det(u,v)|/2\). The three triangle areas are \(\beta K_0,K_0,\gamma K_0\), and hence

\[
|P|=K_0(1+\beta+\gamma).
\tag{G.5.5}
\]

This parameter family covers both triangular and quadrilateral convex hulls.

Here is a full verification of the cancellation hypotheses for \(\delta>0\). One may take \(H(X)=\det(u,X)/\det(u,v)\) as height. It is the signed Euclidean height perpendicular to \(u\), multiplied by a nonzero constant, so only the height scale and possibly its orientation change. No general affine spectral transformation is made. The vertex heights are \(H(B)=-\beta<0\), \(H(A)=H(C)=0\), \(H(D)=\delta>0\), and \(H(E)=1\).

If \(0<\delta<1\), the lowest vertex is \(B\), the highest is \(E\), and the strictly increasing chains are \(BAE\) and \(BCDE\). If \(\delta>1\), the highest vertex is \(D\), and the chains are \(BC D\) and \(BAED\), again strictly increasing. If \(\delta=1\), the top is the horizontal side \(DE\), and the two chains from the lowest vertex are \(BAE\) and \(BCD\). The horizontal side lies only at the upper endpoint of the projection and is not an interior horizontal boundary segment.

In every case, each chain is a continuous piecewise affine graph of height, with no interior horizontal side. Simplicity of the original boundary prevents the chains from meeting at an interior height. Their width is therefore strictly positive and each interior section is a single interval. The nonadjacent vertices \(A,C\) lie on opposite chains; they alone have interior height zero and are respectively strictly convex and strictly reflex. Thus every hypothesis of \ref{lem-steiner} holds.

For \(\alpha>0\), simply exchange the labels \(A\leftrightarrow E\), \(B\leftrightarrow D\), \(u\leftrightarrow v\), \(\alpha\leftrightarrow\delta\), \(\beta\leftrightarrow\gamma\), and read the boundary in reverse. The three cases correspond to \(0<\alpha<1\), \(\alpha=1\), and \(\alpha>1\). The lowest vertex becomes \(D\) and the equal-height pair becomes \(C,E\). In the equality case, the upper horizontal side is \(AB\). Thus the symmetric cancellation satisfies all hypotheses as well.

\end{proof}

\begin{lemma}[Four strict ear-length restrictions at a critical candidate]\label{lem-ear}

Let \(ABCDE\) be an area-critical pentagon as in \ref{def-domain} with the triangulation \(ABC,ACE,CDE\) having disjoint interiors. Then

\[
|AB|,|BC|<|AC|,\qquad |CD|,|DE|<|CE|.
\]

More precisely, if the ear \(ABC\) satisfies \(|AB|\ge|AC|\), and \(t\) parametrizes the side from \(B\) to \(C\), then \(\int_0^1(t-1/2)(\partial_\nu U)^2dt>0\).

\end{lemma}

\begin{proof}[Proof of \ref{lem-ear}]

Use the definitions and notation of \ref{def-domain}, in particular its ground state and criticality. The strict moment established below contradicts the zero-moment identity in \ref{lem-variations}. After a rigid motion, write

\[
B=(0,0),\quad C=(\ell,0),\quad A=(a_0,b_0),\quad b_0>0.
\]

The length condition is equivalent to \(a_0\ge\ell/2\). Put \(T_-=\triangle ABC\cap\{x<\ell/2\}\) and \(\mathcal R(x,y)=(\ell-x,y)\). Triangle geometry gives \(\mathcal R(T_-)\subset\triangle ABC\).

Set \(W=U\circ\mathcal R-U\). On \(T_-\),

\[
(-\Delta-\lambda)W=0.
\]

On the base and reflection axis, \(W=0\). The remaining boundary lies on \(AB\), where the original function vanishes and the reflected point lies either inside the ear or on the interior diagonal \(AC\). Thus the boundary values are nonnegative and not identically zero. Since \(T_-\subsetneq P\), strict domain monotonicity gives \(\lambda_1(T_-)>\lambda_1(P)\), and the maximum principle gives \(W>0\). The inward normal to the base points in the positive \(y\) direction. At each open base point \(0<x<\ell\), the function \(U\) is positive on the \(y>0\) side and zero at \(y=0\). Regularity along the open side gives the normal derivative and hence \(U_y(x,0)=\lim_{y\downarrow0}U(x,y)/y\ge0\). Applying Hopf's lemma to \(W>0\) gives

\[
0< W_y(x,0)=U_y(\ell-x,0)-U_y(x,0)\quad(0<x<\ell/2).
\]

Both normal derivatives are nonnegative and the reflected derivative is strictly larger, so their squares have the same strict ordering. The outer normal is \(-e_y\), and the squared flux is \(g=U_y^2\). Thus

\[
(\partial_\nu U)^2(\ell-x)>(\partial_\nu U)^2(x)
\quad(0<x<\ell/2).
\]

Consequently,

\[
\int_0^\ell(x-\ell/2)g(x)\,dx
=\int_0^{\ell/2}(\ell/2-x)[g(\ell-x)-g(x)]\,dx>0.
\]

This contradicts the central-moment condition (G.2.1).

Exchanging the two endpoints similarly excludes \(|BC|\ge|AC|\). Repeating the argument for the other ear gives the necessary conditions at every area-critical candidate,

\[
\boxed{|AB|,|BC|<|AC|,\qquad |CD|,|DE|<|CE|.}
\tag{G.6.2}
\]

These conditions genuinely exclude minimizers: on the finite-dimensional manifold of vertex coordinates with fixed area, a noncritical point has a first-order descent direction tangent to the area constraint. Local area normalization produces an actual same-area competitor. Differentiability follows from \ref{lem-hadamard-corners}.

For completeness, verify every geometric inclusion used in the reflection argument. Let \(y=b_0\tau\), \(0<\tau<1\). The open horizontal section of the ear triangle is

\[
(L(\tau),R(\tau))=(a_0\tau,\ell+(a_0-\ell)\tau),\qquad
L(\tau)+R(\tau)-\ell=(2a_0-\ell)\tau\ge0.
\]

If \(T_-\) has a nonempty section at this height, then \(L(\tau)<\ell/2\) and \(R(\tau)\ge\ell-L(\tau)>\ell/2\). The section of \(T_-\) is therefore \((L(\tau),\ell/2)\), and the reflected section is \((\ell/2,\ell-L(\tau))\subset(L(\tau),R(\tau))\). This proves \(\mathcal R(T_-)\subset\triangle ABC\) at every height.

The boundary other than the base and axis is the open portion of \(AB\) with \(0<\tau<\min(1,\ell/(2a_0))\). There the reflected horizontal coordinate is \(\ell-L(\tau)>\ell/2>L(\tau)\). If \(a_0>\ell/2\), then \(\ell-L(\tau)<R(\tau)\), so the reflected point lies strictly inside the ear. If \(a_0=\ell/2\), then \(\ell-L(\tau)=R(\tau)\), and the reflected point lies on the open interior diagonal \(AC\). In both cases it is an interior point of \(P\), so \(U\circ\mathcal R>0\), whereas \(U=0\) on \(AB\). This proves that the boundary data of \(W\) are nonnegative and not identically zero.

To exclude the second ear side explicitly, relabel \((A',B',C')=(C,B,A)\). The base is now \(B'C'=BA\) and the interior diagonal is \(A'C'=CA\). The assumption \(|BC|\ge|AC|\) is precisely \(|A'B'|\ge|A'C'|\). The same proof contradicts the central zero-moment identity of \ref{lem-variations} on \(BA\). The analogous two relabelings of ear \(CDE\) give the other two strict inequalities.

\end{proof}

\begin{lemma}[Area-normalized approximation with exactly five active sides]\label{lem-five-edge-approximation}

Use \ref{def-domain}. Every simple polygon \(Q\) of area \(A>0\) and active side count \(m\in\{3,4\}\) admits a sequence of simple polygons \(Q_\varepsilon\) of area exactly \(A\) and exactly five active sides such that \(\lambda_1(Q_\varepsilon)\to\lambda_1(Q)\). Thus every previously obtained strict same-area spectral gap is preserved by sufficiently close such approximations.

\end{lemma}

\begin{proof}[Proof of \ref{lem-five-edge-approximation}]

Use the active-side and spectral definitions in \ref{def-domain}, together with the external triangulation theorem for simple polygons.

To replace a competitor \(Q\) with fewer than five active sides, suppose its side count is \(m=3\) or \(4\). Choose midpoints on \(5-m\) different open sides and disjoint neighborhoods that meet no other side. Move each selected midpoint a distance \(\varepsilon>0\) along its side's outer normal, leaving all original vertices fixed. Each chosen side is replaced by two positive-length noncollinear sides. The old nonflat angles persist for small perturbations. The disjoint local neighborhoods ensure simplicity, so the resulting polygon has exactly five active sides.

To check spectral continuity, take a finite nondegenerate triangulation of a rectangle containing \(Q\), including all original vertices and selected midpoints. Construct a continuous piecewise affine field \(V\) that takes the appropriate outer normal at each selected midpoint and is zero at the original polygon vertices. Additional boundary mesh points take the linearly interpolated boundary values, while all remaining mesh vertices and the exterior of the rectangle have value zero. The field \(V\) is globally Lipschitz, and \(F_\varepsilon=I+\varepsilon V\) is bi-Lipschitz whenever \(\varepsilon\|DV\|_\infty<1\). Its image is exactly the polygon described above. After pulling back the Rayleigh quotient, the Jacobian and energy matrix converge uniformly to \(1,I\). The Rayleigh comparison in both directions gives \(\lambda_1(F_\varepsilon Q)\to\lambda_1(Q)\), and the area tends to \(|Q|\). The scaling factor restoring the area tends to one, so the existing strict spectral gap is preserved.

Here are the precise comparison constants and function spaces. Set \(J_\varepsilon=\det DF_\varepsilon\) and \(B_\varepsilon=J_\varepsilon DF_\varepsilon^{-1}DF_\varepsilon^{-T}\). On the finite mesh, \(DF_\varepsilon=I+\varepsilon DV\) almost everywhere. Hence there is \(\delta_\varepsilon\downarrow0\) such that \(1-\delta_\varepsilon\le J_\varepsilon\le1+\delta_\varepsilon\) and \((1-\delta_\varepsilon)I\le B_\varepsilon\le(1+\delta_\varepsilon)I\). Choose \(\varepsilon\) small enough that \(\delta_\varepsilon<1\). The Sobolev chain rule and change of variables hold for the bi-Lipschitz map: first apply them on each triangle to smooth compactly supported functions, then extend by \(H^1\) density. Pullback and inverse pullback are bounded on \(H^1\). A compactly supported function remains compactly supported in the corresponding domain, where convolution in a neighborhood of the compact support gives smooth \(H^1\) approximations. Thus the two maps carry the respective \(H_0^1\) spaces into one another and are inverse bijections. For nonzero \(u\in H_0^1(Q)\), set \(u_\varepsilon=u\circ F_\varepsilon^{-1}\). Then

\[
\frac{1-\delta_\varepsilon}{1+\delta_\varepsilon}
\frac{\int_Q|\nabla u|^2}{\int_Qu^2}
\le \frac{\int_{F_\varepsilon Q}|\nabla u_\varepsilon|^2}{\int_{F_\varepsilon Q}u_\varepsilon^2}
\le \frac{1+\delta_\varepsilon}{1-\delta_\varepsilon}
\frac{\int_Q|\nabla u|^2}{\int_Qu^2}.
\]

Taking the infimum over nonzero functions, using the bijection, gives the same two-sided eigenvalue bounds and hence \(\lambda_1(F_\varepsilon Q)\to\lambda_1(Q)\). Also \(|F_\varepsilon Q|=\int_QJ_\varepsilon\to A\). Put \(s_\varepsilon=(A/|F_\varepsilon Q|)^{1/2}\) and \(Q_\varepsilon=s_\varepsilon F_\varepsilon Q\). In two dimensions, \(x=sy\) preserves Dirichlet energy and multiplies mass by \(s^2\), with a bijection between the corresponding \(H_0^1\) spaces. Thus \(\lambda_1(s\Omega)=s^{-2}\lambda_1(\Omega)\). It follows that \(|Q_\varepsilon|=A\) and \(\lambda_1(Q_\varepsilon)=(|F_\varepsilon Q|/A)\lambda_1(F_\varepsilon Q)\to\lambda_1(Q)\). Positive scaling preserves simplicity and the number of active sides. Hence every fixed strict spectral gap survives for sufficiently small \(\varepsilon\).

The triangulation above is taken compatible with \(\partial Q\): every boundary segment between consecutive original vertices and selected midpoints is a union of mesh edges. Its existence follows from the permitted triangulation theorem for simple polygons. Choose a generic direction in which the vertex projections are distinct, giving unique leftmost and rightmost vertices, and take a sufficiently large rectangle containing \(Q\) and aligned with that direction. Join these extreme vertices to the left and right rectangle boundaries. The two segments are disjoint and, apart from their initial endpoints, lie outside \(Q\). They split the annular region between the rectangle and \(Q\) into two simple polygons. Triangulate these polygons and \(Q\), then subdivide adjacent triangles at the selected side midpoints. This yields a finite nondegenerate compatible mesh.

On each prescribed boundary segment, the values of \(V\) are interpolated linearly from its endpoints. Any further mesh subdivision uses that same interpolation, rather than assigning zero at a newly inserted boundary point. Consequently, \(F_\varepsilon\) is affine on each entire prescribed boundary segment and creates no additional corners. Its boundary image is precisely the designated pentagon, so the side count and Rayleigh continuity comparison apply.

\end{proof}


\section{Actual sections and candidate parameters}

\begin{lemma}[G.7 Spectral lower bound from the longest connected section]\label{lem-largest-section}

Let \(\Omega\subset\mathbb R^2\) be a bounded Lipschitz open set whose vertical projection is contained in a finite open interval \(I\). Let \(\ell(y)\) be the length of the longest connected component of the horizontal open section \(\Omega_y\), setting \(\ell=0\) on an empty section, and put \(V(y)=\pi^2/\ell(y)^2\), with value \(+\infty\) on empty sections. In integrals, use the convention \(Vv^2=0\) when \(V=+\infty,v=0\). For \(u\in H_0^1(\Omega)\), take its zero extension and define \(v(y)=\|u(\cdot,y)\|_{L^2(\mathbb R)}\). Then \(\ell,V\) are measurable, \(v\in H_0^1(I)\), and

\[
\int_I v^2=\int_\Omega u^2,\qquad
\int_I |v'|^2\le\int_\Omega |u_y|^2,\qquad
\int_I Vv^2\le\int_\Omega |u_x|^2.
\]

Consequently, each nonzero \(u\) produces a nonzero \(v\) in the one-dimensional form domain, and

\[
\lambda_1(\Omega)\ge\inf_{0\ne v\in H_0^1(I),\ \int_I Vv^2<\infty}
\frac{\int_I(|v'|^2+Vv^2)}{\int_Iv^2}.\tag{G.7.4}
\]

\end{lemma}

\begin{proof}[Proof of \ref{lem-largest-section}]

First check measurability of the section length. Each \(\Omega_y\) is a bounded open set and is the union of at most countably many disjoint open intervals \(J\). Their total length has a uniform finite upper bound. If the section is nonempty, its positive supremum of component lengths is attained: if that supremum is \(M>0\), only finitely many components have length at least \(M/2\), and their maximum is \(M\). For \(r>0\), the condition \(\ell(y)>r\) is equivalent to the existence of rational numbers \(a<b\) with \(b-a>r\) and \([a,b]\times\{y\}\subset\Omega\). For each such rational pair, the set of heights satisfying this inclusion is open, by compactness of the segment and openness of \(\Omega\). Thus \(\{y:\ell(y)>r\}\) is a countable union of open sets. Hence \(\ell\) is lower semicontinuous and \(V\) is measurable as an extended-valued function.

Let \(\widetilde u\) be the zero extension. By the definition of \(H_0^1\), one has \(\widetilde u\in H^1(\mathbb R^2)\). Regard \(L^2(\mathbb R_x)\) as the target Hilbert space. Fubini's theorem and the definition of weak derivatives show that \(w(y)=\widetilde u(\cdot,y)\) belongs to \(H^1(\mathbb R_y;L^2(\mathbb R_x))\), with weak derivative \(\widetilde u_y(\cdot,y)\). On every finite interval of heights, take its absolutely continuous representative. The fundamental theorem for the Bochner integral gives

\[
|\|w(t)\|_2-\|w(s)\|_2|\le\|w(t)-w(s)\|_2
\le\int_s^t\|w'(y)\|_2\,dy\quad(s<t).
\]

Therefore \(v=\|w\|_2\) is absolutely continuous, and the Lebesgue differentiation theorem gives \(|v'|\le\|w'\|_2\) almost everywhere. This also handles points where \(v=0\) without division by \(v\). Fubini gives \(\|v\|_2=\|u\|_2\), and hence \(v\in H^1(\mathbb R)\). Since \(w=0\) almost everywhere outside \(I\), its norm has a continuous representative vanishing at both endpoints of \(I\). Thus \(v|_I\in H_0^1(I)\).

For the zero trace on each component, choose \(u_n\in C_c^\infty(\Omega)\) whose zero extensions converge to \(\widetilde u\) in \(H^1(\mathbb R^2)\), and pass to a subsequence with summable squared errors. Fubini then gives \(u_n(\cdot,y)\to\widetilde u(\cdot,y)\) in \(H^1(\mathbb R_x)\) for almost every \(y\). For every such \(y\) and every component \(J\), one has \(u_n(\cdot,y)|_J\in C_c^\infty(J)\): the support stays a positive distance from \(\partial\Omega\) and therefore avoids both endpoints of \(J\). Restriction is continuous in \(H^1\), so closedness of \(H_0^1(J)\) gives \(\widetilde u(\cdot,y)|_J\in H_0^1(J)\). The same subsequence works for every component.

Apply the Dirichlet Poincaré inequality on each interval, then sum the nonnegative terms and use Tonelli's theorem. This yields

\[
\int_\mathbb R |\widetilde u_x(x,y)|^2\,dx
\ge\sum_J\frac{\pi^2}{|J|^2}\int_J\widetilde u^2
\ge\frac{\pi^2}{\ell(y)^2}v(y)^2.
\]

On an empty section, \(v=0\), and the right-hand side is zero by convention. Integration gives a finite upper bound for the potential energy, so the constructed \(v\) lies in the stated form domain. The function \(v\) is nonzero whenever \(u\) is nonzero. Add the vertical derivative bound, compare Rayleigh quotients, and take the infimum over \(u\) to obtain the spectral lower bound.

\end{proof}

\begin{lemma}[G.8 Comparison for every weight from planar criticality]\label{lem-weighted-section}

Use the spectral and profile notation of \ref{def-domain} and \ref{def-candidate}. Suppose an area-critical pentagon is expressed in normalized coordinates with physical height \(\rho y\), area \(S\rho/2\), and longest connected section width \(f(y)\), \(0<y<L\). Then, for every \(A>0\), one has \(\mu(A)\le\rho^2\lambda_1(P)/2+A\lambda_1(P)/(2\pi^2)\). If \(F(P)\le2\pi^2\), then \(\mu(A)\le2(\pi^2\rho+A/\rho)/S\) and \(\rho\ge S/(2L^2)\).

\end{lemma}

\begin{proof}[Proof of \ref{lem-weighted-section}]

Use the ground state, area-criticality, and functional \(F(P)=|P|\lambda_1(P)\) from \ref{def-domain}. Use the quadratic-form construction of \(\mu(A)\) in \ref{def-candidate}, without assuming its spectral necessary condition. Set \(\lambda=\lambda_1(P)\) and let \(U\) be the positive \(L^2\)-normalized ground state. By \ref{lem-energy-tensor}, the energies in the two coordinate directions both equal \(\lambda/2\).

For physical height \(Y\in(0,\rho L)\), let \(w(Y)=\|U(\cdot,Y)\|_{L^2(\mathbb R)}\). By hypothesis, the physical longest component has width \(\ell(Y)=f(Y/\rho)\). The full slicing conclusion of \ref{lem-largest-section} gives \(w\in H_0^1(0,\rho L)\), \(\int w^2=1\), and the bounds \(\int|w'|^2\le\lambda/2\) and \(\pi^2\int f(Y/\rho)^{-2}w(Y)^2dY\le\lambda/2\). Set \(v(y)=\sqrt\rho w(\rho y)\). Changing variables yields

\[
 v\in H_0^1(0,L),\qquad\int_0^L v^2=1,\qquad
\int_0^L|v'|^2\le\rho^2\lambda/2,\qquad
\pi^2\int_0^L f^{-2}v^2\le\lambda/2.
\]

The last bound places \(v\) in the Friedrichs form domain for every \(A>0\). The Rayleigh principle therefore gives

\[
\mu(A)\le\rho^2\lambda/2+A\lambda/(2\pi^2).
\]

Since the area is \(S\rho/2\), the assumption \(F(P)\le2\pi^2\) implies \(\lambda\le4\pi^2/(S\rho)\). Substitution gives \(\mu(A)\le2(\pi^2\rho+A/\rho)/S\). Finally, the vertical Dirichlet Poincaré inequality gives

\[
\pi^2/L^2\le\int|v'|^2\le\rho^2\lambda/2\le2\pi^2\rho/S,
\]

which is \(\rho\ge S/(2L^2)\).

\end{proof}

\begin{lemma}[G.9 A triangular container]\label{lem-triangle-container}

Consider the signed family in \ref{lem-single} and assume the four ear-length bounds in \ref{lem-ear}. Put \(h=\beta+s\ge k=\gamma+t>0\), \(L=1+h\), \(p=\beta/L\), \(q=\gamma/(1+k)\), \(S=1+\beta+\gamma\), and \(d=1-p-q>0\). There is a proper triangular container \(T\supsetneq P\) with \(|P|/|T|=S\min(d,L^{-1})\). If \(F(P)\le2\pi^2\), then \(S\min(d,L^{-1})<\sqrt3/2\).

\end{lemma}

\begin{proof}[Proof of \ref{lem-triangle-container}]

Use the spectral functional in \ref{def-domain} and the signed representation in \ref{lem-single}. Its nonnegative parameters are \(s,t\), so \(h=\beta+s\), \(k=\gamma+t\), and \(\beta\gamma>st\). The four strict ear-length inequalities in \ref{lem-ear} are the prescribed geometric hypotheses.

By a Euclidean similarity alone, write the vertices as \[
A=(0,0),\quad E=(1,0),\quad C=(\xi,\rho),\quad
B=(L\xi-\beta,L\rho),\quad D=(J\xi-t,J\rho),
\] where \(J=1+k\le L=1+h\), \(\rho>0\), and \(|P|=S\rho/2\). At physical height \(\rho y\), the lines through \(AB,ED\) have horizontal coordinates \[
x_-(y)=y(\xi-p),\qquad x_+(y)=1+y(\xi-1+q),
\qquad x_+(y)-x_-(y)=1-dy.
\] The signed representation also gives the strict relation \[
\frac\beta h+\frac\gamma k-1=\frac{\beta\gamma-st}{hk}>0.
\] If \(p+q\le k/J\), then \(L/h\le J/k\) would imply \[
\frac\beta h+\frac\gamma k=\frac Lh p+\frac Jk q
\le\frac Jk(p+q)\le1,
\] a contradiction. Thus \(p+q>k/J\).

If \(d\le1/L\), the two outer side lines meet at height \(y=1/d\ge L\). Together with the base \(AE\), they bound a triangle \(T\). The vertices \(A,E\) are on its base, \(B,D\) are on its two lateral sides, and the horizontal distances from \(C\) to these sides are \(p,q>0\). Thus all five vertices lie in the closed triangle. The triangulation \(ABC,ACE,CDE\) from \ref{lem-single} has disjoint interiors, so the entire polygon \(P\) lies in \(T\). Its area is \(\rho/(2d)\).

If \(d\ge1/L\), take \(T=ABE\). Its right side lies on \[
x_B(y)=1+y(\xi-p-1/L).
\] The horizontal distances from \(C\) to the left and right sides are \(p>0\) and \(h/L-p\ge q>0\). At the height \(J\) of \(D\), the corresponding distances are \[
x_D-x_-(J)=J(p+q)-k>0,\qquad
x_B(J)-x_D=J(d-1/L)\ge0.
\] Hence all vertices lie in \(T\), and the same triangulation gives \(P\subset T\). In this case, \(|T|=\rho L/2\). The two cases agree at \(d=1/L\) and yield \[
\frac{|P|}{|T|}=S\min(d,L^{-1}).\tag{G.9.1}
\]

In both containers, \(C\) is an interior point of \(T\) but a strictly reflex vertex of \(P\), with interior angle less than \(2\pi\). Choose a small disk around that vertex, contained in \(T\), and meeting no other side. The complement of \(P\) in this disk contains a sector of positive angle. Thus \(T\setminus P\) contains a nonempty open set, and strict domain monotonicity gives \(\lambda_1(P)>\lambda_1(T)\). The classical triangle inequality \(|T|\lambda_1(T)\ge4\pi^2/\sqrt3\) applies to this nondegenerate triangle. Consequently, \[
F(P)>\frac{4\pi^2}{\sqrt3}S\min(d,L^{-1}).
\] If \(F(P)\le2\pi^2\), it follows that \[
S\min(d,L^{-1})<\frac{\sqrt3}{2}.\tag{G.9.2}
\]

\end{proof}

\begin{definition}[Parameter data for a low-energy critical candidate]\label{def-candidate}

Candidate data are real numbers \(h,k,\beta,\gamma,\rho\) satisfying

\[
\begin{gathered}
h\ge k>0,\quad0<\beta\le h,\quad0<\gamma\le k,\quad L=1+h,\\
p=\frac\beta L,\quad q=\frac\gamma{1+k},\quad S=1+\beta+\gamma,\quad d=1-p-q>0.
\end{gathered}
\]

\[
\frac S{2L^2}\le\rho<\rho_+=\frac\beta{h(h+2)},\qquad S\min(d,L^{-1})<\sqrt3/2.
\]

Define \(f(y)=1-dy\) for \(0<y<1\), and

\[
f(1+x)=\max\{p(1-x/h)_+,q(1-x/k)_+\},\quad0<x<h.
\]

For \(A>0\), let \(\mu(A)\) be the first eigenvalue on \((0,L)\) of the quadratic form \(\int(|v'|^2+Af^{-2}|v|^2)\) on \(H_0^1(0,L)\cap L^2(f^{-2}dy)\), with the Friedrichs condition at the singular endpoint. The data are also required to satisfy, for every \(A>0\),

\[
\mu(A)\le\frac2S(\pi^2\rho+A/\rho).
\]

This defines a set of conditions and does not assert that the set is nonempty. If \(h=k\), the two ears may be exchanged to arrange \(p\ge q\).

\end{definition}

\begin{proof}[Proof of \ref{def-candidate}]

Set \(f(1)=\max(p,q)\); its value at this single point does not affect the form. Since \(d>0\) and \(p+q>0\), the connecting segment satisfies \(f\ge p+q>0\). On the tail, \(f\ge p(1-x/h)>0\). Hence \(f^{-2}\) is essentially bounded on every compact interior interval. Near the right endpoint, if \(k<h\), the shorter tail vanishes and \(f=p(L-y)/h\); if \(k=h\), then \(f=\max(p,q)(L-y)/h\). Every smooth compactly supported function therefore lies in the form domain, which is dense in \(L^2\).

If a sequence is Cauchy for the norm \(\|v\|_{H_0^1}^2+A\|v/f\|_2^2\), it converges in \(H_0^1\) to some \(v\), and its weighted components converge in \(L^2\) to some \(w\). Passing to almost everywhere convergent subsequences and using \(f>0\) almost everywhere yields \(w=v/f\). Thus the form domain is complete and the quadratic form is closed. Its norm controls the \(H_0^1\) norm. Compactness of \(H_0^1\hookrightarrow L^2\) on a bounded interval gives compact resolvent for the Friedrichs operator, and the Rayleigh infimum is attained by a nonzero normalized function. This establishes the existence of \(\mu(A)\) for every \(A>0\).

\end{proof}

\begin{lemma}[Every geometric critical candidate yields parameter data]\label{lem-geometry-input}

An area-critical pentagon as in \ref{def-domain}, represented as in \ref{lem-single} and satisfying \(F(P)\le2\pi^2\), can be transformed into the data of \ref{def-candidate} using only a Euclidean similarity and an exchange of the two ears.

\end{lemma}

\begin{proof}[Proof of \ref{lem-geometry-input}]

Use area-criticality from \ref{def-domain} and the signed representation from \ref{lem-single}, \[
P=(u,-su-\beta v,0,-\gamma u-tv,v),
\quad s,t\ge0,\quad\beta,\gamma>0,\quad\beta\gamma>st.
\] The target data and one-dimensional form are those of \ref{def-candidate}. Only Euclidean similarities and relabeling of the ears are used below.

\textbf{Ear-length restrictions.} By \ref{lem-ear}, each of the four outer ear sides is strictly shorter than its corresponding interior diagonal. Let \(a_0=|u|\), \(b_0=|v|\), and \(\theta=\angle(u,v)\). The inequality \(|AB|<|AC|\), together with \(s\ge0\), forces \(\cos\theta<0\). Linear independence gives \(c_0=-\cos\theta\in(0,1)\). Set \(X=\beta b_0/a_0\) and \(Y=\gamma a_0/b_0\). The two left-ear inequalities are \[
s^2+X^2-2c_0sX<1,\qquad
(s+1)^2+X^2-2c_0(s+1)X<1.
\] The first gives \(X<s+1\), and the second gives \(X>s\). Similarly, the right ear gives \(t<Y<t+1\). Thus \[
st<\beta\gamma<(s+1)(t+1).\tag{G.6.3}
\] Define \[
h=\beta+s,\quad k=\gamma+t,\quad L=1+h,\quad
p=\frac\beta{1+h},\quad q=\frac\gamma{1+k},\quad S=1+\beta+\gamma.
\] Exchange the ears, if necessary, so that \(h\ge k>0\); when \(h=k\), arrange \(p\ge q\) as well. Since \(s,t\ge0\), one has \(0<\beta\le h\), \(0<\gamma\le k\), and \[
d=1-p-q=
\frac{(s+1)(t+1)-\beta\gamma}{(1+h)(1+k)}>0.\tag{G.6.5}
\]

\textbf{Physical height and sections.} Normalize by a similarity so that \(A=(0,0),E=(1,0),C=(\xi,\rho)\), where \(\rho>0\). Then \[
B=(L\xi-\beta,L\rho),\quad
D=((1+k)\xi-t,(1+k)\rho),\quad |P|=S\rho/2.\tag{G.7.1}
\] Completing the square in \(|AB|^2<|AC|^2\) gives \[
\left(\xi-\frac{\beta(1+h)}{h(h+2)}\right)^2+\rho^2
<\left(\frac\beta{h(h+2)}\right)^2,
\qquad 0<\rho<\rho_+:=\frac\beta{h(h+2)}.\tag{G.7.2}
\] At height \(\rho y\), for \(0<y<1\), the section is the single interval between the outer sides and has width \(1-dy\). At \(y=1+x>1\), it splits into the two ear intervals, with widths \(p(1-x/h)_+\) and \(q(1-x/k)_+\). For \(0<x<k\), the open gap between the two ear intervals has width \[
x\left(\frac\beta h-\frac t k\right)=\frac{x(\beta\gamma-st)}{hk}>0,
\] by strict reflexivity. For \(x\ge k\), only the longer ear can remain. The longest connected component therefore has exactly the profile \[
f(y)=1-dy\quad(0<y<1),\qquad
f(1+x)=\max\{p(1-x/h)_+,q(1-x/k)_+\}\quad(0<x<h).\tag{G.7.3}
\] This is the longest component required by \ref{lem-largest-section}; the intervals on opposite sides of the notch are not combined into a fictitious connected section.

\textbf{Spectral and container restrictions.} Area-criticality and \ref{lem-energy-tensor} give energy \(\lambda_1(P)/2\) in each coordinate direction. Thus \ref{lem-weighted-section} applies to the physical height and actual sections just obtained. Under \(F(P)\le2\pi^2\), it gives the necessary conditions, for every \(A>0\), \[
\mu(A)\le\frac2S\left(\pi^2\rho+\frac A\rho\right),
\qquad \rho\ge\frac S{2L^2}.\tag{G.8.4}
\] Finally, the established representation and ear-length inequalities allow application of \ref{lem-triangle-container}, yielding \[
S\min(d,L^{-1})<\sqrt3/2.
\] All the geometric, aspect-ratio, and spectral conditions for every weight in \ref{def-candidate} have now been obtained.

\end{proof}


\section{The common-interface response principle}

\begin{lemma}[The response criterion at a common interface]\label{lem-interface}

On the finite interval \((0,1+h)\), let \(V\ge0\) be locally integrable on the two open subintervals, and suppose that the Friedrichs form domain embeds compactly into \(L^2\). Impose a Dirichlet condition at the left endpoint and the Friedrichs condition at the right endpoint. Assume that the potential is bounded near the interface and has one-sided limits there. Let \(E\) lie strictly below the first eigenvalue on each side when an additional zero trace is imposed at the interface. Write \(Q_{I,E}(v)=\int_I(|v'|^2+(V-E)|v|^2)\). Let \(g_H,g_T\) be the unique minimizers of \(Q_{I,E}\) with interface trace one and the stated outer endpoint conditions. Define \(\mathfrak d_H=g_H'(1)\) and \(\mathfrak d_T=-g_T'(1)\). Then \[
Q_{H,E}(g_H)=\mathfrak d_H,\qquad Q_{T,E}(g_T)=\mathfrak d_T.
\] Every function in the full form domain has a unique decomposition on the two sides as \(v=cg_H+z_H,cg_T+z_T\), where \(c=v(1)\) and \(z_H,z_T\) belong to the respective zero-trace form domains, and \[
\int_0^{1+h}(|v'|^2+(V-E)|v|^2)
=(\mathfrak d_H+\mathfrak d_T)|c|^2+Q_{H,E}(z_H)+Q_{T,E}(z_T).
\] Both zero-trace remainder forms are strictly positive definite. A nonnegative response sum implies that the first eigenvalue of the full operator is at least \(E\); a positive response sum implies that it is strictly greater than \(E\). Conversely, if the first eigenvalue of the full operator is strictly greater than \(E\), then the response sum is strictly positive.

\end{lemma}

\begin{proof}[Proof of \ref{lem-interface}]

\textbf{G.10. The Robin response criterion at a common interface.}

Consider the full operator \(-\partial_y^2+V\) on \((0,1+h)\), split at \(y=1\) into the connecting segment and the complete tail. Assume that \(E\) lies strictly below the first Dirichlet--Dirichlet eigenvalue on each side. Take the normalized solutions \[
g_H(0)=0,\quad g_H(1)=1,
\qquad g_T(1)=1,\quad g_T\text{ satisfies the right Friedrichs condition},
\] and define \[
\mathfrak d_H(E)=g_H'(1),\qquad
\mathfrak d_T(E)=-g_T'(1).
\tag{G.10.1}
\] Both derivatives refer to the same physical interface.

Given the interface value \(c_1=v(1)\) of a test function, use \(c_1g_H,c_1g_T\) as the energy-minimizing extensions. The remainders vanish at both endpoints of their respective segments. Integration by parts and the differential equation on each segment give \[
\begin{aligned}
\int_0^{1+h}(|v'|^2+(V-E)v^2)
={}&[\mathfrak d_H(E)+\mathfrak d_T(E)]|v(1)|^2\\
&+\mathcal E_{H,E}(z_H)+\mathcal E_{T,E}(z_T),
\end{aligned}
\tag{G.10.2}
\] where the last two forms are strictly positive definite on their zero-trace spaces. Thus a nonnegative response sum gives a first eigenvalue at least \(E\), and a positive response sum gives a first eigenvalue strictly greater than \(E\).

For the strict assertion, the decomposition makes the shifted energy of every nonzero function positive. Compact embedding ensures that the first eigenvalue is attained. If it were equal to \(E\), its eigenfunction would have zero interface value and zero remainders on both sides, hence would vanish identically, a contradiction.

We now justify the minimizing extensions and the treatment of singular endpoints. On either side interval \(I\), let \(\mathcal V_0\) denote the Friedrichs form domain with zero trace at the interface and the outer endpoint, and set \(Q_E(z)=\int_I(|z'|^2+V|z|^2-E|z|^2)\). If the first eigenvalue on that side is \(\lambda_D>E\), then \(Q_E(z)\ge(\lambda_D-E)\|z\|_2^2\). Moreover, \(Q_E\ge(1-E/\lambda_D)Q_0\) when \(E\ge0\), whereas \(Q_E\ge Q_0\) when \(E<0\). Hence \(Q_E\) is coercive in the form norm. Since \(V\) is bounded near the interface, choose a smooth trace lifting \(\psi\), supported near the interface, with interface value one. It has finite energy. The direct minimization argument below gives a unique minimizer \(g\) in \(\psi+\mathcal V_0\), satisfying \(Q_E(g,z)=0\) for every \(z\in\mathcal V_0\).

On either side of the interface, choose an endpoint neighborhood \(I=(1-\delta,1)\) or \(I=(1,1+\delta)\) on which \(V\) is bounded. The minimizer belongs to \(H^1(I)\), and its weak equation gives \(g''=(V-E)g\in L^2(I)\). Thus \(g\in H^2(I)\) and \(g'\in H^1(I)\). The absolutely continuous representative in one dimension extends continuously to the interface, so the one-sided derivative trace exists. Choose the smooth lifting \(\psi\) within this neighborhood, vanishing at its other endpoint. Integration by parts for \(g\in H^2\) and \(\psi\in H^1\) gives \(\int_I(g'\psi'+(V-E)g\psi)=[g'\psi]_{\partial I}\). The weak equation cancels the interior term; \(\psi(1)=1\) and its other endpoint value is zero. Therefore this is precisely the outward derivative from that side. No singular outer endpoint is involved. Testing with this lifting and integrating only within the interface neighborhood gives \(Q_E(g,\psi)=\partial_{\nu_I}g(1)\). Since \(g-\psi\in\mathcal V_0\), we obtain \(Q_E(g)=\partial_{\nu_I}g(1)\) without assuming a classical boundary value at the singular endpoint.

The relation between the side form domains and the full Friedrichs form domain follows explicitly from core approximation. The full domain is the closure, in the form norm, of compactly supported smooth functions on the full open interval. Nonnegativity of \(V\) makes this norm control \(H^1\). Let core functions \(v_n\) converge to \(v\). Continuity of the interface trace gives \(c_n=v_n(1)\to c=v(1)\). Choose smooth liftings \(\psi_H,\psi_T\) with interface value one, supported in the neighborhood where \(V\) is bounded and away from their outer endpoints. Each \(v_n|_I-c_n\psi_I\) is smooth up to the interface, vanishes there, and is supported away from the outer endpoint.

Such a smooth interface-zero function \(z\) can also be cut off near the interface to obtain a compactly supported smooth function in the side interval. Choose a cutoff equal to zero at distance less than \(\varepsilon\) from the interface, equal to one at distance greater than \(2\varepsilon\), and with derivative bounded by \(C/\varepsilon\). Because \(z(1)=0\) and \(z\) is smooth, \(|z|\le C|y-1|\). The integral of the squared gradient error is \(O(\varepsilon)\), and the squared \(L^2\) error is \(O(\varepsilon^3)\). Boundedness of the interface potential also makes the potential-energy error \(O(\varepsilon^3)\). Consequently each \(v_n|_I-c_n\psi_I\) belongs to the side zero-trace Friedrichs domain. These functions converge in the form norm to \(v|_I-c\psi_I\), which belongs to the same domain by closedness. Since \(g_I-\psi_I\) also belongs to it, so does \(v|_I-cg_I\).

The reverse gluing assertion is equally explicit. An element of a side zero-trace form domain is approximated by compactly supported smooth functions on that side. Their zero extensions are compactly supported and smooth on the full interval, so the zero extension belongs to the full form domain. The two liftings have the common interface value one; their gluing is a compactly supported \(H^1\) function in a region where \(V\) is bounded. Local convolution approximates it by compactly supported smooth functions both in \(H^1\) and in the potential-weighted \(L^2\) norm. The glued lifting therefore belongs to the full form domain. Adding the zero-extended remainders shows that any two admissible side functions with the same interface trace glue to a function in the full form domain, in particular the two minimizers \(g_H,g_T\). All cutoffs and convolutions occur where the potential is bounded. Possible singular behavior at the outer endpoints is handled by the original Friedrichs core approximation, without assuming classical derivative traces there.

A full-domain function \(v\) has a unique interface trace \(c\). On each side it can be written as \(v=cg+z\), with \(z\in\mathcal V_0\). Weak orthogonality cancels the cross term and gives exactly (G.10.2). The definition and closedness of the Friedrichs domain extend the identity from the smooth core to the full form domain, including the first eigenfunction of the full operator.

Existence of the minimizer requires only the coercivity and completeness already established. Denote the bilinear form by \(B_E\). On \(\psi+\mathcal V_0\), \[
Q_E(\psi+z)=Q_E(z)+2B_E(\psi,z)+Q_E(\psi).
\] Cauchy--Schwarz applied separately to the gradient, nonnegative potential, and \(L^2\) terms gives \(|B_E(\psi,z)|\le C\|z\|_{\mathrm{form}}\). Together with \(Q_E(z)\ge c\|z\|_{\mathrm{form}}^2\), this gives a finite infimum \(m\). Choose \(w_n\in\psi+\mathcal V_0\) with \(Q_E(w_n)\to m\). Their differences belong to \(\mathcal V_0\), and the parallelogram identity yields \[
\begin{aligned}
0\le Q_E(w_n-w_j)
&=2Q_E(w_n)+2Q_E(w_j)-4Q_E((w_n+w_j)/2)\\
&\le2Q_E(w_n)+2Q_E(w_j)-4m\longrightarrow0.
\end{aligned}
\] Hence \(w_n-\psi\) is Cauchy in the form norm. By closedness it converges to some \(z\in\mathcal V_0\), and \(Q_E(\psi+z)=m\). If two minimizers were distinct, the same identity and coercivity would give a contradiction; thus the minimizer is unique. Finally, differentiating along \(g+t z_0\) for any \(z_0\in\mathcal V_0\) gives \(B_E(g,z_0)=0\), the required weak orthogonality. No additional existence theorem is needed.

For the converse strict assertion, glue the two minimizing extensions with common interface value one to obtain a nonzero full-domain function \(g\). The decomposition gives \(Q_E(g)=\mathfrak d_H+\mathfrak d_T\). If the first eigenvalue of the full operator is \(\mu>E\), the Rayleigh principle gives \(Q_E(g)\ge(\mu-E)\|g\|_2^2>0\).

\end{proof}


\section{Elementary rational bounds}

\begin{lemma}[T.C Rational bounds for the constants]\label{lem-pi-rational}

We have \(493/50<\pi^2<987/100\) and \(6/7<\sqrt3/2<13/15\).

\end{lemma}

\begin{proof}[Proof of \ref{lem-pi-rational}]

Two integral identities with signed remainders provide the required rational bounds. For \(0<x<1\), \[
\begin{aligned}
\arctan x&=x-\frac{x^3}{3}+\int_0^x\frac{t^4}{1+t^2}\,dt,\\
\arctan x&=x-\frac{x^3}{3}+\frac{x^5}{5}
-\int_0^x\frac{t^6}{1+t^2}\,dt.
\end{aligned}
\] Both integrals are strictly positive, and \(0<\arctan x<x\). These identities follow by applying two finite polynomial divisions to \(1/(1+t^2)\) and integrating from \(0\) to \(x\); no infinite-series truncation is needed.

Let \(\theta=\arctan(1/5)\) and \(\varphi=\arctan(1/239)\). The tangent double-angle and subtraction identities give

\[
\tan(2\theta)=\frac5{12},\quad
\tan(4\theta)=\frac{120}{119},\quad
\tan(4\theta-\varphi)=1.
\]

Since \(\theta=\arctan(1/5)\in(0,\pi/4)\), we have \(2\theta\in(0,\pi/2)\). Moreover, \(\tan(2\theta)=5/12<1\) implies \(2\theta<\pi/4\), hence \(4\theta<\pi/2\). The inequality \(0<1/239<1/5\) gives \(0<\varphi<\theta\), so \(0<4\theta-\varphi<\pi/2\). Tangent is strictly increasing on this interval; consequently \(\tan(4\theta-\varphi)=1\) implies \(4\theta-\varphi=\pi/4\).

The signed integral remainders therefore yield

\[
\ell_\pi=
16\left(\frac15-\frac1{3\cdot5^3}\right)-\frac4{239}
<\pi
\]

and

\[
\pi<
u_\pi=
16\left(\frac15-\frac1{3\cdot5^3}+\frac1{5\cdot5^5}\right)
-4\left(\frac1{239}-\frac1{3\cdot239^3}\right).
\]

Both endpoints are positive rational numbers. Squaring and cross-multiplying gives

\[
\ell_\pi^2>\frac{493}{50},\qquad
u_\pi^2<\frac{987}{100}.
\]

This proves the two asserted rational bounds for \(\pi^2\).

The remaining bounds

\[
\frac67<\frac{\sqrt3}{2}<\frac{13}{15}
\]

also follow by squaring:

\[
\frac34-\frac{36}{49}=\frac3{196}>0,\qquad
\frac{169}{225}-\frac34=\frac1{900}>0.
\]

\end{proof}


\section{Middle-height parameters and explicit responses}

\begin{definition}[Intermediate-height parameter domains and explicit response functions]\label{def-analytic}

Write

\[
\Pi=\pi^2,\qquad c=\frac{\sqrt3}{2},\qquad L=1+h,\qquad
D_h=h^2+2h+2,\qquad C_h=\frac{D_h}{h+2}.
\]

Assume

\[
\frac38<h\le4,\qquad 0<a\le1,\quad0\le b\le1,
\]

\[
p=\frac{ha}{L},\quad q=\frac{hb}{L},\quad
w=p+q<1,\quad d=1-w,\quad m=\max(p,q),
\]

\[
S_0=ha+\frac1{1-q},\qquad
S_0<\frac{2aL^2}{h+2}.                                      \tag{M.1.1}
\]

The value \(b=0\) is allowed only for continuous extension to the boundary. Denote the domain defined by these conditions by \(\mathcal D_*\). The original target domain additionally requires

\[
S_0\min(d,L^{-1})<c.                                       \tag{M.1.2}
\]

Condition (M.1.1) is equivalent to

\[
a(1-q)>\frac1{C_h}.                                        \tag{M.1.3}
\]

Define

\[
\rho_+=\frac a{h+2},\quad A=\frac{\Pi\rho_+^2}{2},\quad
E=\frac{3\Pi\rho_+}{S_0},\quad K=E-A,\quad T=h^2E,
\]

\[
R=\frac{\rho_+h}{m},\qquad
\sigma=\frac{1+\sqrt{1+2\Pi R^2}}2.
\]

Combine the two gains on the connecting interval into

\[
\mathcal B(A,w)=\frac{A(1-w)}{2w}
+\frac{A(2-A)(1-w)^2}{20w}.                                \tag{M.1.4}
\]

Define

\[
\mathsf H(K)=1-\frac{2K}{\Pi-K}-\frac{2K}{4\Pi-K}
-\frac{2K(1/6-5/(4\Pi))}{1-K/(9\Pi)}.                          \tag{M.1.5}
\]

Further, set

\[
r=\sqrt{(2\sigma+5)^2-4T},\quad
C=\sigma+\frac12+\frac r2,\quad
D=2\sigma+1-\frac TC,
\]

\[
\mathsf T(\sigma,T)=\sigma-\frac TD.                       \tag{M.1.6}
\]

The function to be estimated is

\[
\mathcal G(h,a,b)=\mathsf H(K)+\mathcal B(A,w)
+\frac{\mathsf T(\sigma,T)}h.                              \tag{M.1.7}
\]

Here \(\mathcal D_*\) is defined only by (M.1.1) and the preceding parameter inequalities; \(\mathcal D\) additionally imposes (M.1.2). The responses involving radicals are initially defined where the radicals are real and the denominators are nonzero. Subsequent lemmas establish their validity throughout the domains used below. Define the boundary quantities, using \(p=ha/(1+h),\beta=ha,c=\sqrt3/2\):

\[
v_c=\frac{\beta p-(1-c)+\sqrt{(\beta p-(1-c))^2+4\beta p}}{2\beta},\quad q_c=1-v_c,\quad b_c=(1+h)q_c/h.
\]

Define \(\Psi(h)=\mathcal G(h,1,b_c(h,1))\), \(a_\triangle(h)=(c(1+h)-1)/h\), and \(\Phi(h)=\mathcal G(h,a_\triangle(h),0)\) whenever the corresponding parameters belong to \(\mathcal D_*\).

\end{definition}


\section{Rational terminal constructions}

\begin{definition}[T: explicit rational constructions]\label{def-terminal-rationals}

For real \(h\) for which the denominators below are nonzero, put \(L=1+h\) and \(\Pi_+=987/100\), and define \[
\begin{gathered}
s=\frac{1381h+1712}{500(h+2)},\quad v=\frac{2(h+7)}{67},\quad B=2s+5,\\
\mathcal C_h(\tau)=2s+3-\frac{13\tau}{10B},\qquad
\mathcal D_h(\tau)=2s+1-\frac{\tau}{\mathcal C_h(\tau)},\\
\mathcal L_h(\tau)=\frac1h\left(s-\frac{\tau}{\mathcal D_h(\tau)}\right).
\end{gathered}
\] Also define \[
\begin{aligned}
H_2(k)&=1-\frac{2k}{1-k}-\frac{2k}{4-k}-\frac{79k}{100(1-k/9)},\\
H_1(k)&=1-\frac{2k}{1-k}-\frac{129k}{100(1-k/4)}.
\end{aligned}
\] \[
\begin{gathered}
w=\frac{60h^3+143h^2+16h+8}{15hL(4h+9)},\quad
 d=\frac{52h^2+119h-8}{15hL(4h+9)},\\
 e=\frac{3(60h^2+127h-8)}{hL(h+2)(60h+127)},\\
 k=e-\frac1{2(h+2)^2},\quad A_- =\frac{493}{100(h+2)^2}.
\end{gathered}
\] Using \(\mathcal B\) from \ref{def-analytic}, put \[
\begin{aligned}
J&=1-\frac{2k}{1-k}-\left(\frac{129}{100}+\frac1{4L}\right)k,\\
G_* &=\frac{A_-(26h-5)}{60h^2}\left(1+\frac{33}{400L}\right),\\
Z_\Psi&=s+h(J+G_*)-v\Pi_+h^2e.
\end{aligned}
\] These rational functions are defined on their common domain; their positivity and the validity of their denominators are established separately in the terminal comparison chapter.

\end{definition}


\section{Exclusion of small-height candidates}

\begin{lemma}[Nonexistence of small-height candidates]\label{lem-small}

Any data satisfying \ref{def-candidate} must have \(h>3/8\).

\end{lemma}

\begin{proof}[Proof of \ref{lem-small}]

Use all the parameter conditions in \ref{def-candidate}. Put \(a=\beta/h\in(0,1]\), \(v=1-q\), and define

\[
C_h=\frac{h^2+2h+2}{h+2},\qquad
\underline S=ha+\frac1v.
\]

Since \(\gamma\le k\), one has \(\gamma\ge q/(1-q)\), and hence \(S\ge\underline S\). The candidate inequalities \(S/(2L^2)\le\rho<a/(h+2)\) give

\[
\underline S\le S<\frac{2aL^2}{h+2},
\qquad v>\frac1{aC_h},\quad a>\frac1{C_h}.
\tag{R.2.1}
\]

Also \(d=v-ha/L>0\). Therefore

\[
\begin{aligned}
Sd&\ge(ha+v^{-1})(v-ha/L)\\
&=1-\frac{h^2a^2}{L}+ha\left(v-\frac1{Lv}\right)\\
&>1+\frac h{C_h}-\frac{2hL}{h+2}a^2\\
&\ge R(h):=1+\frac{h(h+2)}{h^2+2h+2}-\frac{2h(h+1)}{h+2}.
\end{aligned}\tag{R.2.2}
\]

For the third line, the preceding expression is strictly increasing in \(v\). On the other hand,

\[
\frac SL>\frac{1+h/C_h}{L}=
Z(h):=\frac{2(h+1)}{h^2+2h+2},
\]

\[
Z(h)-R(h)=\frac{2h^2(h+1)^2}{(h+2)(h^2+2h+2)}>0.
\]

It follows that

\[
S\min(d,L^{-1})>R(h).\tag{R.2.3}
\]

The function \(R\) is strictly decreasing. Indeed, with \(D=h^2+2h+2>h+2\),

\[
R'(h)=\frac{4(h+1)}{D^2}-\frac{2(h^2+4h+2)}{(h+2)^2}
< -\frac{2h}{h+2}<0.
\]

Moreover,

\[
R(3/8)=\frac{12287}{14060}>\frac{13}{15}>\frac{\sqrt3}{2};
\]

The last comparison follows from \((13/15)^2-3/4=1/900\). If \(h\le3/8\), (R.2.3) contradicts the strict triangle-container bound in the candidate definition.

Thus every low-energy area-critical candidate has \(h>3/8\).

\end{proof}


\section{Head and tail responses and pole-free comparison}

\begin{lemma}[R.5 Two gains in the connecting-segment response]\label{lem-head-response}

Let \(0<A\le2\), \(0<d<1\), \(w=1-d\), and \(0<E-A<\pi^2\). Suppose that \(E\) lies below the first Dirichlet--Dirichlet eigenvalue for \(V(y)=A/(1-dy)^2\) on \((0,1)\). Let \(u\) solve \(-u''+Vu=Eu\), with \(u(0)=0,u(1)=1\). Then \(u'(1)\ge\sqrt{E-A}\cot\sqrt{E-A}+Ad/(2w)+A(2-A)d^2/(20w)\).

\end{lemma}

\begin{proof}[Proof of \ref{lem-head-response}]

Write \(f_0(y)=1-dy\), \(\kappa^2=E-A\), and \(c=(2-A)/20\), and set \[
b_1(y)=\frac{Ady^2}{2f_0(y)},\qquad
b_2(y)=\frac{cAd^2y^3}{f_0(y)},\qquad b=b_1+b_2.
\] Since \(0<A\le2\), we have \(c\ge0\) and \(cA\le1/20\). Differentiating the displayed rational functions gives \[
A(f_0^{-2}-1)-(b_1'+2b_1/y+b_1^2)
=\frac{Ad^2y^2}{f_0^2}\left(\frac12-\frac{Ay^2}{4}\right),
\] and \[
\frac{f_0^2}{Ad^2y^2}
\left(b_2'+\frac{2b_2}y+2b_1b_2+b_2^2\right)
=5c+cdy(-4+Ay^2+cAdy^3)\le5c=\frac{2-A}{4}.
\] The last inequality uses \(-4+Ay^2+cAdy^3\le-4+2+1/20<0\). Subtracting yields the residual estimate on the whole interval: \[
A(f_0^{-2}-1)-b'-2b/y-b^2
\ge\frac{A^2d^2y^2(1-y^2)}{4f_0^2}\ge0.\tag{R.5.0}
\]

Because \((\sin t-t\cos t)'=t\sin t>0\) for \(0<t<\pi\), we have \(t\cot t<1\). Thus \(0<\kappa<\pi\) implies \(\kappa\cot(\kappa y)\le1/y\). The positive function \[
g(y)=\sin(\kappa y)\exp\!\left(\int_0^y b(t)\,dt\right)
\] satisfies \[
\frac{-g''+(V-E)g}{g}
=A(f_0^{-2}-1)-b'-2\kappa\cot(\kappa y)b-b^2\ge0.
\]

To justify positivity of the compared solution, let \(u_0\) be the regular initial-value solution with \(u_0(0)=0,u_0'(0)=1\). If its first zero occurred at \(z\in(0,1]\), its restriction to \((0,z)\), extended by zero, would be a nonzero Dirichlet test function on \((0,1)\) with Rayleigh quotient \(E\), contradicting the assumed Dirichlet--Dirichlet spectral bound. Hence \(u_0>0\) on \((0,1]\), and the prescribed two-endpoint solution is \(u=u_0/u_0(1)\). For \(W=u'g-ug'\), \[
W'=u[-g''+(V-E)g]\ge0,\qquad W(0)=0.
\] Consequently \(u'(1)/u(1)\ge g'(1)/g(1)\). Since \(u(1)=1\) and \(f_0(1)=w\), \[
u'(1)\ge\sqrt{E-A}\cot\sqrt{E-A}
+\frac{Ad}{2w}+\frac{A(2-A)d^2}{20w}.\tag{R.5.1}
\] Both gains come from the explicit differential residual of the connecting-segment potential.

\end{proof}

\begin{lemma}[R.6.1 Regular tail solutions from a Volterra equation]\label{lem-regular-tail}

For real \(a>1\), let \(\lambda_D(a)\) be the first eigenvalue of the Friedrichs realization of \(-\partial_z^2+a(a-1)z^{-2}\) on \((0,1)\) with a Dirichlet condition at \(1\). Its form domain is \(H_0^1(0,1)\), and \(\lambda_D(a+1)\ge\lambda_D(a)\).

For every \(T>0\), the differential equation \[
-u_a''+\frac{a(a-1)}{z^2}u_a=Tu_a
\] has a unique solution on \((0,1]\) satisfying \(u_a/z^a\to1\) at zero. It has finite form energy at the singular endpoint, and \[
u_a=z^a(1+O(z^2)),\qquad
u_a'=az^{a-1}+O(z^{a+1}),\qquad
u_a'-\frac a z u_a=-\frac{T}{2a+1}z^{a+1}+O(z^{a+3}).
\] If \(0<T<\lambda_D(a)\), then \(u_a>0\) on \((0,1]\), and \(u_a/u_a(1)\) is the unique minimizer of the shifted form \[
\int_0^1\left(|v'|^2+\frac{a(a-1)}{z^2}|v|^2-T|v|^2\right)\,dz
\] over \(v\in H^1(0,1)\) with \(v(0)=0,v(1)=1\). Its outward response at \(1\) is \(D_a=u_a'(1)/u_a(1)\).

\end{lemma}

\begin{proof}[Proof of \ref{lem-regular-tail}]

For a real compactly supported smooth function, integration by parts and Cauchy--Schwarz give \[
\int_0^1\frac{v^2}{z^2}
=2\int_0^1\frac{vv'}z
\le2\left(\int_0^1\frac{v^2}{z^2}\right)^{1/2}
\left(\int_0^1|v'|^2\right)^{1/2},
\qquad
\int_0^1\frac{v^2}{z^2}\le4\int_0^1|v'|^2.
\] Apply this to the real and imaginary parts for complex functions. Approximation extends the inequality to \(H_0^1(0,1)\). Since \(a(a-1)>0\), the form norm is equivalent to the \(H_0^1\) norm, its closure has exactly that domain, and the smooth compactly supported functions are a form core. Compact embedding gives compact resolvent and a positive first eigenvalue. The potential increases with \(a>1\), so the variational principle on this common domain gives \(\lambda_D(a+1)\ge\lambda_D(a)\).

Write \(u_a=z^av\). For \(v\in C[0,1]\), define \[
(\mathcal K_av)(z)
=\int_0^z\xi^{-2a}\int_0^\xi\eta^{2a}v(\eta)\,d\eta\,d\xi
=\frac1{2a-1}\int_0^z
\eta\left[1-\left(\frac\eta z\right)^{2a-1}\right]v(\eta)\,d\eta
\] for \(z>0\), and \((\mathcal K_av)(0)=0\). The second equality follows by changing the order of integration. The kernel is nonnegative, is bounded above by \(\eta/(2a-1)\), and extends continuously at the origin.

Set \(\beta=T/(2a-1)>0\) and \(\|v\|_\beta=\sup_{[0,1]}e^{-\beta z^2}|v(z)|\). This norm is equivalent to the uniform norm, so \(C[0,1]\) is complete. The kernel bound gives \[
\begin{aligned}
e^{-\beta z^2}|T\mathcal K_a(v-w)(z)|
&\le\frac{T}{2a-1}\|v-w\|_\beta
e^{-\beta z^2}\int_0^z\eta e^{\beta\eta^2}\,d\eta\\
&=\frac12(1-e^{-\beta z^2})\|v-w\|_\beta
\le\frac12\|v-w\|_\beta.
\end{aligned}
\] Thus the successive iterates of \(v\mapsto1-T\mathcal K_av\) have geometrically decreasing differences and converge uniformly to its unique fixed point \[
v=1-T\mathcal K_av.
\] Differentiating its integral formula for \(z>0\) gives \[
v'(z)=-Tz^{-2a}\int_0^z\eta^{2a}v(\eta)\,d\eta.
\] Boundedness first gives \(v'=O(z)\) and \(v=1+O(z^2)\); substituting that estimate back gives \[
v'(z)=-\frac{T}{2a+1}z+O(z^3).
\] The derivative formula also yields \((z^{2a}v')'=-Tz^{2a}v\). Hence \(u_a=z^av\) solves the differential equation and has all the stated asymptotics.

To check uniqueness from the normalization alone, let \(\widetilde u/z^a=\widetilde v\to1\). Then \(p=z^{2a}\widetilde v'\) satisfies \(p'=-Tz^{2a}\widetilde v\). Since the latter is integrable at zero, \(p\) has a finite limit \(\ell\), and \[
p(z)=\ell+O(z^{2a+1}),\qquad
\widetilde v'(z)=\ell z^{-2a}+O(z).
\] If \(\ell\ne0\), integration produces a nonzero term \(\ell z^{1-2a}/(1-2a)\), contradicting boundedness of \(\widetilde v\) because \(a>1\). Thus \(\ell=0\). Integrating from zero now gives exactly the same fixed-point equation, whose uniqueness proves \(\widetilde u=u_a\).

The asymptotics give \(u_a(0)=0\) and \[
|u_a'|^2+z^{-2}|u_a|^2=O(z^{2a-2}),
\] so \(u_a\in H^1(0,1)\) has finite potential energy. Suppose \(T<\lambda_D(a)\). Its leading regular term is positive. If its first zero were \(z_0\le1\), extension by zero from \((0,z_0)\) would give a nonzero member of \(H_0^1(0,1)\) with Rayleigh quotient \(T\). Indeed, integration of the equation has no boundary term: \(u_a(z_0)=0\) and \(u_au_a'=O(z^{2a-1})\to0\) at zero. This contradicts \(T<\lambda_D(a)\), so \(u_a>0\) on \((0,1]\).

Normalize \(g=u_a/u_a(1)\). Its differential equation gives the weak equation against compactly supported smooth zero-trace tests. Finite energy, Hardy's inequality, and form-core density extend it to every \(w\in H_0^1(0,1)\). Every other function with the same traces is \(g+w\), and its shifted energy equals that of \(g\) plus that of \(w\), since the mixed term vanishes. If \(q_a\) denotes the unshifted form, then \[
q_a(w)-T\|w\|_2^2
\ge\left(1-\frac{T}{\lambda_D(a)}\right)q_a(w)>0
\qquad(w\ne0).
\] This proves the unique minimizing property. Integration by parts, with the already established zero boundary product at the singular endpoint, gives its shifted energy as \(g'(1)=D_a\).

\end{proof}

\begin{lemma}[R.6.2 Adjacent-order recurrence and a direct response bound]\label{lem-tail-recurrence}

Let \(s>1\), and let \(T>0\) lie below the first Friedrichs--Dirichlet eigenvalue of \(-\partial_z^2+s(s-1)z^{-2}\) on \((0,1)\). Normalize the regular solution by \(u_s/z^s\to1\), and write \(D_s=u_s'(1)/u_s(1)\). Then \(D_s=s-T/(s+D_{s+1})\), with positive denominator. If \(s=\sigma\), \(r=\sqrt{(2\sigma+5)^2-4T}>0\), \(C=\sigma+1/2+r/2>0\), and \(D=2\sigma+1-T/C>0\), then \(D_\sigma\ge\sigma-T/D\).

\end{lemma}

\begin{proof}[Proof of \ref{lem-tail-recurrence}]

Use the regular solutions, asymptotics, and spectral ordering in \ref{lem-regular-tail}. Write \(L_a=-\partial_z^2+a(a-1)z^{-2}\) for the differential expression, and set \[
A_a=\partial_z-\frac a z,\qquad B_a=\partial_z+\frac a z.
\] Composition gives \[
L_a=-B_aA_a,\qquad
L_{a+1}=-A_aB_a,\qquad
L_{a+1}A_a=A_aL_a.
\] These are differential identities on \((0,1]\); no Dirichlet operator-domain claim is made for the individual factors. Thus \(A_au_a\) solves the order-\(a+1\) equation at the same energy \(T\). The sharper asymptotic in \ref{lem-regular-tail} and uniqueness of its normalized solution give \[
A_au_a=-\frac{T}{2a+1}u_{a+1}.
\] Applying \(B_a\) and using \(B_aA_au_a=-Tu_a\) then gives \[
B_au_{a+1}=(2a+1)u_a.
\] When \(T<\lambda_D(s)\), spectral ordering gives \(T<\lambda_D(s+1)\), so both regular interface values are positive. Consequently \[
s+D_{s+1}=(2s+1)\frac{u_s(1)}{u_{s+1}(1)}>0,
\qquad
D_s=s-\frac{T}{s+D_{s+1}}.
\]

For the second assertion put \(a=\sigma\), \(A=2a+1\), \(M=2a+3\), and \(c=(2a+5-r)/2>0\). The given radical and denominators satisfy \[
T=c(2a+5-c),\qquad
C=M-c>0,\qquad D=A-\frac TC>0.
\] Let \(B=Ac+T\) and \(H(z)=AM-Bz^2\). Then \[
H(z)\ge H(1)=AM-Ac-T=CD>0,\qquad
M-cz^2\ge C>0
\quad(0\le z\le1).
\] Define the positive comparison function \[
w(z)=z^a\exp\left(-\int_0^z
\frac{T\eta(M-c\eta^2)}{H(\eta)}\,d\eta\right),
\qquad
\frac{w'}w=\frac a z-\frac{Tz(M-cz^2)}{H(z)}.
\] Its endpoint response is \[
\frac{w'(1)}{w(1)}=a-\frac{TC}{CD}=a-\frac TD.
\] To compute its residual, write \(x=z^2\), \(h(x)=AM-Bx\), and \(f(x)=T(M-cx)/h(x)\). Direct differentiation gives \[
f'=\frac{T^2M}{h^2},\qquad Af-T=\frac{T^2x}{h}.
\] Since \(M=A+2\), \(2M-A=2a+5\), and \(T=c(2a+5-c)\), \[
h+2M-(M-cx)^2
=[(2M-A)c-T]x-c^2x^2=c^2x(1-x).
\] Therefore \[
\frac{(L_a-T)w}{w}
=Af+2xf'-xf^2-T
=\frac{T^2x}{h^2}\bigl[h+2M-(M-cx)^2\bigr]
=\frac{T^2c^2z^4(1-z^2)}{H(z)^2}\ge0.
\] The inequality is strict for \(0<z<1\). By \ref{lem-regular-tail}, \(u_a>0\) at the given energy. Also \(w=z^a(1+O(z^2))\) and \(w'=az^{a-1}+O(z^{a+1})\), so the Wronskian \(W=u_a'w-u_aw'\) tends to zero at the singular endpoint. Its derivative is \[
W'=u_a(L_a-T)w>0\qquad(0<z<1).
\] Integration to one and division by positive interface values give \[
D_\sigma>\frac{w'(1)}{w(1)}
=\sigma-\frac TD=\mathsf T(\sigma,T),
\] where the last notation is from \ref{def-analytic}. In particular, the stated non-strict bound holds.

\end{proof}

\begin{lemma}[R.7 A rational lower bound for the cotangent response]\label{lem-cot-rational}

With \(\mathsf H\) defined in \ref{def-analytic}, one has \(\sqrt K\cot\sqrt K>\mathsf H(K)\) for \(0<K<\pi^2\).

\end{lemma}

\begin{proof}[Proof of \ref{lem-cot-rational}]

Write

\[
t=\frac{x^2}{\Pi},\qquad
a=\frac{3\Pi}{2}-\frac{45}{4},\qquad
p(x)=x(1-t)(1-t/4)(1-t/9)^a
\qquad(0\le x\le\pi).
\tag{R.7.1}
\]

Use the response definition in \ref{def-analytic}. The rational bounds in \ref{lem-pi-rational} imply

\[
\frac{49}{5}<\Pi<\frac{99}{10},\qquad
\frac{69}{20}<a<\frac{18}{5}.
\]

On the entire closed interval, \(1-t/9\ge8/9>0\). Thus its real power is a smooth function of \(x\) there. The remaining factors are polynomials. Consequently \(p\in C^\infty[0,\pi]\), with \(p(0)=p(\pi)=0\), finite endpoint derivatives, and \(p(x)>0\) for \(0<x<\pi\). In particular, the real exponent causes no endpoint singularity.

Logarithmic differentiation on the open interval gives

\[
\frac{xp'(x)}{p(x)}
=1-\frac{2t}{1-t}-\frac{t/2}{1-t/4}
-\frac{2at/9}{1-t/9}
=\mathsf H(x^2),
\tag{R.7.2}
\]

because \(a/(9\Pi)=1/6-5/(4\Pi)\). Every denominator is positive for \(0<x<\pi\).

\textbf{The differential residual.}

Define

\[
\begin{aligned}
c_0(\Pi)&=\frac{665}{72}+\frac{5\Pi}{27}-\frac{\Pi^2}{9},\\
c_1(\Pi)&=\frac{175}{432}-\frac{239\Pi}{162}+\frac{5\Pi^2}{36},\\
c_2(\Pi)&=-\frac{1295}{1296}+\frac{131\Pi}{324}-\frac{\Pi^2}{36},\\
R_\Pi(t)&=c_0+c_1t+c_2t^2-\frac{\Pi t^3}{324}.
\end{aligned}
\tag{R.7.3}
\]

The exact residual identity is

\[
-p''(x)-p(x)
=\frac{xt}{\Pi}(1-t/9)^{a-2}R_\Pi(t).
\tag{R.7.4}
\]

For a direct verification, put \(Q(t)=(1-t)(1-t/4)\), \(w(t)=1-t/9\), and \(f(t)=Q(t)w(t)^a\), so \(p(x)=xf(t)\). The chain rule gives

\[
p''(x)+p(x)=\frac{x}{\Pi}\bigl(6f'(t)+4tf''(t)+\Pi f(t)\bigr).
\]

Using \(Q'=-5/4+t/2\), \(Q''=1/2\), and \(w'=-1/9\) yields

\[
\begin{aligned}
&-\frac{6f'+4tf''+\Pi f}{w^{a-2}}\\
&\quad=-\left(\Pi Q-\frac{15}{2}+5t\right)w^2
+\frac a9\left(6-\frac{35t}{2}+\frac{11t^2}{2}\right)w
-\frac{4a(a-1)}{81}\,tQ\\
&\quad=tR_\Pi(t).
\end{aligned}
\]

The constant term cancels because \(2a/3=\Pi-15/2\). Substituting \(a=3\Pi/2-45/4\) in the remaining cubic gives precisely (R.7.3). This establishes the differential residual identity directly.

We now prove the residual sign on the full interval using only three quadratic functions of \(\Pi\). On \(49/5<\Pi<99/10\),

\[
c_0'(\Pi)=\frac{5-6\Pi}{27}<0,\qquad
c_2'(\Pi)=\frac{131-18\Pi}{324}<0.
\]

Evaluation at the single upper constant \(99/10\) therefore gives

\[
c_0(\Pi)>c_0(99/10)=\frac{323}{1800}>0,
\qquad
c_2(\Pi)>c_2(99/10)=\frac{4553}{16200}>0.
\tag{R.7.5}
\]

Also,

\[
c_1+2c_2=\frac{\Pi^2}{12}-\frac{2\Pi}{3}-\frac{2065}{1296},
\qquad
\frac{d}{d\Pi}(c_1+2c_2)=\frac{\Pi-4}{6}>0.
\]

Consequently

\[
c_1+2c_2<(c_1+2c_2)(99/10)
=-\frac{419}{16200}<0.
\tag{R.7.6}
\]

Since \(c_2>0\), for every \(0\le t\le1\) this proves

\[
R_\Pi'(t)=c_1+2c_2t-\frac{\Pi t^2}{108}
<c_1+2c_2<0.
\tag{R.7.7}
\]

The first inequality is strict also at the endpoints: at \(t=0\) it follows from \(c_2>0\), while at \(t=1\) it follows from \(\Pi>0\). Finally,

\[
R_\Pi(0)=c_0>0,
\qquad
R_\Pi(1)=\frac{700}{81}-\frac{8\Pi}{9}
<\frac{700}{81}-\frac89\frac{49}{5}
=-\frac{28}{405}<0.
\tag{R.7.8}
\]

The intermediate value theorem and strict monotonicity give exactly one \(t_*\in(0,1)\) with \(R_\Pi(t_*)=0\). Hence \(-p''-p\) is strictly positive on \(0<x<x_*:=\sqrt{\Pi t_*}\) and strictly negative on \(x_*<x<\pi\).

\textbf{Comparison using both zero endpoints.}

Let

\[
W(x)=\cos x\,p(x)-\sin x\,p'(x).
\]

The endpoint regularity established above makes \(W\) continuously differentiable on \([0,\pi]\). Since both \(p\) and \(\sin\) vanish at both endpoints,

\[
W(0)=W(\pi)=0.
\]

Differentiation gives

\[
W'(x)=\sin x\,[-p''(x)-p(x)].
\tag{R.7.9}
\]

Because \(\sin x>0\) on the open interval, \(W'>0\) before \(x_*\) and \(W'<0\) after \(x_*\). Therefore

\[
W(x)=\int_0^x W'(y)\,dy>0\quad(0<x\le x_*),
\]

and

\[
W(x)=-\int_x^\pi W'(y)\,dy>0\quad(x_*\le x<\pi).
\]

Dividing the strictly positive Wronskian by the positive product \(p(x)\sin x\) now yields

\[
x\cot x-\mathsf H(x^2)
=\frac{xW(x)}{p(x)\sin x}>0
\qquad(0<x<\pi).
\]

Taking \(x=\sqrt K\) proves the required strict inequality on the entire original interval \(0<K<\pi^2\).

\end{proof}

\begin{lemma}[Algebraic quantities and pole-free conditions in the middle-height region]\label{lem-pole-free}

Throughout \(\mathcal D_*\) of \ref{def-analytic}, including \(b=0\), one has \(0<A<1\), \(0<K<\pi^2\), \(1<\sigma<5/2\), and \(T<\pi^2(2\sigma+1)/3\). The quantities \(r,C,D\) are all strictly positive. The first Dirichlet--Dirichlet eigenvalue of the connecting-segment potential \(A/(1-dy)^2\) exceeds \(E\). The first Friedrichs--Dirichlet eigenvalue of the long-tail model potential \(Ah^2/(m^2z^2)\) on \((0,h)\) also exceeds \(E\).

\end{lemma}

\begin{proof}[Proof of \ref{lem-pole-free}]

Use the notation of \ref{def-analytic} and the bounds \(9<\Pi<10\) from \ref{lem-pi-rational}. Only the conditions defining \(\mathcal D_*\) are used below, so \(b=0\) is allowed.

\textbf{The head.} From \(0<a\le1\), \(h>3/8\), and \(S_0\ge1+ha\), \[
0<A=\frac{\Pi a^2}{2(h+2)^2}<\frac{320}{361}<1,
\qquad E\le\frac{3\Pi}{(1+h)(h+2)}<\Pi.
\] Here we used \(a/(1+ha)\le1/(1+h)\). Also \(q\le h/(1+h)\), so \(S_0\le1+2h\) and hence \(\rho_+S_0<2\). Thus \(E/A=6/(\rho_+S_0)>3\), giving \(0<K=E-A<\Pi\). The head potential is at least \(A\), so the interval Dirichlet Poincare inequality gives a first Dirichlet--Dirichlet eigenvalue at least \(\Pi+A>E\).

\textbf{The geometric energy bound for the tail.} Set \[
R=\frac{\rho_+h}{m}=\frac{a(1+h)}{(h+2)\max(a,b)},
\qquad \theta=\frac{mh}{S_0}.
\] Then \(0<R\le(1+h)/(h+2)\le5/6\), and \[
T=3\Pi R\theta,\qquad
T\le\frac{3\Pi h^2}{(1+h)(h+2)}\le\frac{8\Pi}{5}.\tag{R.4.3}
\] The last rational function increases for \(h>0\), since the numerator of its derivative is \(3h^2+4h>0\).

If \(p\ge q\), then \[
\theta\le\frac{ph}{1+(1+h)p}\le\frac{h^2}{(1+h)^2}\le\frac{16}{25}.
\] If \(q>p\), the bound \(p\le h/(1+h)\) and the strict area condition in the definition imply \[
S_0<\frac{2p(1+h)^3}{h(h+2)}
=2ph+2p+\frac{2p(1+h)}{h(h+2)}<2ph+2p+1.
\] On the other hand, \(S_0=1+(1+h)p+q+q^2/(1-q)\), so \[
\frac{q^2}{1-q}<ph+p-q<ph.
\] Consequently \[
\theta<\frac{h q(1-q)}{1+((1+h)/h)q^2}
\le\frac{h}{2(1+\sqrt{2+1/h})}\le\frac45.\tag{R.4.4}
\] The middle maximum follows directly by differentiation: for \(c_1=(1+h)/h\), the derivative of \(q(1-q)/(1+c_1q^2)\) has numerator \(1-2q-c_1q^2\), whose unique positive zero is \((1+\sqrt{1+c_1})^{-1}\). Substitution gives the displayed maximum. The final expression increases with \(h\) and equals \(4/5\) at \(h=4\).

Let \(B(R)=2+\sqrt{1+2\Pi R^2}=2\sigma+1\). Its derivative satisfies \(0<B'(R)<\sqrt{2\Pi}<36/5\). Hence \(B(R)-36R/5\) decreases, whereas \(B(R)-24/5\) increases. At \(R=2/3\), the bound \(\Pi>9\) gives \(B(R)>5\). Combining the two bounds for \(T\) above yields \[
\frac{3T}{\Pi}\le\min\!\left(\frac{36R}{5},\frac{24}{5}\right)
<2\sigma+1-\frac15.\tag{R.4.5}
\] In particular, \(T<\Pi(2\sigma+1)/3\). Moreover, \(R\le5/6\) and \(\Pi<10\) imply \(1<\sigma<5/2\).

\textbf{The tail Dirichlet spectrum and the singular endpoint.} On \((0,h)\), the model potential is \(\sigma(\sigma-1)/z^2\). The Hardy inequality \[
\int_0^h\frac{|u|^2}{z^2}\le4\int_0^h|u'|^2
\] follows first for compactly supported smooth functions by integration by parts and Cauchy--Schwarz, then extends by density to \(H_0^1(0,h)\). Therefore the Friedrichs form is closed with domain exactly \(H_0^1\), and compact embedding on the finite interval ensures attainment of its first eigenvalue.

For \(0<x<\pi\), define \(J(x)=\sin x-x\cos x\) and \(F(x)=(1-x^2/3)\sin x-x\cos x\). The identities \(J(0)=F(0)=0\), \(J'(x)=x\sin x>0\), and \(F'(x)=xJ(x)/3>0\) show directly that \[
1-x\cot x>\frac{x^2}{3}\qquad(0<x<\pi).
\] Taking \(g(z)=z^{\sigma-1}\sin(\pi z/h)>0\) and differentiating gives \[
-\frac{g''}{g}+\frac{\sigma(\sigma-1)}{z^2}
=\frac{\Pi}{h^2}+\frac{2(\sigma-1)}{z^2}
\left(1-\frac{\pi z}{h}\cot\frac{\pi z}{h}\right)
\ge\frac{\Pi(2\sigma+1)}{3h^2}.
\] For a compactly supported smooth function \(u=g\phi\), integration by parts gives Picone's identity: \[
\int_0^h\left(|u'|^2+\frac{\sigma(\sigma-1)}{z^2}|u|^2\right)
=\int_0^h g^2|\phi'|^2+
\int_0^h\left(-\frac{g''}{g}+\frac{\sigma(\sigma-1)}{z^2}\right)|u|^2.
\] Form-core density extends the resulting lower bound to the entire form domain. Since \(T=h^2E<\Pi(2\sigma+1)/3\), the first tail Dirichlet--Dirichlet eigenvalue is strictly greater than \(E\). No numerical information on Bessel zeros is used.

\textbf{The radical and recurrence denominators.} Since \(\Pi<10\) and \(\sigma>1\), \[
r^2-9>(2\sigma+5)^2-\frac{40}{3}(2\sigma+1)-9
=\frac43(3\sigma-2)(\sigma-1)>0.
\] Thus \(r>3\) and \(C=\sigma+1/2+r/2>\sigma+2\). The identity \(\sigma(\sigma-1)=\Pi R^2/2\) and \(\Pi<10\) give, uniformly for \(\sigma>1\), \[
\Pi R=\sqrt{2\Pi}\sqrt{\sigma(\sigma-1)}
<\frac92\left(\sigma-\frac12\right).
\] Indeed, \(2\Pi<20<81/4\) and \(\sigma(\sigma-1)=(\sigma-1/2)^2-1/4\). Since \(\theta\le4/5\), it follows that \[
T=3\Pi R\theta<\frac{54}{5}\left(\sigma-\frac12\right).
\] Consequently the original denominator has the uniform positive product bound \[
\begin{aligned}
CD&=(2\sigma+1)C-T\\
&>(2\sigma+1)(\sigma+2)-\frac{54}{5}\left(\sigma-\frac12\right)\\
&=2\left(\sigma-\frac{29}{20}\right)^2+\frac{639}{200}>0.
\end{aligned}
\] As \(C>0\), this proves \(D>0\). All radicals and denominators are therefore admissible without an order split.

\end{proof}

\begin{lemma}[Positive explicit response excludes middle-height candidates]\label{lem-robin}

For candidate data in \ref{def-candidate} with \(3/8<h\le4\), setting \(a=\beta/h,b=(1+h)q/h\) gives a point in \(\mathcal D\) of \ref{def-analytic}. If \(\mathcal G\ge0\) at this point, a contradiction follows. Thus every such candidate must satisfy \(\mathcal G<0\).

\end{lemma}

\begin{proof}[Proof of \ref{lem-robin}]

Use the parameters and the all-weight spectral condition of \ref{def-candidate}, together with the explicit functions of \ref{def-analytic}. We first check the parameter map. Since \(0<\beta\le h\) and \(0<\gamma\le k\le h\), \[
a=\beta/h\in(0,1],\qquad
0<q=\gamma/(1+k)\le k/(1+k)\le h/(1+h),
\qquad b=(1+h)q/h\in(0,1].
\] The inequality \(\gamma\le k\) gives \(k\ge q/(1-q)\) and \(\gamma=q(1+k)\ge q/(1-q)\), hence \[
S_0=ha+\frac1{1-q}\le1+\beta+\gamma=S.
\] The candidate inequalities \(S/(2L^2)\le\rho<a/(h+2)\) imply \[
S_0\le S<\frac{2aL^2}{h+2}.
\] Using \(d>0\) and the strict container bound also gives \[
w=p+q<1,\qquad S_0\min(d,L^{-1})\le S\min(d,L^{-1})<\sqrt3/2.
\] These are exactly all the conditions defining \(\mathcal D\).

Choose the single potential coefficient \(A_0=\Pi\rho_+^2/2\). Since \(\beta\le h\), \[
\frac\rho{\rho_+}\ge\frac{Sh(h+2)}{2\beta L^2}
\ge\frac{(1+\beta)h(h+2)}{2\beta L^2}
\ge\frac{h+2}{2L}>\frac12,
\] whereas the ear-length restriction gives \(\rho/\rho_+<1\). For \(x=\rho/\rho_+\in(1/2,1)\), \[
\frac32-x-\frac1{2x}=\frac{(1-x)(2x-1)}{2x}>0.
\] Thus the same \(A_0\) gives the strict upper bound throughout the admissible aspect-ratio interval: \[
\mu(A_0)\le\frac2S\left(\Pi\rho+\frac{A_0}{\rho}\right)
<\frac{3\Pi\rho_+}{S}\le\overline E:=\frac{3\Pi\rho_+}{S_0}.\tag{R.3.3}
\]

Compare the responses of the full operator at the same energy \(\overline E\). By \ref{lem-pole-free}, the first Dirichlet eigenvalues of both the head and the inverse-square model tail exceed \(\overline E\); also \(0<A_0<1\), \(0<K=\overline E-A_0<\Pi\), and all recurrence denominators are positive. The true tail profile satisfies \[
f(1+x)\le m(1-x/h),\qquad 0<x<h.
\] Therefore the true tail potential is at least the model potential \(A_0/[m^2(1-x/h)^2]\), and its first Dirichlet eigenvalue also exceeds \(\overline E\).

Both tail potentials are nonnegative, bounded on compact interior subintervals, and at most \(A_0h^2/[p^2(L-y)^2]\). The right-end Hardy inequality makes both zero-trace form domains equal to \(H_0^1(1,L)\). Choose the same smooth trace-one lifting \(\psi\), supported near the interface where both potentials are bounded. Both affine trace-one form domains are then \(\psi+H_0^1(1,L)\). Comparing the shifted energies function by function and taking minima, \ref{lem-interface} shows that the true tail response is no smaller than the model tail response. The full form is likewise closed and compactly embedded, and the one-sided interface limits of the potentials are \(A_0/w^2\) and \(A_0/m^2\). Thus every hypothesis of the interface lemma is satisfied.

By \ref{lem-head-response} and \ref{lem-cot-rational}, the head response obeys \[
\mathfrak d_H(\overline E)\ge\mathsf H(K)+\mathcal B(A_0,w).
\] For the model tail, set \(z=(h-x)/h\). Its equation becomes \[
-u''+\sigma(\sigma-1)z^{-2}u=T u,
\qquad T=h^2\overline E,\quad\sigma(\sigma-1)=A_0h^2/m^2.
\] By \ref{lem-regular-tail}, the normalized regular solution is positive, has finite endpoint form energy, and is the unique trace-one minimizing extension for this equation. Returning it to the physical interval multiplies its shifted energy and outward response by \(1/h\), so its physical response is \(D_\sigma(T)/h\). Applying the direct bound in \ref{lem-tail-recurrence} gives \[
\mathfrak d_T^{\rm coupled}(\overline E)
\ge\frac{D_\sigma(T)}h\ge\frac{\mathsf T(\sigma,T)}h.
\] The point \(z=1\) is always the original physical interface, and both sides retain the same function value.

The full response sum is therefore bounded below by \[
\mathfrak d_H(\overline E)+\mathfrak d_T^{\rm coupled}(\overline E)
\ge\mathsf H(K)+\mathcal B(A_0,w)+\frac{\mathsf T(\sigma,T)}h
=\mathcal G(h,a,b).
\] If \(\mathcal G\ge0\), \ref{lem-interface} gives \(\mu(A_0)\ge\overline E\), contradicting (R.3.3). Every candidate must therefore satisfy \(\mathcal G<0\).

\end{proof}


\section{Parameter monotonicity and boundary reduction}

\begin{lemma}[M.3.4 A uniform lower bound from head-response monotonicity]\label{lem-m34-identity}

Let \(3/8<h\le\sqrt2\), \(L=1+h\), and \(D_h=h^2+2h+2\), and set \[
v_h=\frac3{2L^2}-\frac1{2(h+2)^2},\qquad
R_2=\frac{2L^6}{5h^2(h+2)^2D_h}.
\] Then \(0<v_h<1\), and \[
\frac2{15}+\frac1{5(1-v_h)^2}+\frac h6-R_2
>\frac{321929}{15162000}>0.                    \tag{M.3.4}
\]

\end{lemma}

\begin{proof}[Proof of \ref{lem-m34-identity}]

Set \(y=L^{-2}\). Since \(h>3/8\) and \(h\le\sqrt2\), \[
\frac16<y<\frac{64}{121}<\frac45.
\] The lower bound follows from \((1+\sqrt2)^2=3+2\sqrt2<6\). Direct factorization gives \[
R_2=\frac2{5(1-y)^2(1+y)},\qquad
v_h-\frac54y=\frac{2-h^2}{4L^2(h+2)^2}\ge0.
\] Since \(h+2>L>11/8\), the definition gives the direct bounds \[
0<v_h<\frac3{2L^2}<\frac{96}{121}<1.
\] Thus all denominators are strictly positive.

For \(0<y<4/5\), define \[
f(y)=\frac1{(1-5y/4)^2},\qquad
g(y)=\frac2{(1-y)^2(1+y)}.
\] The identity \[
(1+y)^2-\frac89(1+3y)=\frac{(3y-1)^2}{9}\ge0
\] and the inequality \(1-5y/4<1-y\) give \[
f'(y)>\frac{5}{2(1-y)^3},\qquad
g'(y)=\frac{2(1+3y)}{(1-y)^3(1+y)^2}
\le\frac9{4(1-y)^3}.
\] Hence \(f-g\) is strictly increasing. Since \(v_h\ge5y/4\) and \(h/6>1/16\), \[
\begin{aligned}
\frac2{15}+\frac1{5(1-v_h)^2}+\frac h6-R_2
&>\frac{47}{240}+\frac{f(y)-g(y)}5\\
&>\frac{47}{240}+\frac15\left(\frac{576}{361}-\frac{432}{175}\right)
=\frac{321929}{15162000}>0.
\end{aligned}
\] This derivative comparison for a rational function covers the entire stated height interval.

\end{proof}

\begin{lemma}[Derivatives of the head response, tail response, and gain]\label{lem-derivatives}

On \(\mathcal D_*\) from \ref{def-analytic}, let \(X=-\mathsf H'(K),M=-\partial_T\mathsf T(\sigma,T),N=\partial_\sigma\mathsf T(\sigma,T)\). Then

\[
X>1/3,\quad M>1/(2\sigma+1)>1/6,\quad N>0,
\]

\[
\frac13+\frac2\Pi[(1-K/\Pi)^{-2}-1]\le X\le\frac1{3(1-K/\Pi)^2},
\]

\[
\mathcal B_A=\frac d{2w}+\frac{(1-A)d^2}{10w},\quad \mathcal B_w=-A\left[\frac{12-A}{20w^2}-\frac{2-A}{20}\right]>-\frac{3A}{5w^2}.
\]

The same tail-response monotonicity holds on every connected region where \(T>0,r>0,C>0,D>0\).

\end{lemma}

\begin{proof}[Proof of \ref{lem-derivatives}]

We use the rational constant bounds in \ref{lem-pi-rational}; any weaker bounds below follow by inclusion or by squaring. Use all definitions and notation from \ref{def-analytic}.

The notation is from \ref{def-analytic}, and the expressions are well defined by \ref{lem-pole-free}. Write

\[
X=-\mathsf H'(K),\quad
M=-\partial_T\mathsf T(\sigma,T),\quad
N=\partial_\sigma\mathsf T(\sigma,T).
\]

Direct differentiation gives

\[
X>\frac13,\qquad M>\frac1{2\sigma+1}>\frac16,\qquad N>0.  \tag{M.2.1}
\]

Indeed,

\[
C_T=-1/r,\quad D_T=-1/C-T/(rC^2)<0,
\]

\[
M=1/D-TD_T/D^2>1/D>1/(2\sigma+1),
\]

whereas

\[
C_\sigma=1+(2\sigma+5)/r>0,\quad
D_\sigma=2+TC_\sigma/C^2>0,\quad
N=1+TD_\sigma/D^2>0.
\]

For the head response,

\[
X\ge\frac13+\frac2\Pi\left[(1-K/\Pi)^{-2}-1\right],\qquad
X\le\frac1{3(1-K/\Pi)^2}.                                  \tag{M.2.2}
\]

More explicitly, setting \(x=K/\Pi\in(0,1)\) gives

\[
X=\frac2\Pi(1-x)^{-2}+\frac1{2\Pi}(1-x/4)^{-2}
+\left(\frac13-\frac5{2\Pi}\right)(1-x/9)^{-2}.
\]

All three coefficients are positive because \(\Pi>9\), and their sum is \(1/3\). For \(n=1,4,9\), we have \(1-x\le1-x/n<1\), hence \(1\le(1-x/n)^{-2}\le(1-x)^{-2}\). Since \(x>0\), the first term is strictly larger than its value at zero, and the other two are no smaller than their values at zero. This proves \(X>1/3\) and the lower estimate. Bounding all three terms by the common factor \((1-x)^{-2}\) gives the upper estimate.

These expressions come from differentiating the terms of (M.1.5); their positive coefficients sum to \(1/3\) at \(K=0\).

The partial derivatives of the two-term gain are

\[
\mathcal B_A=\frac d{2w}+\frac{(1-A)d^2}{10w},
\]

\[
\mathcal B_w=-A\left[\frac{12-A}{20w^2}-\frac{2-A}{20}\right]
>-\frac{3A}{5w^2}.                                        \tag{M.2.3}
\]

The full derivative, including both gain terms, is used throughout. For the larger region in the final assertion, we do not assume \(\sigma>1\). Nevertheless, \(D>0\) and \(T/C>0\) imply \(2\sigma+1=D+T/C>0\), hence \(\sigma>-1/2\). Thus \(2\sigma+5>0\), and the sign computations \(C_\sigma=1+(2\sigma+5)/r>0\), \(D_\sigma>0\), and \(N>0\) remain valid throughout that region. Likewise, \(\partial_T\mathsf T<0\) uses only \(T,r,C,D>0\).

\end{proof}

\begin{lemma}[Strict monotonicity and one-sided derivatives when the long tail dominates]\label{lem-q-branch-monotone}

Use \(\mathcal D_*\) from \ref{def-analytic}. For fixed \(h,a\), the function \(q\mapsto\mathcal G(h,a,(1+h)q/h)\) is strictly increasing over all admissible \(q\le p\). At admissible interior points with \(q<p\), the ordinary partial derivative satisfies \(\partial_q\mathcal G>0\). If \(q=p\) is admissible, the left derivative along the branch \(q\le p\) is strictly positive. No two-sided partial derivative at \(q=p\) is asserted.

\end{lemma}

\begin{proof}[Proof of \ref{lem-q-branch-monotone}]

We use the rational constant bounds in \ref{lem-pi-rational}; any weaker bounds below follow by inclusion or by squaring. Use all definitions and notation from \ref{def-analytic}.

Fix \(h,a\) and restrict to the branch \(q\le p\). Here \(m=p\), so \(R=L/(h+2)\) and \(\sigma\) are independent of \(q\). Write \(v=1-q\). All derivatives below refer to the extension of this branch; at \(q=p\) they give the left derivative. Use the definitions in \ref{def-analytic}, the estimates in \ref{lem-pole-free}, and the derivative bounds in \ref{lem-derivatives}.

\[
-E_q=\frac{3\Pi a}{(h+2)(1+hav)^2}>0,
\]

\[
\mathcal G_q=(-E_q)(X+hM)+\mathcal B_w.
\]

By (M.2.3),

\[
\mathcal G_q>
(-E_q)\left[X+\frac h6-\mathcal R\right],\quad
\mathcal R=\frac{a(1+hav)^2}{10(h+2)w^2}.                 \tag{M.3.2}
\]

Using \(w\ge ha/L,\ v\le1,\ a>1/C_h\), we obtain

\[
\mathcal R\le
\frac{L^2}{10h^2(h+2)}
\max_{1/C_h\le a\le1}\left(\frac1a+2h+h^2a\right).
\]

The function in parentheses is convex, so its maximum occurs at an endpoint. Consequently,

\[
\mathcal R\le\max(R_1,R_2),
\]

\[
R_1=\frac{L^4}{10h^2(h+2)},\qquad
R_2=\frac{2L^6}{5h^2(h+2)^2D_h}.                         \tag{M.3.3}
\]

Since \(h^2-C_h=L(h^2-2)/(h+2)\), the larger term is \(R_1\) when \(h\ge\sqrt2\) and \(R_2\) when \(h\le\sqrt2\).

If \(h\ge\sqrt2>4/3\), then \(L>7/3\) and \(1-L^{-2}>40/49>0\). Since \(h(h+2)=L^2-1\),

\[
\frac{6R_1}{h+2}
=\frac3{5(1-L^{-2})^2}
<\frac{7203}{8000}<1.
\]

Consequently \(R_1<(h+2)/6=1/3+h/6\).

In the other height range, the geometric condition gives

\[
K>\Pi v_h,\qquad v_h=\frac3{2L^2}-\frac1{2(h+2)^2}\in(0,1).
\]

Indeed, (M.1.1) gives \(E=3\Pi a/((h+2)S_0)>3\Pi/(2L^2)\), while \(a\le1\) gives \(A\le\Pi/[2(h+2)^2]\). Subtraction yields \(K>\Pi v_h\). Since \(h+2>L>11/8\), we have \(3/L^2>1/(h+2)^2\) and \[
0<v_h<\frac3{2L^2}<\frac{96}{121}<1.
\]

By (M.2.2) and \(\Pi<10\),

\[
X>\frac2{15}+\frac1{5(1-v_h)^2}.
\]

The uniform lower bound in \ref{lem-m34-identity} gives

\[
\frac2{15}+\frac1{5(1-v_h)^2}+\frac h6-R_2
>\frac{321929}{15162000}>0.
\]

Substitution into (M.3.2) proves the asserted sign on this branch.

The argument applies on the larger domain without the container restriction (M.1.2). Thus, when \(q\) is subsequently decreased to zero or to the container boundary, intermediate points need not satisfy (M.1.2). For fixed \(h,a\), admissibility of \(q\) is equivalent to \(q\ge0\), \(q\le\min(p,h/L)\), \(q<1-p\), and \(q<1-1/(aC_h)\). The admissible set is therefore an interval. The branch formula is smooth in its interior and has the same left derivative limit at an admissible point \(q=p\). Given two admissible points \(q_1<q_2\le p\), integrating the strictly positive branch derivative gives \(\mathcal G(q_2)>\mathcal G(q_1)\). This proves strict monotonicity and the left-derivative assertion.

\end{proof}

\begin{lemma}[Uniform positivity when the short tail is wider]\label{lem-other-branch}

On \(\mathcal D_*\cap\{q>p\}\) from \ref{def-analytic}, we have \(\mathcal G>1/150+\mathcal B(A,w)>1/150\).

\end{lemma}

\begin{proof}[Proof of \ref{lem-other-branch}]

Use the definitions in \ref{def-analytic}, the pole-free conclusions in \ref{lem-pole-free}, and the rational bounds in \ref{lem-pi-rational}. In particular, \[
\frac{49}{5}<\Pi<\frac{99}{10}<10,
\qquad 0<A<1,\qquad 0<K<\Pi.
\] Suppose \(q>p\). All the estimates below use only \(\mathcal D_*\); the additional container restriction defining \(\mathcal D\) is not used.

\textbf{The geometric branch forces \(h>2\).}

Write \(L=1+h\) and \(D_h=h^2+2h+2\). The area condition is equivalent to \[
p(1-q)>j(h):=\frac{h(h+2)}{LD_h}.
\tag{M.4.2}
\] If \(h\le1\), then \(q\le h/L\le1/2\). The function \(q(1-q)\) is increasing on \([0,1/2]\), so \[
p(1-q)<q(1-q)\le\frac{h}{L^2}<j(h),
\] where the last difference is \(j(h)-h/L^2=h^2/(L^2D_h)>0\). This contradicts (M.4.2), and hence \(h>1\).

For \(1\le h\le2\), the identity \[
h^3-h^2-4h+2=(h-2)(h-1)(h+2)-2\le-2
\] gives \(j(h)>1/4\). This again contradicts \(p(1-q)<q(1-q)\le1/4\). Therefore \[
2<h\le4.
\tag{M.4.3}
\]

\textbf{A common order and energy range.}

Since \(m=q\) on this branch, (M.4.2) gives \[
R=\frac{Lp}{(h+2)q}
>\frac{h}{D_hq(1-q)}
\ge\frac{4h}{D_h}\ge\frac8{13}.
\] The last inequality follows because \((4h/D_h)'=4(2-h^2)/D_h^2<0\) on \([2,4]\). Consequently, \[
\sigma(\sigma-1)=\frac{\Pi R^2}{2}
>\frac{1568}{845}
>\frac{741}{400}
=\frac{39}{20}\frac{19}{20}.
\] As \(\sigma>1\), this implies \[
\sigma>\frac{39}{20}.
\tag{M.4.4}
\]

Also \(S_0=Lp+(1-q)^{-1}>Lp+(1-p)^{-1}\). Therefore \[
E<\frac{3\Pi L}{h(h+2)}
\frac{p(1-p)}{1+Lp(1-p)}
\le\frac{3\Pi L}{h(h+2)(h+5)}.
\tag{M.4.5}
\] Here \(0<p<1\), \(p(1-p)\le1/4\), and \(t/(1+Lt)\) is increasing for \(t\ge0\).

Set \[
e(h)=\frac{3L}{h(h+2)(h+5)},\qquad
u(h)=h^2e(h)=\frac{3hL}{(h+2)(h+5)}.
\] The logarithmic derivatives are \[
\frac{e'}e=\frac1L-\frac1h-\frac1{h+2}-\frac1{h+5}<0,
\] \[
\frac{u'}u=\frac1h+\frac1L-\frac1{h+2}-\frac1{h+5}>0.
\] In the second line, pair \(1/h\) with \(-1/(h+2)\) and \(1/L\) with \(-1/(h+5)\). Thus (M.4.3) and (M.4.5) imply \[
\frac K\Pi<e(h)<e(2)=\frac9{56}<\frac16,
\qquad
T<\Pi u(h)\le\Pi u(4)=\frac{10\Pi}{9}<11.
\tag{M.4.6}
\]

\textbf{One affine lower bound for the tail response.}

By (M.4.4) and (M.4.6), \[
r^2=(2\sigma+5)^2-4T
>\frac{3521}{100}>\left(\frac{59}{10}\right)^2.
\] It follows that \[
C=\sigma+\frac12+\frac r2>\frac{27}{5},
\] \[
D=2\sigma+1-\frac TC
>\frac{49}{10}-\frac{55}{27}
=\frac{773}{270}>\frac{20}{7}.
\] All the quantities used in these divisions are strictly positive. Hence the same estimate holds throughout the branch: \[
\mathsf T(\sigma,T)=\sigma-\frac TD
>\frac{39}{20}-\frac{7T}{20}.
\tag{M.4.7}
\] \textbf{Reduce the whole response to one rational function.}

Let \(x=K/\Pi\). The three coefficients \[
\frac2\Pi,\qquad \frac1{2\Pi},\qquad
\frac13-\frac5{2\Pi}
\] are positive and sum to \(1/3\). Applying their common denominator bound in the formula for \(\mathsf H\) gives \[
\mathsf H(K)\ge1-\frac{K}{3(1-x)}
>1-\frac{2K}{5}>1-\frac{2E}{5},
\] where (M.4.6) was used. Combining this with (M.4.7) and \(T=h^2E\) yields \[
\mathcal G-\mathcal B(A,w)
>1+\frac{39}{20h}-\frac{8+7h}{20}E.
\] Use (M.4.5) and \(\Pi<10\) in this expression. The exact partial-fraction identity gives \[
\begin{aligned}
\mathcal G-\mathcal B(A,w)
&>1+\frac{39}{20h}
-\frac{3(8+7h)L}{2h(h+2)(h+5)}\\
&=F(h):=1+\frac3{4h}+\frac3{2(h+2)}-\frac{54}{5(h+5)}.
\end{aligned}
\tag{M.4.8}
\]

\textbf{A single tangent proves the uniform positive margin.}

For \(2\le h\le4\), \[
\frac{(h+5)^3}{2}F''(h)
=\frac34\left(\frac{h+5}{h}\right)^3
+\frac32\left(\frac{h+5}{h+2}\right)^3-\frac{54}{5}.
\] The first ratio is at least \(9/4>2\), and the second is at least \(3/2\). Therefore \[
\frac{(h+5)^3}{2}F''(h)
>6+\frac{81}{16}-\frac{54}{5}
=\frac{21}{80}>0.
\] Thus \(F\) is strictly convex on this full interval. At the single interior point \(h=3\), \[
F(3)=\frac15,\qquad F'(3)=\frac{61}{2400}>0.
\] Its supporting tangent gives, for every \(2<h\le4\), \[
F(h)\ge\frac15+\frac{61}{2400}(h-3)
>\frac{419}{2400}>\frac16>\frac1{150}.
\] Substitution into (M.4.8) proves the first stated inequality. Finally, \(0<A<1\) and \(0<w<1\) imply \(\mathcal B(A,w)>0\), giving its last strict inequality as well.

\end{proof}

\begin{lemma}[Monotonicity on the boundary with vanishing short-tail width]\label{lem-zero-monotone}

On the admissible boundary \(q=0\) from \ref{def-analytic}, with \(3/8<h\le4\) fixed, the function \(a\mapsto\mathcal G(h,a,0)\) is strictly decreasing.

\end{lemma}

\begin{proof}[Proof of \ref{lem-zero-monotone}]

We use the rational constant bounds in \ref{lem-pi-rational}; any weaker bounds below follow by inclusion or by squaring. Use all definitions and notation from \ref{def-analytic}.

Use the notation of \ref{def-analytic} together with \ref{lem-pole-free} and \ref{lem-derivatives}.

\[
A=\frac{\Pi L^2p^2}{2h^2(h+2)^2},\qquad
E=\frac{3\Pi Lp}{h(h+2)(1+Lp)}.
\]

The order \(\sigma\) is constant. Define

\[
r_0=\frac{E_p}{A_p}
=\frac{3h(h+2)}{Lp(1+Lp)^2}
=\frac{3(h+2)}{a(1+ha)^2}.
\]

By \ref{lem-derivatives},

\[
\frac p{2A}\frac{d\mathcal G(h,a,0)}{dp}
=(1-r_0)X-r_0hM+J_0,                                    \tag{M.5.2}
\]

where the full two-term gain gives

\[
J_0=\mathcal B_A+\frac p{2A}\mathcal B_w
=\frac{12-3A}{40p}-\frac{7-2A}{10}+\frac{6-5A}{40}p.      \tag{M.5.3}
\]

If \(p<1/2\), then

\[
J_0\le \frac3{10p}-\frac7{10}+\frac{3p}{20}
=:J_*(p).                                                \tag{M.5.4}
\]

Indeed, the difference is \(A(1-p)(5p-3)/(40p)<0\). For fixed \(0<A<1\), the expression in (M.5.3) satisfies \[
\partial_pJ_0=-\frac{12-3A}{40p^2}+\frac{6-5A}{40}
<-\frac{6+2A}{40}<0\qquad(0<p<1).
\] Thus, if \(p\ge1/2\), \[
J_0(p,A)\le J_0(1/2,A)=-\frac{2+A}{80}<-\frac1{40}.
\] This pointwise bound holds for every \(0<A<1\), hence also for the geometric value \(A=A(p)\).

First assume \(3/8<h\le1\). Then \(p\le1/2,r_0>1\). By \(X>1/3,M>1/6\), it suffices to prove

\[
r_0(h+2)/6-1/3-J_*(p)>0.
\]

After multiplication by \(p\), the left-hand side becomes

\[
\frac{h(h+2)^2}{2L(1+Lp)^2}
-\frac3{10}+\frac{11p}{30}-\frac{3p^2}{20}.
\]

Since \(p\le h/L\), the first term is at least \(h(h+2)^2/(2L^3)\). This function is increasing on \([3/8,1]\) and has value \(1083/2662>3/10\) at \(3/8\). The sum of the remaining two terms is positive as well. Thus (M.5.2) is negative.

Next assume \(1\le h\le4,r_0\ge1\). If \(p\ge1/2\), the conclusion is immediate. If \(p<1/2\), then

\[
\frac{r_0hp}{6}
=\frac{h^2(h+2)}{2L(1+Lp)^2}
\ge\frac3{4(1+2p)^2}.
\]

Here the left-hand side is strictly increasing in \(h\). Set \(f(p)=3/[4(1+2p)^2]\). Since \[
f''(p)=\frac{18}{(1+2p)^4}>0,
\] the tangent at \(p=1/2\) gives \[
f(p)\ge\frac3{16}-\frac38\left(p-\frac12\right)
=\frac38-\frac{3p}{8}.
\] Consequently, \[
\frac3{4(1+2p)^2}-pJ_*(p)
\ge\frac3{40}+\frac{p(13-6p)}{40}>\frac3{40}.
\] Thus \(r_0h/6>J_0\), so (M.5.2) is again negative.

Finally assume \(r_0<1\). Since \(a\le1\), this forces \(h>5/2\), and \(r_0\ge18/25\). We must also have \(p>1/2\); otherwise

\[
r_0\ge\frac{24h(h+2)}{L(h+3)^2}>1
\quad(1\le h\le4),
\]

a contradiction. The numerator minus denominator in the last comparison is \(-h^3+17h^2+33h-9>0\). Thus \(J_0<0\). Also \(K/\Pi<4/21\), so (M.2.2) gives \(X<147/289<3/5\). Consequently the right-hand side of (M.5.2) is less than

\[
\frac7{25}\frac35-\frac{18}{25}\frac{5/2}{6}
=-\frac{33}{250}<0.
\]

This establishes the required derivative sign in every case. We record the monotonicity computations and the bounds used in the final branch. Let \(F(h)=h(h+2)^2/[2(1+h)^3]\). Then

\[
F'(h)/F(h)=\frac1h+\frac2{h+2}-\frac3{1+h}=\frac{2-h}{h(1+h)(h+2)}>0\quad(3/8\le h\le1).
\]

For fixed \(p>0\), let \(B(h)=h^2(h+2)/[2(1+h)(1+(1+h)p)^2]\). Since \(p/(1+(1+h)p)<1/(1+h)\),

\[
B'(h)/B(h)=\frac2h+\frac1{h+2}-\frac1{1+h}-\frac{2p}{1+(1+h)p}>\frac{h+4}{h(1+h)(h+2)}>0.
\]

These identities justify the endpoint lower bounds in the first two cases.

Since \(a(1+ha)^2\) is increasing for \(a>0\), the bound \(a\le1\) gives \(r_0\ge R_0(h):=3(h+2)/(1+h)^2\). Moreover, \(R_0'(h)=-3(h+3)/(1+h)^3<0\). Therefore \(r_0\ge54/49>1\) when \(h\le5/2\), and \(r_0\ge18/25\) when \(h\le4\). This proves both bounds used in the final case. On \(q=0\), we also have

\[
K/\Pi<\frac{3a}{(h+2)(1+ha)}\le\frac3{(h+2)(1+h)}<\frac4{21}\quad(h>5/2),
\]

The second inequality uses the increase of \(a/(1+ha)\) with \(a\); the last uses \((h+2)(h+1)>(9/2)(7/2)\). The admissible values of \(a\) form exactly the interval \(1/C_h<a\le1\). The strictly negative derivative therefore proves the asserted strict decrease.

\end{proof}

\begin{lemma}[Monotonicity on the container equality boundary]\label{lem-container-monotone}

Fix \(h\). In \(\mathcal D_*\cap\{q\le p\}\) from \ref{def-analytic}, along the curve \(F(p,q)=(Lp+(1-q)^{-1})(1-p-q)=c\), the function \(\mathcal G\) strictly decreases as \(a\) increases.

\end{lemma}

\begin{proof}[Proof of \ref{lem-container-monotone}]

We use the rational constant bounds in \ref{lem-pi-rational}; any weaker bounds below follow by inclusion or by squaring. Use all definitions and notation from \ref{def-analytic}.

Apply \ref{def-analytic}, \ref{lem-pole-free}, and \ref{lem-derivatives}. Fix \(h\) and set

\[
F(p,q):=\left(Lp+\frac1{1-q}\right)(1-p-q)=c,\qquad q\le p.
                                                                  \tag{M.6.1}
\]

The conditions of the larger geometric domain (M.1.1) remain in force. Write \(v=1-q,d=v-p\). We have \(F_q=-Lp-p/v^2<0\). Furthermore, \(F=c<1\) implies \(Lv(v-p)<1\), whence

\[
F_p=Lv-2Lp-1/v<-Lp<0.
\]

Thus \(q_p<0,v_p>0\) along the equality curve.

On the admissible part of this equality boundary, we shall prove

\[
\frac d{da}\mathcal G(h,a,b_c(h,a))<0.                    \tag{M.6.2}
\]

Since \(S_0=c/d\), we have

\[
E=\frac{3\Pi Lpd}{ch(h+2)},\qquad
A=\frac{\Pi L^2p^2}{2h^2(h+2)^2}.
\]

Write \(\alpha=Lv^2\). Implicit differentiation gives

\[
d_p=\frac{d(1-\alpha)}{p(1+\alpha)},\qquad
E_p=\frac{2E}{p(1+\alpha)}.
\]

Let

\[
r_c=\frac{E}{A(1+\alpha)}.
\]

The derivative of the full gain is

\[
J_c=\mathcal B_A-\frac{d\,\mathcal B_w}{2A}
       \frac{1-\alpha}{1+\alpha}.
\]

Consequently,

For fixed \(h\), along the admissible branch \(q\le p\) we have \(m=p=ha/L\), so \(R=\rho_+h/m=(a/(h+2))h/(ha/L)=L/(h+2)\). Thus \(\sigma\) is independent of \(p\) on this branch, and \(\sigma_p=0\): no derivative of the tail order occurs in the chain rule. If an endpoint satisfies \(q=p\), only the one-sided derivative within the admissible branch is intended. Monotonicity follows from the derivative in the interior of the admissible interval and continuity at its endpoints.

\[
\frac p{2A}\frac{d\mathcal G}{dp}
=(1-r_c)X-r_chM+J_c.                                     \tag{M.6.3}
\]

We first obtain uniform geometric bounds. Set \(x=hav=Lp v\). Then (M.1.3) gives

\[
x>\frac{h(h+2)}{D_h}.
\]

The identity \(c=(1+x)(1-p/v)\) gives

\[
\alpha=\frac{x(1+x)}{x+1-c}.                              \tag{M.6.4}
\]

Let \(y=d/v=c/(1+x)>c/L\). Direct substitution yields

\[
r_c=\frac{6h(h+2)y^2}{cL(c-y^2)}
>\frac{6h(h+2)}{L(L^2-c)}>1.                             \tag{M.6.5}
\]

The last inequality holds for \(3/8<h\le4\): its numerator minus denominator is

\[
-h^3+3h^2+(9+c)h+c-1
\ge(5+c)h+c-1>0.
\]

Also \(S_0>1+h/C_h=2L^2/D_h\), so

\[
d<\frac{cD_h}{2L^2}.                                    \tag{M.6.6}
\]

First suppose \(h\ge1/2\). Then \(x>5/13\) and \(x^2>25/169>1/7>1-c\), hence \(\alpha>1\). Therefore

\[
J_c\le\mathcal B_A\le\frac{3d}{5w}
<\frac{3cD_h}{5(2L^2-cD_h)}.
\]

Using \(X>1/3,M>1/6,r_c>1\), a sufficient condition for negativity in (M.6.3) is

\[
\frac{h(h+2)^2}{L(L^2-c)}
-\frac13-\frac{3cD_h}{5(2L^2-cD_h)}>0.                    \tag{M.6.7}
\]

We prove the stronger uniform lower bound \(1/8\) directly. Set \(y=L^{-1}\), and extend to \(0\le y\le2/3\) the function \[
\mathcal F_c(y)=
\frac{(1+y)(1-y^2)}{1-cy^2}
-\frac13-\frac{3c(1+y^2)}{5(2-c-cy^2)}.
\] This is exactly the left-hand side of (M.6.7). By \(6/7<c<7/8\), \[
1-cy^2>\frac{11}{18},\qquad
2-c-cy^2>\frac{53}{72}.
\] For \(\eta>0\), wherever the denominator is positive, write \(W_\eta(y)=y^2/(\eta-cy^2)\). Direct differentiation gives \[
W_\eta'(y)=\frac{2\eta y}{(\eta-cy^2)^2}\ge0,\qquad
W_\eta''(y)=\frac{2\eta(\eta+3cy^2)}{(\eta-cy^2)^3}>0.
\] The identities \[
\frac{(1+y)(1-y^2)}{1-cy^2}
=1+y-(1-c)(1+y)W_1(y),
\] \[
\frac{c(1+y^2)}{2-c-cy^2}
=\frac{c}{2-c}+\frac{2c}{2-c}W_{2-c}(y)
\] show that \[
\mathcal F_c''(y)
=-(1-c)\bigl(2W_1'(y)+(1+y)W_1''(y)\bigr)
-\frac{6c}{5(2-c)}W_{2-c}''(y)<0.
\] Thus \(\mathcal F_c\) lies above its endpoint chord. The endpoint bounds are \[
\mathcal F_c(0)=\frac23-\frac{3c}{5(2-c)}
>\frac15,
\] \[
\begin{aligned}
\mathcal F_c(2/3)
&=\frac{25}{3(9-4c)}-\frac13-\frac{39c}{5(18-13c)}\\
&>\frac{175}{117}-\frac13-\frac{273}{265}
=\frac{4099}{31005}>\frac18.
\end{aligned}
\] In the second line, the first rational term increases with \(c\) and is bounded below using \(c>6/7\); the subtracted term increases with \(c\) and is bounded above using \(c<7/8\). Hence \(\mathcal F_c>1/8\) on \([0,2/3]\), proving (M.6.7) for every \(h\ge1/2\).

Next suppose \(3/8<h<1/2\). We have \(x>57/185,c>6/7\).

To justify substituting these two lower bounds, set \(f(x,c)=x(1+x)/(x+1-c)\). Its denominator is positive for \(x>0,0<c<1\), and

\[
f_x=\frac{x^2+(1-c)(2x+1)}{(x+1-c)^2}>0,\qquad
f_c=\frac{x(1+x)}{(x+1-c)^2}>0.
\]

Therefore \(f(x,c)>f(57/185,6/7)=48279/54020\), giving the required strict comparison.

Hence

\[
\alpha>\frac{48279}{54020}>\frac78.
\]

By (M.6.6) and \(c<13/15\),

\[
d<\frac{481}{726}<\frac23,\qquad w>\frac13.
\]

By (M.6.5), \(h>3/8\), \(h<1/2\), and \(c>6/7\),

\[
r_c>\frac{6h(h+2)}{L(L^2-c)}
>\frac{6(3/8)(19/8)}{(3/2)(9/4-6/7)}
=\frac{133}{52}>\frac52.
\]

Moreover \(K>\Pi v_h\) and \(v_h\ge44/75\), so (M.2.2) gives \(X>5/4\). Since \((1-\alpha)/(1+\alpha)<1/15\),

\[
J_c\le\frac{3d}{5w}+\frac{d}{50w^2}
<\frac65+\frac3{25}=\frac{33}{25}.
\]

The right-hand side of (M.6.3) is consequently less than

\[
-\frac32\frac54+\frac{33}{25}=-\frac{111}{200}<0.
\]

This proves (M.6.2) in all cases.

For completeness, we justify the quantities used at small heights and the shape of the admissible interval. Define

\[
v_h=\frac3{2(1+h)^2}-\frac1{2(h+2)^2}.
\]

From \(S_0<2aL^2/(h+2)\) we obtain \(E>3\Pi/(2L^2)\); together with \(A\le\Pi/[2(h+2)^2]\), this gives \(K>\Pi v_h\). Since \(v_h'=-3/(1+h)^3+1/(h+2)^3<0\), for \(3/8<h<1/2\) we have \(v_h>v_{1/2}=44/75\). The lower bound in \ref{lem-derivatives} therefore yields

\[
X>\frac13+\frac15[(75/31)^2-1]=\frac{18797}{14415}>\frac54,
\]

where we used \(\Pi<10\).

For fixed \(h\) and \(0<p\le h/L\), we have \(F(p,0)=1+Lp(h-Lp)/L\ge1>c\), while \(F(p,1-p)=0\) and \(F_q<0\). Hence the root \(q_c(p)\in(0,1-p)\) is unique. Along the curve \(q_c'(p)<0\), so \(q_c(p)-p\) strictly decreases and \(a(1-q_c)\) strictly increases. Both \(q_c\le p\) and \(a(1-q_c)>1/C_h\) are therefore preserved when \(p\) increases. Moreover \(q_c\le p\le h/L\) gives \(b\le1\), and all other domain conditions are automatically preserved. The admissible part is thus an interval; its strictly negative derivative proves strict monotonicity between any two admissible points.

\end{proof}

\begin{lemma}[Reduction to two terminal curves]\label{lem-boundary}

On \(\mathcal D\) from \ref{def-analytic}, the following hold. If \(q>p\), then \(\mathcal G>1/150\). If \(q\le p\) and \(S_0d<c\), then \(\mathcal G(h,a,b)>\Psi(h)\), and the terminal point belongs to \(\mathcal D_*\) and satisfies \(q_c\le p\). If \(q\le p\) and \(S_0/L<c\), then \(\mathcal G(h,a,b)>\Phi(h)\), with \(\sqrt3-1<h\le4\) and the terminal point in \(\mathcal D_*\). At least one of these cases applies.

\end{lemma}

\begin{proof}[Proof of \ref{lem-boundary}]

We use the rational constant bounds in \ref{lem-pi-rational}; any weaker bounds below follow by inclusion or by squaring. Use all definitions and notation from \ref{def-analytic}.

Use the boundary definitions of \ref{def-analytic}, continuous well-definedness from \ref{lem-pole-free}, and the respective monotonicity statements of \ref{lem-q-branch-monotone}, \ref{lem-other-branch}, \ref{lem-zero-monotone}, and \ref{lem-container-monotone}. For fixed \(h,a\), put \(\beta=ha,p=\beta/L,\epsilon=1-c\), and define

\[
v_c(h,a)=
\frac{\beta p-\epsilon+
\sqrt{(\beta p-\epsilon)^2+4\beta p}}{2\beta},
\]

\[
q_c(h,a)=1-v_c(h,a),\qquad b_c(h,a)=\frac Lh q_c(h,a).
                                                                  \tag{M.7.1}
\]

This gives the unique root of \(S_0d=c\). Indeed,

\[
F(p,0)=(1+\beta)(1-\beta/L)
=1+\frac{\beta(h-\beta)}L\ge1>c,
\]

while \(F\) strictly decreases with \(q\) and vanishes at \(q=1-p\). Hence \(0<q_c<1-p\).

The first boundary function is

\[
\boxed{\Psi(h)=\mathcal G(h,1,b_c(h,1)).}                 \tag{M.7.2}
\]

Only the geometrically admissible heights inherited from an actual candidate need to be considered. Equivalently, within \(3/8<h\le4\) we require

\[
\mathscr R(h)<c,\qquad
\mathscr R(h)=
1+\frac{h(h+2)}{D_h}-\frac{2hL}{h+2}.                    \tag{M.7.3}
\]

This is exactly the condition that \(a=1,q=q_c\) satisfy (M.1.3). The inequality \(q_c\le p\) holds automatically in this height range.

At the terminal point \(a=1\), if \(3/8<h<1\), then

\[
F(p,p)=1-\frac{2h^2}{1+h}
\le\frac{35}{44}<c,
\]

so \(q_c<p\). If \(h\ge1\), then \(q_c<1-p\le p\). Thus no additional comparison between the two ear widths is needed beyond (M.7.3).

The second boundary function is

\[
a_\triangle(h)=\frac{cL-1}{h},\qquad
\boxed{\Phi(h)=\mathcal G(h,a_\triangle(h),0).}           \tag{M.7.4}
\]

Its relevant height range is

\[
\sqrt3-1<h\le4.                                         \tag{M.7.5}
\]

The condition \(b=0\) is used only to extend the comparison function continuously; it does not assert that the boundary still represents a domain with five active sides.

\textbf{Reduction to two one-variable boundary functions.} For every parameter in the original domain (M.1.1)--(M.1.2), at least one of the following cases applies:

\begin{itemize}
\item
  If \(q>p\), then \(\mathcal G>1/150\).
\item
  If \(q\le p\) and \(S_0d<c\), then
\end{itemize}

\[
  \mathcal G(h,a,b)>\Psi(h);
\]

\begin{itemize}
\tightlist
\item
  If \(q\le p\) and \(S_0/L<c\), then
\end{itemize}

\[
  \mathcal G(h,a,b)>\Phi(h).
\]

The first case is \ref{lem-other-branch}. In the second case, \(q>q_c(h,a)\), so \ref{lem-q-branch-monotone} gives \(\mathcal G(h,a,b)>\mathcal G(h,a,b_c(h,a))\). Increase \(a\) to \(1\) along \(S_0d=c\). Then \(q_c\) decreases and \(a(1-q_c)\) increases, so all conditions of the larger geometric domain remain valid. Applying \ref{lem-container-monotone} proves the second case.

In the third case, first decrease \(q\) to zero using \ref{lem-q-branch-monotone}. The condition \(S_0/L<c\) implies

\[
a<a_\triangle(h)<1.
\]

Then increase \(a\) to \(a_\triangle(h)\) using \ref{lem-zero-monotone}, proving the third case. For this branch to be nonempty, we also need \(a_\triangle>1/C_h\), equivalently,

\[
c>\frac{2L}{L^2+1},
\]

Since \(L>1\), this is equivalent to \(L>\sqrt3\), namely (M.7.5).

We finish by verifying the equality root and both admissible paths in detail. Write \(v=1-q\). Then

\[
F_q=-Lp-p/v^2<0,\qquad F_p=Lv-2Lp-1/v.
\]

On \(F=c<1\), the identity \(F-1=p[L(v-p)-1/v]<0\) gives \(Lv(v-p)<1\), hence \(F_p<-Lp<0\). The implicit-function formula \(q_c'(p)=-F_p/F_q<0\) proves monotonicity of the equality root. As \(a\) increases along the root curve, \(p\) increases, \(q_c\) decreases, and \(a(1-q_c)\) strictly increases. Thus \(q_c\le p\) and \(a(1-q_c)>1/C_h\) are preserved. The root itself lies in \((0,1-p)\), and \(q_c\le p\le h/L\), so all remaining \(\mathcal D_*\) conditions hold up to \(a=1\).

In the third branch, first fix \(h,a\) and decrease \(q\) from its original value to zero. The inequalities \(q\ge0\), \(q\le p\), and \(q\le h/L\) are preserved; \(w=p+q\) does not increase and stays below one, while \(a(1-q)\) does not decrease and remains strictly above \(1/C_h\). The whole segment therefore lies in \(\mathcal D_*\). From \(S_0/L<c\) and \(S_0\ge1+ha\), we obtain \(0<a<a_\triangle\); moreover \(c<1\) implies \(a_\triangle<1\). Keeping \(q=0\) fixed and increasing \(a\) to \(a_\triangle\) preserves \(a>1/C_h\) and \(p=ha/L\le h/L<1\), so the second segment also lies in \(\mathcal D_*\). If the original point has \(q=0\), the first comparison may be an equality, but the second segment has positive length. Its strict decrease from \ref{lem-zero-monotone} makes the final comparison strict.

\end{proof}


\section{Analytic positivity on the terminal curves}

\begin{lemma}[T.2 A common affine lower bound for the tail response]\label{lem-common-tail-line}

Let \(3/8\le h\le4\) and \(L=1+h\), and put \[
s=\frac{1381h+1712}{500(h+2)},\qquad B=2s+5,\qquad
v=\frac{2(h+7)}{67},\qquad
T_0=\frac{987}{100}\frac{3h^2}{L(h+2)}.
\] Define \[
\mathcal C(\tau)=2s+3-\frac{13\tau}{10B},\qquad
\mathcal D(\tau)=2s+1-\frac{\tau}{\mathcal C(\tau)},\qquad
\mathcal L_h(\tau)=\frac1h\left(s-\frac{\tau}{\mathcal D(\tau)}\right).
\] Then \(s>1\) and \(0<T_0/B^2<30/169\); moreover \(\mathcal C,\mathcal D>0\) for \(0\le\tau\le T_0\), with \(\mathcal D(\tau)>1/v+1/2000\), and \[
\mathcal L_h(\tau)>\frac{s}{h}-\frac{v\tau}{h}
\qquad(0<\tau\le T_0).
\]

\end{lemma}

\begin{proof}[Proof of \ref{lem-common-tail-line}]

Write \(R=L/(h+2)\). We have \(s=331/500+(21/10)R>1\) and \(B>(51h+82)/[5(h+2)]\). Define \[
f(h)=\frac{h^2(h+2)}{(1+h)(51h+82)^2}.
\] Its logarithmic derivative is \[
\frac{f'(h)}{f(h)}
=\frac2h+\frac1{h+2}-\frac1{h+1}-\frac{102}{51h+82}
=\frac{113h^2+410h+328}{h(h+1)(h+2)(51h+82)}>0.
\] Therefore \(f(h)\le f(4)=24/102245<1/4225\), and \[
\frac{T_0}{B^2}<\frac{2961}{4}f(h)
<\frac{3(987/100)}{169}
=\frac{2961}{16900}<\frac{30}{169}.
\] For \(0\le\tau\le T_0\), successive substitution gives \[
\mathcal C(\tau)>\frac{20s+24}{13}>0,\qquad
\mathcal D(\tau)>\frac{400s^2+284s-438}{13(20s+24)}>0.
\] The last numerator is increasing for \(s\ge1\) and equals \(246\) at \(s=1\).

We prove the stronger uniform estimate \[
\mathcal D(T_0)>\frac1v+\frac1{2000}.
\] Use the compact height coordinate \[
x=\frac h{h+2}\in[\xi_0,\xi_1],\qquad
\xi_0=\frac3{19},\quad \xi_1=\frac23,
\] and put \[
\begin{gathered}
\kappa=\frac{2961}{50},\qquad \gamma=\frac{13\kappa}{10},\qquad
b(x)=\frac{803}{125}+\frac{21x}{10},\qquad
c(x)=\frac{1053}{125}+\frac{21x}{10},\\
W(x)=(1+x)b(x)-\frac{\gamma x^2}{c(x)},\qquad
\widehat D(x)=\frac{553}{125}+\frac{21x}{10}-\frac{\kappa x^2}{W(x)}.
\end{gathered}
\] The change of variable gives \[
T_0=\frac{\kappa x^2}{1+x},\qquad
W(x)=(1+x)\mathcal C(T_0)>0,\qquad
\widehat D(x)=\mathcal D(T_0).
\] Thus the denominators were already proved positive before the change of variable.

The essential estimate is a uniform curvature bound for the original continued fraction. Let \(c_*=1228/125\). Since \(x\le2/3\), we have \(c(x)\le c_*\), and the elementary rational comparisons \[
\frac\gamma{c_*}=\frac{38493}{4912}>\frac{39}{5},\qquad
\frac{\gamma c(0)}{c_*^2}=\frac{40533129}{6031936}>\frac{33}{5},\qquad
\frac{c(0)}{c_*}=\frac{1053}{1228}>\frac45
\] yield, by direct differentiation, \[
\begin{aligned}
W-xW'&=\frac{803}{125}
+x^2\left(\frac{\gamma c(0)}{c(x)^2}-\frac{21}{10}\right)
>\frac{32}{5}+\frac92x^2=:U(x),\\
W''&=\frac{21}{5}-\frac{2\gamma c(0)^2}{c(x)^3}<0,\\
0<W(x)&<\frac{13}{2}+\frac{43}{5}x-\frac{57}{10}x^2=:V(x).
\end{aligned}
\] For the second inequality, the displayed bounds give \(\gamma c(0)^2/c(x)^3>(33/5)(4/5)>21/10\). For the third, bound the constant and linear coefficients of \((1+x)b(x)\) by \(13/2\) and \(43/5\), and use \(\gamma/c(x)>39/5\) for the quadratic coefficient.

We next compare the positive quadratic \(U\) with \(V\) without a high-degree expansion. The function \(f(x)=U(x)^{2/3}\) is strictly convex, since \[
f''(x)=\frac{6(32/5+3x^2/2)}{U(x)^{4/3}}>0.
\] Its tangent at \(x_*=4/(3\sqrt5)\), where \(U(x_*)=8\), gives \[
\frac52 U(x)^{2/3}\ge\frac{22}{3}+2\sqrt5\,x.
\] Moreover \(\sqrt5>11/5\) and \[
\frac{22}{3}+\frac{22x}{5}-V(x)
=\frac{57}{10}\left(x-\frac7{19}\right)^2+\frac{17}{285}>0.
\] Consequently \(V(x)<(5/2)U(x)^{2/3}\). The exact second derivative identity now gives \[
\begin{aligned}
\left(\frac{\kappa x^2}{W(x)}\right)''
&=\kappa\left\{\frac{2(W-xW')^2}{W^3}-\frac{x^2W''}{W^2}\right\}\\
&>\frac{2\kappa U(x)^2}{V(x)^3}
>\frac{16\kappa}{125}
=\frac{23688}{3125}>\frac{15}{2}.
\end{aligned}
\] Hence \(\widehat D''<-15/2\) on the whole compact height interval.

It remains to compare this curvature with the affine-bound threshold \[
g(x)=\frac1v=\frac{67(1-x)}{2(7-5x)}.
\] For any function \(f\), let \(\ell_f\) denote its affine interpolant at \(\xi_0,\xi_1\), and write \(\omega(x)=(x-\xi_0)(\xi_1-x)\). Strong concavity gives \[
\widehat D(x)\ge\ell_{\widehat D}(x)+\frac{15}{4}\omega(x).
\] The fractional-linear function \(g\) has the exact chord remainder \[
g(x)-\ell_g(x)
=\frac{335\omega(x)}{(7-5\xi_0)(7-5\xi_1)(7-5x)},
\] whose coefficient is bounded by \[
\frac{57285}{14278}<\frac{201}{50}.
\] Direct substitution into the continued fraction gives \[
\begin{aligned}
\widehat D(\xi_0)-\frac{4561}{1000}
&=\frac{23710670971}{60245409817000}>0,\\
\widehat D(\xi_1)-\frac{307}{100}
&=\frac{3931}{146701500}>0.
\end{aligned}
\] As \(g(\xi_0)=268/59\) and \(g(\xi_1)=67/22\), both endpoint gaps exceed \(37/2000\): \[
\frac{4561}{1000}-\frac{268}{59}=\frac{1099}{59000}>\frac{37}{2000},\qquad
\frac{307}{100}-\frac{67}{22}=\frac{27}{1100}>\frac{37}{2000}.
\] Finally, \[
0\le\omega(x)\le\frac{(\xi_1-\xi_0)^2}{4}
=\frac{841}{12996}<\frac1{15}.
\] Subtracting the chord formulas therefore proves \[
\widehat D(x)-g(x)
>\frac{37}{2000}-\frac{27}{100}\omega(x)
>\frac1{2000}.
\]

To finish, differentiate with respect to the energy variable: \[
\mathcal D'(\tau)=-\frac{2s+3}{\mathcal C(\tau)^2}<0.
\] Thus \(\mathcal D(\tau)\ge\mathcal D(T_0)>1/v+1/2000\) for \(0\le\tau\le T_0\). Substitution into \(\mathcal L_h\) proves the strict affine lower bound. At \(\tau=0\), the affine comparison is an equality.

\end{proof}

\begin{lemma}[T.2 A pole-free comparison preserving the shared energy]\label{lem-terminal-comparison}

On \(\mathcal D_*\cap\{q\le p\}\) from \ref{def-analytic}, let \[
e=\frac{3a}{(h+2)S_0},\quad
k=e-\frac{a^2}{2(h+2)^2}=\frac K\Pi,\quad
A_- =\frac{493a^2}{100(h+2)^2},\quad
\overline T=\frac{987}{100}h^2e.
\] Use \(H_1,H_2\) from \ref{def-terminal-rationals} and \(\mathcal L_h\) from \ref{lem-common-tail-line}. All functions below are well defined, and \[
\begin{aligned}
\mathcal G(h,a,b)&>H_2(k)+\mathcal B(A_-,w)+\mathcal L_h(\overline T),\\
\mathcal G(h,a,b)&>H_1(k)+\mathcal B(A,w)+\mathcal L_h(\overline T).
\end{aligned}
\]

\end{lemma}

\begin{proof}[Proof of \ref{lem-terminal-comparison}]

Use the notation of \ref{def-analytic}. By \ref{lem-pole-free}, \(0<K<\Pi\) and \(0<A<1\), and the initial radicals and denominators are positive. By \ref{lem-pi-rational}, \(\Pi_-=493/50<\Pi<987/100=\Pi_+\). Since \(S_0\ge1+ha\) and \(a\le1\), \[
0<e\le\frac3{L(h+2)}=:e_0,
\qquad 0<T=\Pi h^2e<\overline T\le T_0=\Pi_+h^2e_0.
\] The parameter \(k\) belongs to \((0,1)\), so the denominators of \(H_1,H_2\) are strictly positive. Writing out the head response and using \(\Pi<\Pi_+\) gives \[
\mathsf H(\Pi k)
=1-\frac{2k}{1-k}-\frac{2k}{4-k}
-\frac{(\Pi/3-5/2)k}{1-k/9}>H_2(k)>H_1(k).
\] The second strict comparison follows from \[
H_2(k)-H_1(k)=\frac{79k}{100}
\left(\frac1{1-k/4}-\frac1{1-k/9}\right)>0.
\] Also \(0<A_-<A<1\) and \(d=1-w>0\), whence \[
\partial_A\mathcal B(A,w)=\frac d{2w}+
\frac{(1-A)d^2}{10w}>0.
\] Therefore \(\mathcal B(A,w)>\mathcal B(A_-,w)>0\).

We next compare the tails corresponding to the same \(e\). Since \(q\le p\), we have \(m=p>0\), and hence \(R=L/(h+2)\). Take \(s=s(h)>1\) from \ref{lem-common-tail-line}. Completing the square yields the identity \[
\frac{\Pi_-R^2}{2}-s(s-1)
=\frac{13}{25}\left(R-\frac{1701}{2600}\right)^2
+\frac{15427}{13000000}>0
\] which, together with \(\sigma(\sigma-1)=\Pi R^2/2\), gives \(\sigma>s\).

Let \(B=2s+5\). By \ref{lem-common-tail-line}, \(T_0/B^2<30/169\). For \(0\le t\le30/169\), \[
(1-4t)-(1-13t/5)^2=t(6/5-169t/25)\ge0,
\] and the proposed lower bound for the square root is positive. Therefore \[
\sqrt{B^2-4\overline T}\ge B-\frac{13\overline T}{5B}>0.
\] The original algebraic response at \((s,\overline T)\) consequently satisfies \(C\ge\mathcal C(\overline T)>0\), \(D\ge\mathcal D(\overline T)>0\), and hence \[
\frac{\mathsf T(s,\overline T)}h\ge\mathcal L_h(\overline T).
\] Starting at \((s,\overline T)\), increase the order to \(\sigma\) and then decrease the energy to \(T\). The radical, \(C\), and \(D\) are nondecreasing along this path, so positivity is preserved throughout. The general tail monotonicity in \ref{lem-derivatives} gives \[
\frac{\mathsf T(\sigma,T)}h>
\frac{\mathsf T(s,\overline T)}h\ge\mathcal L_h(\overline T).
\] Adding the head and gain comparisons proves both inequalities.

\end{proof}

\begin{lemma}[T.3 Reciprocal comparisons and curvature on the first terminal curve]\label{lem-psi-payment}

Let \(3/8\le h\le4\). Use \(L,s,v,w,d,e,k,A_-,H_2\) and \(J,G_*,\mathcal L_h,Z_\Psi\) from \ref{def-terminal-rationals}, and \(\mathcal B\) from \ref{def-analytic}. All denominators used are nonzero; moreover \(0<k<1/L\), \(Z_\Psi>27/16000\), and \[
H_2(k)+\mathcal B(A_-,w)+\mathcal L_h((987/100)h^2e)
>\frac{Z_\Psi}{h}>\frac{27}{16000h}>0.
\]

\end{lemma}

\begin{proof}[Proof of \ref{lem-psi-payment}]

This lemma concerns only the rational functions in \ref{def-terminal-rationals}. Write \(A=A_-\), \(e_0=3/[L(h+2)]\), and \(\Pi_+=987/100\). Common-denominator calculations give \(d=1-w\) and \[
\frac e{e_0}=1-\frac8{h(60h+127)}\in(0,1).
\] Indeed, \(h(60h+127)-8\) is strictly increasing and equals \(769/16>0\) at \(h=3/8\). Thus \(0<\Pi_+h^2e<T_0\), and \ref{lem-common-tail-line} applies.

The head energy is \[
k=\frac{300h^3+1295h^2+1349h-96}
{2hL(h+2)^2(60h+127)}>0.
\] Its numerator is strictly increasing and equals \(77799/128>0\) at \(h=3/8\). Furthermore, \[
\frac1L-k=\frac{P_4(h)}{2hL(h+2)^2(60h+127)},\qquad
P_4(h)=120h^4+434h^3+201h^2-333h+96.
\] Since \[
P_4''(h)=1440h^2+2604h+402>0,\quad
P_4'(3/8)=\frac{837}{32}>0,\quad
P_4(3/8)=\frac{12621}{512}>0,
\] we obtain \(0<k<1/L\le8/11<3/4\). Every head denominator is strictly positive. For \(0\le t\le3/4\), \[
\frac14-\frac1{8(1-t/4)}-\frac{79}{900(1-t/9)}
=\frac{25t^2-196t+134}{100(t-9)(t-4)}>0.
\] The denominator is positive; the numerator is decreasing on this interval and equals \(17/16\) at \(t=3/4\). Consequently, \[
H_2(k)>1-\frac{2k}{1-k}-\frac{129k}{100}-\frac{k^2}{4}
\ge1-\frac{2k}{1-k}-\left(\frac{129}{100}+\frac1{4L}\right)k=J.
\]

For the gain, put \(W=60h^3+143h^2+16h+8>0\). We have \[
\frac dw-\frac{26h-5}{30h^2}
=\frac{P_3(h)}{30h^2W},\qquad P_3(h)=152h^3+59h^2-128h+40,
\] and the tangent-line identity \[
P_3(h)-\left(45h-\frac{51}{4}\right)
=\frac{(2h-1)^2(152h+211)}4\ge0
\] together with \(45h-51/4\ge33/8\) gives \(P_3>0\). Moreover, \[
d-\frac{11}{15L}=\frac{4(2h^2+5h-2)}{15hL(4h+9)}>0,
\] because \(2h^2+5h-2\) is strictly increasing and equals \(5/32\) at \(h=3/8\). Thus \(d>0\), \(w>0\), and \(A\le7888/9025<7/8\). Using the full gain from \ref{def-analytic}, we obtain \[
\begin{aligned}
\mathcal B(A,w)
&=\frac{Ad}{2w}\left(1+\frac{(2-A)d}{10}\right)\\
&>\frac{A(26h-5)}{60h^2}
\left(1+\frac{33}{400L}\right)=G_*.
\end{aligned}
\] Both gain terms are retained together in this estimate.

\textbf{A common reciprocal form.} Keep the same \(e\) in the head and tail. Put \[
N=300h^3+1295h^2+1349h-96,\qquad
D=2hL(h+2)^2(60h+127).
\] Thus \(k=N/D\), and \(D-N>0\) has already been established. Define \[
S(h)=\frac{2hk}{1-k}=\frac{2hN}{D-N},\qquad R_0(h)=Z_\Psi(h)+S(h).
\] Elementary partial fractions of the original response give \[
\begin{aligned}
R_0(h)={}&\frac{389h}{3350}-\frac{25913}{33500}
-\frac{b_0}{h}+\frac{b_1}{(h+1)^2}-\frac{b_2}{h+1}\\
&+\frac{b_3}{(h+2)^2}+\frac{b_4}{h+2}-\frac{b_5}{h+127/60},
\end{aligned}
\tag{T.3.1}
\] where all six constants are positive: \[
\begin{gathered}
b_0=\frac{213469}{1920000},\quad b_1=\frac{225}{268},\quad
b_2=\frac{11452440743}{3591200000},\\
b_3=\frac{1773689}{1600000},\quad
b_4=\frac{74489644841}{4502400000},\quad
b_5=\frac{23869857}{3927875}.
\end{gathered}
\] We bound \(S\) with one reciprocal majorant and control the resulting difference by convexity.

\textbf{One majorant on the full interval.} Use the single majorant

\[
M(h)=\frac9{2(h+2)}+\frac6{5(h+2)^2}-\frac1{9h},\qquad G(h)=R_0(h)-M(h).
\]

Since \(h>0\) and \(D-N>0\), the weighted gap

\[
\mathcal W(h)=\frac{15(h+2)^2(D-N)}{2h}\,[M(h)-S(h)]
\]

has the same sign as \(M-S\). Combining the three reciprocal terms of \(M\) gives the unexpanded formula

\[
\mathcal W(h)=\frac{(395h^2+878h-40)(D-N)}{12h^2}
-15(h+2)^2N.
\]

Three differentiations yield

\[
\mathcal W'''(h)=\frac{7680}{h^5}-\frac{48804}{h^4}
+58071-12820h-33000h^2.
\]

For \(h\ge3/8\), we have \(7680/h\le20480\) and \(58071-12820h<59000\). Thus

\[
h^4\mathcal W'''(h)<-28324+h^4(59000-33000h^2).
\]

The scalar function \(g(y)=y^2(59000-33000y)\) for \(y\ge0\) has derivative \(g'(y)=y(118000-99000y)\). It increases up to \(y=118/99\) and decreases afterwards, so

\[
g(y)\le\frac{4\cdot59000^3}{27\cdot33000^2}
=\frac{821516000}{29403}<28000.
\]

Applying this with \(y=h^2\) gives

\[
\mathcal W'''(h)<-\frac{324}{h^4}<0.
\]

Consequently \(\mathcal W''\) is strictly decreasing and \(\mathcal W'\) is strictly concave. On the lower interval \(3/8\le h\le2\),

\[
\mathcal W''(h)\ge\mathcal W''(2)=\frac{140594}{3}>45000.
\]

At \(h=1/2\), the value and slope are

\[
\mathcal W(1/2)=\frac{5979}{32},\qquad
\mathcal W'(1/2)=-\frac{11609}{4}.
\]

The function \(\mathcal W(h)-22500(h-1/2)^2\) is strictly convex. Its tangent at \(1/2\), with \(z=h-1/2\), gives

\[
\begin{aligned}
\mathcal W(h)
&\ge\frac{5979}{32}-\frac{11609}{4}z+22500z^2\\
&\ge\frac{5979}{32}-\frac{(11609/4)^2}{90000}
=\frac{134286119}{1440000}>90.
\end{aligned}
\]

For the upper interval \(2\le h\le4\), use

\[
\mathcal W'(2)=\frac{537851}{4}>0,\qquad
\mathcal W(2)=113604,\qquad \mathcal W(4)=227460.
\]

The positive set of the concave function \(\mathcal W'\) is an interval containing \(2\). Once \(\mathcal W'\) becomes nonpositive, it cannot become positive again: otherwise concavity between \(2\) and a later point with positive derivative would force positivity at the intervening point. Thus \(\mathcal W\) first increases and may then decrease, and its minimum on \([2,4]\) occurs at an endpoint. Both endpoint values are positive, so \(\mathcal W>0\) also on this upper interval.

Finally,

\[
M-S=\frac{2h}{15(h+2)^2(D-N)}\mathcal W(h)>0.
\]

Hence \(S<M\) and \(Z_\Psi>G\) throughout \([3/8,4]\).

\textbf{A uniform curvature bound.}

Put

\[
x=\frac1{h+2}\in\left[\frac16,\frac8{19}\right],\qquad
\epsilon=b_0-\frac19=\frac{407}{5760000}>0,\qquad C=b_3-\frac65<0.
\]

Differentiating the retained partial fractions gives

\[
\frac{(h+2)^3G''(h)}2=Q(x),
\]

\[
Q(x)=b_4-\frac92+3Cx-f(x)
-\frac{b_5}{(1+7x/60)^3}-\frac{\epsilon}{(1-2x)^3},
\qquad
f(x)=\frac{b_2(1-x)-3b_1x}{(1-x)^4}.
\]

All denominators are positive on this interval. We will prove

\[
Q(x)>-\frac12\quad\left(\frac16\le x\le\frac8{19}\right),
\qquad
Q(x)>0\quad\left(\frac16\le x\le\frac27\right).
\]

First consider \(1/6\le x\le2/5\). We have

\[
f'(x)=\frac{3\{(b_2-b_1)-(b_2+3b_1)x\}}{(1-x)^5},
\qquad
f''(x)=\frac{12\{(b_2-2b_1)-(b_2+3b_1)x\}}{(1-x)^6}.
\]

The numerator controlling \(f''\) is decreasing and affine. Thus \(f'\) either is monotone or first increases and then decreases; its minimum on a compact interval occurs at an endpoint. The bounds \(b_2>3\) and \(b_1<1\) give

\[
f'(1/6)>\frac{3}{(5/6)^5}>3.
\]

At the other endpoint, use \(b_2>637/200\) and \(b_1<21/25\):

\[
f'(2/5)=\frac{625}{81}(3b_2-11b_1)
>\frac{625}{81}\left(3\frac{637}{200}-11\frac{21}{25}\right)
=\frac{175}{72}>\frac{12}{5}.
\]

Consequently \(f'>12/5\) on the whole interval. Since \(C<0\), \(\epsilon>0\), and \(b_5<61/10\), differentiation of \(Q\) yields

\[
Q'(x)<-\frac{12}{5}+\frac7{20}b_5
<-\frac{12}{5}+\frac{427}{200}=-\frac{53}{200}<0.
\]

For \(2/5\le x\le8/19\), bound the coupled head term by its elementary scalar maximum. For positive \(a,b\),

\[
\frac{d}{dx}\frac{a-(a+3b)x}{(1-x)^4}
=\frac{3\{a-b-(a+3b)x\}}{(1-x)^5},\qquad x<1.
\]

It increases up to \(x=(a-b)/(a+3b)\) and decreases afterwards, so its global maximum for \(x<1\) is \((a+3b)^4/(256b^3)\). Use

\[
b_2<\frac{319}{100},\qquad b_1>\frac{67}{80}.
\]

Since \(0<x<1\), these inequalities increase the numerator defining \(f\). They give

\[
f(x)<\frac{2281^4}{12800000\,67^3}<\frac{176}{25}.
\]

The other three bounds are

\[
\begin{aligned}
b_4-\frac92+3Cx
&\ge b_4-\frac92+\frac{24C}{19}
=\frac{1020466792283}{85545600000}>\frac{119}{10},\\
\frac{b_5}{(1+7x/60)^3}
&\le b_5\left(\frac{150}{157}\right)^3
=\frac{644486139000}{121603647739}<\frac{16}{3},\\
\frac{\epsilon}{(1-2x)^3}
&\le\epsilon\left(\frac{19}{3}\right)^3
=\frac{2791613}{155520000}<\frac1{50}.
\end{aligned}
\]

Thus

\[
Q(x)>\frac{119}{10}-\frac{176}{25}-\frac{16}{3}-\frac1{50}
=-\frac{37}{75}>-\frac12.
\]

In particular this holds at \(x=2/5\). Since \(Q\) is decreasing on \([1/6,2/5]\), the same lower bound extends to the full interval \([1/6,8/19]\).

To obtain positive curvature in the upper height range, it remains to bound \(Q(2/7)\). The affine contribution and the small-pole bound above hold on all of \([1/6,8/19]\). In addition,

\[
\begin{aligned}
f(2/7)&=\frac{343}{625}(5b_2-6b_1)
<\frac{343}{625}\left(5\frac{319}{100}-6\frac{67}{80}\right)
=\frac{149891}{25000}<6,\\
\frac{b_5}{(1+7(2/7)/60)^3}
&=b_5\left(\frac{30}{31}\right)^3
<\frac{61}{10}\left(\frac{30}{31}\right)^3<\frac{28}{5}.
\end{aligned}
\]

Therefore

\[
Q(2/7)>\frac{119}{10}-6-\frac{28}{5}-\frac1{50}
=\frac7{25}>\frac14.
\]

Monotonicity gives \(Q>0\) on \([1/6,2/7]\), as required.

\textbf{One convex correction and one tangent.}

Define

\[
F(h)=G(h)+\frac1{2(h+2)}.
\]

The uniform bound just proved gives

\[
F''(h)=\frac{2}{(h+2)^3}\left(Q(x)+\frac12\right)>0
\quad\left(\frac38\le h\le4\right).
\]

Convexity and the tangent at \(h=1\) imply

\[
G(h)\ge G(1)+G'(1)(h-1)-\frac{(h-1)^2}{18(h+2)}.
\]

The scalar data are

\[
G(1)=\frac{208251137}{16401600000}>\frac1{80},
\qquad
-\frac1{400}<G'(1)=-\frac{42620371163}{18402595200000}<0.
\]

For \(3/8\le h\le3/2\), the negative contribution from the linear term is at most \(1/800\): it is nonnegative for \(h\le1\), while \(0\le h-1\le1/2\) otherwise. Also \(|h-1|\le5/8\) and \(h+2\ge19/8\), so

\[
\frac{(h-1)^2}{18(h+2)}\le\frac{(5/8)^2}{18(19/8)}=\frac{25}{2736}.
\]

It follows that

\[
G(h)>\frac1{80}-\frac1{800}-\frac{25}{2736}
=\frac{289}{136800}>\frac{27}{16000}.
\]

For \(h\ge3/2\), we have \(x\le2/7\) and therefore \(G''(h)>0\). Moreover,

\[
G'(3/2)=\frac{5115313959269}{1118040013125000}>0.
\]

Hence \(G\) is strictly increasing on \([3/2,4]\), and the same strict lower bound extends to that entire interval. Together with \(Z_\Psi>G\), this proves the original strict margin on the full original range.

\textbf{The original response.} Combining \ref{lem-common-tail-line} with the head and gain estimates above yields \[
H_2(k)+\mathcal B(A,w)+\mathcal L_h(\Pi_+h^2e)
>J+G_*+\frac{s-v\Pi_+h^2e}{h}=\frac{Z_\Psi}{h}>\frac{27}{16000h}>0.
\]

\end{proof}

\begin{lemma}[T.4 Convexity and one tangent on the second terminal curve]\label{lem-phi-payment}

Let \(\sqrt3-1<h\le4\) and \(L=1+h\), and put \[
e=\frac{3(13h-2)}{13hL(h+2)},\quad
z=\frac{13h-2}{15h(h+2)},\quad k=e-z^2/2,\quad r=15e/16.
\] Use \(H_1\) from \ref{def-terminal-rationals} and \(s,v,\mathcal L_h\) from \ref{lem-common-tail-line}. Set \[
J_\Phi(r)=1-\frac{2r}{1-r}-\frac{3r}{2},\qquad
Z_\Phi=s+hJ_\Phi(r)-v\frac{987}{100}h^2e.
\] All denominators are strictly positive, and \[
H_1(k)+\mathcal L_h((987/100)h^2e)>\frac{Z_\Phi}{h}>\frac{3}{500h}>0.
\]

\end{lemma}

\begin{proof}[Proof of \ref{lem-phi-payment}]

Put \(\Pi_+=987/100\) and \(e_0=3/[L(h+2)]\). Since \(h>\sqrt3-1>2/3\), \[
0<\frac{z^2}{2e}=\frac{13(13h-2)L}{1350h(h+2)}<\frac{169}{1350}<1,
\] we have \(e,k>0\). Also \(e_0-e=6/[13hL(h+2)]>0\), so \(0<\Pi_+h^2e<T_0\); \ref{lem-common-tail-line} guarantees positivity of all tail denominators. The inequality \(3>(19/11)^2\) gives \(h>8/11\). In the identity \[
\frac{z^2}{2e}-\frac1{16}
=\frac{677h^2-206h-208}{10800h(h+2)}
\] the numerator equals \(32/121\) at \(h=8/11\), and its derivative for \(h\ge8/11\) is at least \(8566/11>0\). Hence \(k<r=15e/16\). Furthermore, \[
\frac{27}{52}-e
=\frac{3(h+4)(3h-2)(3h-1)}{52hL(h+2)}>0.
\] Thus \(0<k<r<405/832<1/2\), and every head denominator is positive. The function \(H_1\) is strictly decreasing on this interval, and \[
\frac{129}{100(1-r/4)}<\frac{258}{175}<\frac32,
\] so \(H_1(k)>J_\Phi(r)\).

Retaining the same \(e\), set \[
Q=208h^3+624h^2-169h+90=208hL(h+2)(1-r)>0.
\] The only non-elementary-looking reciprocal in \(Z_\Phi\) admits a simple majorant: \[
\begin{aligned}
\frac{45}{8(h+3)}-\frac{2hr}{1-r}
&=\frac{45(32h^2-73h+90)}{8(h+3)Q}>0,\\
32h^2-73h+90
&=32\left(h-\frac{73}{64}\right)^2+\frac{6191}{128}>0.
\end{aligned}
\] The identity follows by subtracting \(16h(h+3)(13h-2)\) from \(Q\). Substitution into \(Z_\Phi\) and elementary partial fractions give \[
\begin{aligned}
Z_\Phi(h)>G(h):={}&\frac{389h}{3350}-\frac{277649}{435500}
-\frac{174393}{139360(h+1)}\\
&+\frac{54705}{6968(h+2)}-\frac{45}{8(h+3)}.
\end{aligned}
\tag{T.4.7}
\] We prove positivity of \(G\) by curvature on the whole interval \([8/11,4]\). Put \[
c_1=\frac{174393}{139360}<\frac{21}{16},\qquad
c_2=\frac{54705}{6968}>\frac{15}{2},\qquad c_3=\frac{45}{8}.
\] For \(x=1/(h+2)\in[1/6,11/30]\), direct differentiation yields \[
\frac{(h+2)^3G''(h)}2
=c_2-f(x),\qquad
f(x)=\frac{c_1}{(1-x)^3}+\frac{c_3}{(1+x)^3}.
\] Both summands of \(f\) are strictly convex, so its maximum on this interval occurs at an endpoint. The elementary inequalities \[
\left(\frac65\right)^3<2,\quad
\left(\frac67\right)^3<\frac23,\quad
\left(\frac{30}{19}\right)^3<4,\quad
\left(\frac{30}{41}\right)^3<\frac25
\] give \[
f(1/6)<\frac{51}{8}<\frac{15}{2},\qquad
f(11/30)<4\frac{21}{16}+\frac25\frac{45}{8}=\frac{15}{2}.
\] Consequently \(G''>0\) throughout \([8/11,4]\).

One tangent now suffices. At \(h_*=9/4\), substitution gives \[
G(h_*)=\frac{9786437}{673718500}>\frac7{500},\qquad
0<G'(h_*)=\frac{209699639}{52112125975}<\frac1{200}.
\] Since \(G\) is convex and \(G'(h_*)>0\), every \(h\in[8/11,4]\) satisfies \[
\begin{aligned}
G(h)&\ge G(h_*)+G'(h_*)(h-h_*)\\
&\ge G(h_*)-\frac{67}{44}G'(h_*)
>\frac7{500}-\frac{67}{8800}
=\frac{281}{44000}>\frac3{500}.
\end{aligned}
\] Thus \(Z_\Phi>3/500\). Finally, \ref{lem-common-tail-line} gives \[
H_1(k)+\mathcal L_h(\Pi_+h^2e)
>J_\Phi(r)+\frac{s-v\Pi_+h^2e}{h}=\frac{Z_\Phi}{h}>\frac3{500h}>0.
\]

\end{proof}

\begin{lemma}[Strict positivity on the two terminal curves]\label{lem-terminal}

For the \(\Psi\) parameters in \ref{def-analytic} that belong to \(\mathcal D_*\) and satisfy \(q_c(h,1)\le h/(1+h)\), we have \(\Psi(h)>0\). For \(\Phi\) on \(\sqrt3-1<h\le4\), we have \(\Phi(h)>0\). Consequently \(\mathcal G>0\) throughout \(\mathcal D\) from that definition.

\end{lemma}

\begin{proof}[Proof of \ref{lem-terminal}]

Use the notation of \ref{def-analytic} and the constant bounds in \ref{lem-pi-rational}, and put \(\Pi_+=987/100\).

First consider \(\Psi\). Here \(a=1\) and \(p=h/L\), and the container equality \(S_0d=c\) is equivalent to \[
1-c=q_c\left(h+\frac{p}{1-q_c}\right).
\] The right-hand side is strictly increasing for \(0<q_c<1\). Since \[
\frac2{15}<1-c<\frac17,\qquad
\frac15\left(h+\frac54p\right)\ge\frac{63}{440}>\frac17
\quad(h\ge3/8),
\] we obtain \(q_c<1/5\), and then \[
q_c>\frac{2/15}{h+5p/4}
=\frac{8L}{15h(4h+9)}=:q_0(h)>0.
\] At the geometrically admissible heights in the statement, \(q_0<q_c\le p\). When \(q\) decreases from \(q_c\) to \(q_0\), \(w\) does not increase and \(a(1-q)\) increases. The entire segment therefore remains in \(\mathcal D_*\cap\{q\le p\}\). By \ref{lem-q-branch-monotone}, \[
\Psi(h)>\mathcal G\left(h,1,\frac Lh q_0(h)\right).
\] Direct calculation at this endpoint gives \[
\begin{aligned}
w&=\frac{60h^3+143h^2+16h+8}{15hL(4h+9)},\qquad
S_0=\frac{hL(60h+127)}{60h^2+127h-8},\\
e&=\frac{3(60h^2+127h-8)}{hL(h+2)(60h+127)},\qquad
k=e-\frac1{2(h+2)^2},\qquad A_- =\frac{493}{100(h+2)^2}.
\end{aligned}
\] These are precisely the \(\Psi\) quantities of \ref{def-terminal-rationals}. Thus \ref{lem-terminal-comparison} and \ref{lem-psi-payment} give \[
\Psi(h)>H_2(k)+\mathcal B(A_-,w)+\mathcal L_h(\Pi_+h^2e)>0.
\] The rational lemma holds on \([3/8,4]\); the geometric path above is used only within \(\mathcal D_*\) as required in the statement. Points introduced by rational continuation are not treated as actual pentagons.

Next consider \(\Phi\). If \(\sqrt3-1<h\le4\), then \(L>\sqrt3\). Set \(a_\triangle=(cL-1)/h\). We have \(0<a_\triangle<1\), and \[
a_\triangle>\frac1{C_h}\quad\Longleftrightarrow\quad
c(L^2+1)-2L=c(L-\sqrt3)(L-1/\sqrt3)>0.
\] Thus \((h,a_\triangle,0)\in\mathcal D_*\), and \(\Phi\) is defined on the entire stated interval. Define \[
a_1(h)=\frac{13h-2}{15h},\qquad
a_1-a_\triangle=\left(\frac{13}{15}-c\right)\frac Lh>0.
\] We have \(a_\triangle<a_1<1\). Increasing \(a\) to \(a_1\) while keeping \(q=0\) preserves \(a>1/C_h\) and \(p=ha/L<1\), so the whole segment lies in \(\mathcal D_*\). By \ref{lem-zero-monotone}, \[
\Phi(h)>\mathcal G(h,a_1(h),0).
\] At this endpoint, \[
S_0=\frac{13L}{15},\qquad
e=\frac{3(13h-2)}{13hL(h+2)},\qquad
\rho_+=\frac{13h-2}{15h(h+2)},\qquad k=e-\rho_+^2/2.
\] These are exactly the parameters of \ref{lem-phi-payment}. The second inequality of \ref{lem-terminal-comparison} therefore gives \[
\Phi(h)>H_1(k)+\mathcal L_h(\Pi_+h^2e)+\mathcal B(A,w)
>\frac{Z_\Phi}{h}+\mathcal B(A,w)>0.
\] Both original gain terms remain as an additional positive remainder.

Finally, take any parameter in \(\mathcal D\). By \ref{lem-boundary}, three branches cover this domain: if \(q>p\), then \(\mathcal G>1/150\); if \(q\le p\) and \(S_0d<c\), then \(\mathcal G>\Psi(h)>0\); and if \(q\le p\) and \(S_0/L<c\), then \(\mathcal G>\Phi(h)>0\). Since \(S_0\min(d,L^{-1})<c\), these branches cover all of \(\mathcal D\). Hence \(\mathcal G>0\).

\end{proof}


\section{Large heights: direct Sturm comparison}

\begin{definition}[H.3 The two-interval comparison operator]\label{def-high-operator}

For real \(2\le s\le3\) and \(\ell>0,b\ge0\), let \(K_{s;\ell,b}\) be the Friedrichs operator associated with \[
-\partial_z^2+s(s-1)z^{-2}\mathbf1_{(0,\ell)}
\quad\hbox{on }(0,\ell+b),
\] with a Dirichlet condition at the right endpoint. Its form domain is \(H_0^1(0,\ell+b)\), and the potential energy is finite on that domain. Functions in the operator domain and their first derivatives are continuous across the interface. The free interval is empty when \(b=0\).

\end{definition}

\begin{proof}[Proof of \ref{def-high-operator}]

Let \(J=(0,\ell+b)\) and \(V(z)=s(s-1)z^{-2}\mathbf1_{(0,\ell)}(z)\). The space \(C_c^\infty(J)\) is dense in \(L^2(J)\), and every one of its functions has finite potential energy. For a real compactly supported smooth \(u\), integration by parts followed by Cauchy--Schwarz gives

\[
\int_J\frac{u^2}{z^2}=2\int_J\frac{uu'}z
\le2\left(\int_J\frac{u^2}{z^2}\right)^{1/2}\left(\int_J|u'|^2\right)^{1/2},
\quad \int_Ju^2/z^2\le4\int_J|u'|^2.
\]

Applying the same estimate to the real and imaginary parts covers complex functions. Approximation in \(H_0^1\) extends this Hardy inequality to every \(u\in H_0^1(J)\). The form domain is therefore exactly \(H_0^1(J)\); its form norm, including the potential, is equivalent to the Sobolev norm. The form is closed and nonnegative, and \(C_c^\infty(J)\) is a form core. The representation theorem for densely defined closed nonnegative symmetric forms (Kato, Chapter VI, Section 2 \cite{nc:kato}) gives a unique positive self-adjoint operator, which is the Friedrichs extension used here.

Compact embedding on the finite interval gives a compact resolvent. The nonnegative potential and Dirichlet Poincare inequality give a strictly positive first eigenvalue; min--max yields \(\lambda_n\ge n^2\pi^2/(\ell+b)^2\) for the full spectrum. When \(b>0\), the interface conditions concern functions in the operator domain. For any \(u\in D(K_{s;\ell,b})\), the representation theorem provides \(f=K_{s;\ell,b}u\in L^2(J)\) such that \(\int(u'\psi'+Vu\psi)=\int f\psi\) for every form-domain test function. Taking \(\psi\in C_c^\infty(I)\) on an open neighborhood \(I\) of the interface that avoids zero gives the distributional equation \(-u''+Vu=f\). Since \(V\) is bounded there and \(u\in H^1(I)\), we have \(u''=Vu-f\in L^2(I)\), hence \(u\in H^2(I)\). One-dimensional Sobolev embedding makes \(u,u'\) continuous across the whole interface neighborhood. Values and first derivatives therefore match, with no additional point potential. This applies to the entire operator domain, not merely to eigenfunctions. The right zero trace comes from \(H_0^1\); the left endpoint is determined by closure of the smooth core, with no additional classical derivative condition. When \(b=0\), the interface is the right endpoint and has Dirichlet zero trace; the free interval is empty.

\end{proof}

\begin{lemma}[H.2 The area cost of the actual crossing point]\label{lem-high-area-cost}

Use \ref{def-candidate} and set \(\delta=1/h\) and \(\eta=ph/S\), where \(h\ge4\). If \(p<q\), then \(h>k\), the tails intersect at \(x_c=hk(q-p)/(hq-kp)\), \(\alpha=x_c/h\in(0,1/2)\), and \(\eta<[1+4\alpha(1-\alpha)+3\delta+\delta^2]^{-1}\). If \(p\ge q\), then \(\eta<(1+\delta)^{-2}\).

\end{lemma}

\begin{proof}[Proof of \ref{lem-high-area-cost}]

Use the rational bounds in \ref{lem-pi-rational}; the weaker bounds invoked below follow by inclusion and comparison of squares. Use the height ordering, the strict upper bound \(\rho<\beta/[h(h+2)]\), and the lower bound \(\rho\ge S/[2(1+h)^2]\) from \ref{def-candidate}.

\textbf{H.2 The actual crossing point incurs an area cost}

If \(p\ge q\), the long tail \(p(1-x/h)\) dominates throughout the tail interval. Suppose \(p<q\). First establish the area estimate without assuming that the height ordering is strict. The candidate inequalities \(S/(2L^2)\le\rho<\beta/[h(h+2)]\) imply \(S<2pL^3/[h(h+2)]=2ph+2p+2pL/[h(h+2)]\). Since \(p\le h/L\) and \(h\ge4\), the last remainder is at most \(2/(h+2)\le1/3<1\), so \(S<2ph+2p+1\). On the other hand, \(S=1+p+q+ph+qk\) and \(q>p\) imply \(S>1+2p+ph+qk\). Comparison yields \(qk<ph\). If \(h=k\), this would give \(q<p\), a contradiction. The given ordering \(h\ge k\) therefore gives \(h>k\). This argument does not mistake the permission to exchange the ears for an ordering assumption already imposed on the original data. The crossing point is

\[
x_c=\frac{hk(q-p)}{hq-kp},\qquad
\alpha=\frac{x_c}{h},\qquad v=\frac kh.
\tag{H.2.1}
\]

The denominator and location of the crossing point are valid as follows. The inequalities \(p>0\), \(q>p\), and \(h>k>0\) give \(hq-kp=p(h-k)+h(q-p)>0\), hence \(x_c>0\). Moreover, \(h(q-p)<hq-kp\) is equivalent to \(hp>kp\), so \(x_c<k\), that is, \(0<\alpha<v=k/h<1\). Since \(p\le h/(h+1)\), we also have \(2p(h+1)/[h(h+2)]\le2/(h+2)\le1/3<1\), which makes the remainder in the following area estimate strictly smaller than one.

The longitudinal lower bound and the outer-ear upper bound for a low-energy candidate give

\[
S<\frac{2p(h+1)^3}{h(h+2)}
=2ph+2p+\frac{2p(h+1)}{h(h+2)}<2ph+2p+1.
\]

Since \(S=1+p+q+ph+qk\) and \(q>p\), it follows that

\[
qk<ph.
\tag{H.2.2}
\]

The crossing equation also gives

\[
\frac{qk}{ph}=\frac{v^2(1-\alpha)}{v-\alpha}.
\tag{H.2.3}
\]

If \(\alpha\ge1/2\), then \(\alpha<v<1\) and

\[
v^2(1-\alpha)-(v-\alpha)
=(v-1)((1-\alpha)v-\alpha)\ge0,
\]

contradicting (H.2.2). Therefore

\[
0<\alpha<\frac12.
\tag{H.2.4}
\]

The inequality \(v-\alpha>0\) has already been checked. Since \(v^2-4\alpha(v-\alpha)=(v-2\alpha)^2\ge0\), division by \(v-\alpha\) gives the estimate used next.

Using \(v^2/(v-\alpha)\ge4\alpha\), \(1/p\ge1+\delta\), and \(q/p>1\), we obtain

\[
\begin{aligned}
\frac S{ph}
&=1+\frac{qk}{ph}+\delta\left(\frac1p+1+\frac qp\right)\\
&>1+4\alpha(1-\alpha)+3\delta+\delta^2
=:D(\alpha,\delta).
\end{aligned}
\tag{H.2.5}
\]

In the absence of a crossing, \(\beta\le h\) and \(\gamma>0\) give directly

\[
\eta=\frac{\beta h}{(h+1)(1+\beta+\gamma)}
<(1+\delta)^{-2}.
\tag{H.2.6}
\]

\end{proof}

\begin{lemma}[H.4 An elementary free-interval response bound]\label{lem-free-sturm-bound}

For every \(0<q\le4\), \[
\sqrt q\cot\sqrt q\ge1-\frac{q}{3-q/4}.
\tag{H.4.1}
\] All denominators are positive on this interval.

\end{lemma}

\begin{proof}[Proof of \ref{lem-free-sturm-bound}]

Use the positive comparison function \[
p_q(z)=z\left(1-\frac{qz^2}{12}\right)^2
\qquad(0<z\le1).
\] Direct differentiation gives the factored identity \[
-p_q''-qp_q=\frac{q^2z^3}{144}(4-qz^2)\ge0.
\tag{H.4.2}
\] The solution \(w(z)=\sin(\sqrt q\,z)/\sqrt q\) is positive on \((0,1]\), since \(q\le4<\pi^2\) by \ref{lem-pi-rational}. The Wronskian vanishes at zero, and \[
(w'p_q-wp_q')'=w(-p_q''-qp_q)\ge0.
\] Integrating from zero to one and dividing by the positive interface values gives \[
\sqrt q\cot\sqrt q=\frac{w'(1)}{w(1)}
\ge\frac{p_q'(1)}{p_q(1)}
=1-\frac{q}{3-q/4}.
\] Since \(q\le4\), the last denominator is at least \(2\).

\end{proof}

\begin{lemma}[H.5.1 Reduction by concavity of the area cost]\label{lem-high-cost-reduction}

Let \(0\le\alpha\le1/2\), \(0\le\delta\le1/4\), and \(\ell=1-\alpha,b=\alpha+\delta,D=1+4\alpha(1-\alpha)+3\delta+\delta^2,t=b/\ell\). Then \(0\le t\le3/2\) and \(D/\ell^2\ge F_0(t)\), where \(F_0(t)=1+3t+t^2\) on \([0,1/4]\) and \(F_0(t)=(9+138t+29t^2)/25\) on \([1/4,3/2]\). At each fixed ratio, the minimum is attained at \(\delta=\min(t,1/4)\).

\end{lemma}

\begin{proof}[Proof of \ref{lem-high-cost-reduction}]

Use the rational bounds in \ref{lem-pi-rational}; the weaker bounds invoked below follow by inclusion and comparison of squares.

\textbf{H.5.1 The concave quadratic area cost at a fixed ratio}

Put \(t=b/\ell\) and \(z=1/(1+\delta)\). Since \(\ell\ge1/2\), we have \(0\le t\le3/2\). Inverting these relations gives

\[
\alpha=\frac{t-\delta}{1+t},\qquad
\ell=\frac{1+\delta}{1+t},\qquad
\max\{0,(t-1)/2\}\le\delta\le\min\{t,1/4\}.
\]

Direct substitution yields

\[
Q:=\frac D{\ell^2}
=(1+t)^2(1+z-z^2)+4(1+t)z-4.
\tag{H.5.2}
\]

For fixed \(t\), the expression \(Q\) is a strictly concave quadratic in \(z\), while \(\delta\mapsto z\) is strictly monotone. Its minimum therefore occurs at an endpoint of the admissible \(\delta\) interval, including when that interval is a singleton. The four boundary segments give

\[
\begin{array}{c|c|c}
\text{Boundary}&Q&\text{Range}\\\hline
\delta=0&1+6t+t^2&0\le t\le1\\
\alpha=0&1+3t+t^2&0\le t\le1/4\\
\delta=1/4&(9+138t+29t^2)/25&1/4\le t\le3/2\\
\alpha=1/2&t^2+4t+3&1\le t\le3/2
\end{array}
\]

On \([0,1/4]\), \([1/4,1]\), and \([1,3/2]\), respectively, the value of \(Q\) at the lower admissible endpoint minus its value at the upper endpoint is

\[
3t,\qquad \frac{4(4-t)(1+t)}{25},\qquad
\frac{2(3-2t)(t+11)}{25},
\]

All three differences are nonnegative. Thus the exact minimum occurs at \(\delta=\min\{t,1/4\}\) and equals

\[
F_0(t)=\begin{cases}
1+3t+t^2,&0\le t\le1/4,\\
(9+138t+29t^2)/25,&1/4\le t\le3/2.
\end{cases}
\tag{H.5.3}
\]

\end{proof}

\begin{lemma}[H.5 Direct Sturm bounds throughout the crossing range]\label{lem-crossing-sturm}

Use \(K_{s;\ell,b}\) from \ref{def-high-operator} and \(F_0\) from \ref{lem-high-cost-reduction}. Put \[
C_s=\frac{25s-11}{2}.
\] For every real \(2\le s\le3\) and every \(0\le t\le3/2\), \[
\lambda_1(K_{s;1,t})>\frac{C_s}{F_0(t)}.
\tag{H.5.4}
\] Consequently, for the geometric parameters \(\ell,b,D\) of \ref{lem-high-cost-reduction}, \[
\lambda_1(K_{2;\ell,b})>\frac{39}{2D},\qquad
\lambda_1(K_{3;\ell,b})>\frac{32}{D}.
\]

\end{lemma}

\begin{proof}[Proof of \ref{lem-crossing-sturm}]

Use \ref{def-high-operator} and write \(F(t)=F_0(t)\) as defined in \ref{lem-high-cost-reduction}.

\textbf{1. A positive comparison function and a uniform pole-free bound.}

For the fixed order \(2\le s\le3\), use the positive function

\[
v_s(y)=y^s(1-y^2)e^{-s y^2/2}\qquad(0<y<1).
\]

Direct logarithmic differentiation gives

\[
-\frac{v_s''}{v_s}+\frac{s(s-1)}{y^2}
=s^2+5s+s^2(1-y^2)+\frac{2}{1-y^2}
\ge B_s:=s^2+5s+2\sqrt2\,s.
\tag{H.5.5}
\]

The inequality is AM--GM applied to the last two terms. For a compactly supported smooth function \(f\), integration by parts gives the ground-state identity \[
\int_0^1\left(|f'|^2+\frac{s(s-1)}{y^2}|f|^2\right)
=\int_0^1v_s^2\left|\left(\frac f{v_s}\right)'\right|^2
+\int_0^1\left(-\frac{v_s''}{v_s}+\frac{s(s-1)}{y^2}\right)|f|^2.
\] Discarding the first nonnegative term gives a lower bound by \(B_s\|f\|_2^2\). The Hardy inequality and form-core density in \ref{def-high-operator} extend that inequality to \(H_0^1(0,1)\). Thus the first Friedrichs--Dirichlet eigenvalue on the singular interval satisfies \(\lambda_D(s)\ge B_s\).

The inequality \(\sqrt2>45/32\) follows by squaring. Hence, for \(2\le s\le3\),

\[
B_s-C_s>
 s^2-\frac{75s}{16}+\frac{11}{2}
=\left(s-\frac{75}{32}\right)^2+\frac7{1024}>0.
\tag{H.5.6}
\]

The case \(t=0\) of (H.5.4) follows immediately.

For \(t>0\), write \(F=F_0(t)\) and set \(E=C_s/F\), so \(F>1\) and \(0<E<C_s<B_s\le\lambda_D(s)\). Apply \ref{lem-regular-tail} with \(a=s>1\) and \(T=E<\lambda_D(s)\). Its Volterra construction gives the unique normalized regular solution \(u\), which has finite form energy, is positive on \((0,1]\), and satisfies \[
u(z)=z^s(1+O(z^2)),\qquad
u'(z)=s z^{s-1}+O(z^{s+1}).
\] Write \(D_s(E)=u'(1)/u(1)\). The hypotheses needed for the regular solution and its positivity have all been verified by the preceding separated spectral bound.

Use the rescaled comparison function

\[
g(z)=z^s\left(1-\frac{z^2}{F}\right)e^{-s z^2/(2F)}\qquad(0<z\le1).
\]

Its effective potential is the expression before the inequality in (H.5.5), at \(y=z/\sqrt F\), divided by \(F\). It therefore strictly exceeds \(E\). The Wronskian identity

\[
(u'g-ug')'=u\left[-g''+\left(\frac{s(s-1)}{z^2}-E\right)g\right]>0
\]

and the regular asymptotics at zero give

\[
D_s(E)>\frac{g'(1)}{g(1)}
=s\left(1-\frac1F\right)-\frac2{F-1}.
\tag{H.5.7}
\]

All functions have positive value at the interface. The displayed regular asymptotics of \(u\) and the explicit form of \(g\) give \(u'g-ug'=O(z^{2s+1})\to0\), which justifies integration at the singular endpoint.

Let \(q=Et^2=C_st^2/F\). The quantity \(t^2/F(t)\) increases on either quadratic piece, since its derivative is \(t(2a+bt)/F^2\) when \(F=a+bt+ct^2\). The rational bounds in \ref{lem-pi-rational} therefore give

\[
0<q\le\frac{32}{29}<4\quad(0<t\le1/4),
\qquad
0<q\le\frac{32}{5}<\pi^2\quad(1/4\le t\le3/2).
\tag{H.5.8}
\]

In particular, \(E<\pi^2/t^2\), so the free segment also stays below its first separated Dirichlet eigenvalue. Its outward response at the shared physical energy is

\[
D_{\mathrm{free}}(E,t)=\sqrt E\cot(\sqrt E\,t)
=\frac1t\sqrt q\cot\sqrt q.
\]

It remains to prove that the sum of the two outward responses is positive.

\textbf{2. The first interval: a single concave quadratic bound.}

By \ref{lem-free-sturm-bound}, \[
\sqrt q\cot\sqrt q\ge1-\frac{q}{3-q/4}
\qquad(0<q\le4).
\tag{H.5.9}
\]

For \(0<t\le1/4\), put \(F=1+3t+t^2\). Combining (H.5.7) and (H.5.9) gives

\[
\begin{aligned}
t(D_s+D_{\mathrm{free}})
&>1-\frac2{3+t}+\frac{s t(F-1)}F
-\frac{C_s t^2}{3F-C_st^2/4}.
\end{aligned}
\]

Here \(F\le29/16\), while \(3F-C_st^2/4\ge3+9t-5t^2\ge3\). Therefore

\[
\begin{aligned}
t(D_s+D_{\mathrm{free}})
&>\frac13+\frac{8t}{39}
-\left(\frac{C_s}{3}-\frac{48s}{29}\right)t^2\\
&\ge P(t):=\frac13+\frac{8t}{39}-\frac{496t^2}{87}.
\end{aligned}
\tag{H.5.10}
\]

The last step uses

\[
\frac{C_s}{3}-\frac{48s}{29}
=\frac{437s-319}{174}\le\frac{496}{87}.
\]

The quadratic \(P\) is concave. Its endpoint values are

\[
P(0)=\frac13,\qquad P(1/4)=\frac{32}{1131}>0.
\]

Consequently the total interface response is strictly positive on the entire first interval, uniformly in \(2\le s\le3\).

\textbf{3. A simple rational lower bound for the free response.}

For the second interval use the rational response function \(\mathsf H\) from \ref{def-analytic}. The Sturm comparison with two zero endpoints in \ref{lem-cot-rational} gives \(\sqrt q\cot\sqrt q>\mathsf H(q)\), and \ref{lem-pi-rational} gives \(493/50<\Pi=\pi^2<987/100\). Put \(k=q/\Pi\in(0,1)\). Then

\[
\begin{aligned}
\mathsf H(q)
&=1-\frac{2k}{1-k}-\frac{k/2}{1-k/4}
-\frac{(\Pi/3-5/2)k}{1-k/9}\\
&>1-\frac{2k}{1-k}-\frac{129k}{100(1-k/4)}.
\end{aligned}
\]

Indeed, \(0<\Pi/3-5/2<79/100\) and \(1-k/9>1-k/4>0\). The last rational function decreases strictly in \(k\). Since \(\Pi>49/5\) and \(129/100<13/10\), it follows that

\[
\sqrt q\cot\sqrt q>
1-\frac{2q}{49/5-q}-\frac{13q/10}{49/5-q/4}
\qquad\left(0<q\le\frac{32}{5}\right).
\tag{H.5.11}
\]

All denominators are positive. The estimate follows from the scalar Sturm comparison in \ref{lem-cot-rational}.

\textbf{4. The second interval: concavity of the response lower bound.}

Throughout \(1/4\le t\le3/2\), write

\[
F=a+bt+ct^2,\qquad a=\frac9{25},\ b=\frac{138}{25},\ c=\frac{29}{25}.
\]

Define

\[
A_s(t)=1+s t-\frac{s t}{F}-\frac{2t}{F-1},
\]

\[
B_C(t)=\frac{10C t^2}{49F-5Ct^2}
+\frac{(13C/2)t^2}{49F-(5C/4)t^2},
\qquad G_s(t)=A_s(t)-B_{C_s}(t).
\tag{H.5.12}
\]

Equations (H.5.7) and (H.5.11) give \(t(D_s+D_{\mathrm{free}})>G_s(t)\).

\textbf{4a. The singular-side expression is strictly concave.}

For a quadratic \(Q=a+bt+ct^2\),

\[
\left(\frac tQ\right)''
=-\frac{2[ab+3ac t-c^2t^3]}{Q^3}.
\]

Put \(H=b+3ct>0\) and \(X=c^2t^3>0\). Then

\[
A_s''=
\frac{2s(aH-X)}{F^3}
+\frac{4((a-1)H-X)}{(F-1)^3}.
\]

Since \(F>F-1>0\), \(2\le s\le3\), and

\[
2sa+4(a-1)\le10a-4=-\frac25,
\]

we obtain

\[
A_s''<-
\frac{2H}{5(F-1)^3}<0.
\tag{H.5.13}
\]

The omitted \(X\) terms have negative signs. The inequality \(10a-4<0\) comes from the constant term \(a=9/25\) in the area-cost quadratic.

\textbf{4b. The sum of the two free fractions is strictly convex.}

For \(39/2\le C\le32\), put

\[
A=49a,\quad B=49b,\quad
d=5C-49c,\quad e=49c-\frac{5C}{4},
\]

\[
Q_-=A+Bt-dt^2,\qquad Q_+=A+Bt+et^2.
\]

The bounds on \(q=Ct^2/F\) give \(Q_->0\) and \(Q_+>Q_-\). Moreover

\[
d\ge\frac{2033}{50}>40,\qquad
0<\frac{421}{25}\le e\le\frac{6493}{200}<33.
\]

In particular, \(d>e>0\). Set \(\alpha=10C\) and \(\beta=13C/2\), so \(\alpha>\beta>0\).

The second elementary identity is

\[
\left(\frac{t^2}{a+bt+ct^2}\right)''
=\frac{2[a^2-3ac t^2-bc t^3]}{(a+bt+ct^2)^3}.
\]

With \(M=A^2>0\) and \(N=3At^2+Bt^3>0\), it gives

\[
\frac12 B_C''
=\frac{\alpha(M+dN)}{Q_-^3}
+\frac{\beta(M-eN)}{Q_+^3}.
\]

Using \(Q_-<Q_+\) and \(\alpha d>\beta e\) yields

\[
\frac12 B_C''
\ge\frac{\alpha M}{Q_-^3}+\frac{\beta M}{Q_+^3}
+\frac{(\alpha d-\beta e)N}{Q_+^3}>0.
\tag{H.5.14}
\]

Thus \(G_s''<0\) for every \(2\le s\le3\) on the entire second interval.

\textbf{5. The endpoint comparisons are short and uniform.}

At \(t=1/4\), \(F=29/16\). The function \(B_C\) increases in \(C\), so \(C\le32\) gives

\[
B_C(1/4)\le\frac{320}{1261}+\frac{208}{1381}
<\frac13+\frac16=\frac12.
\]

Since \(s\ge2\),

\[
A_s(1/4)=1+\frac{13s}{116}-\frac8{13}
>\frac13+\frac{13}{58}.
\]

Therefore

\[
G_s(1/4)>\frac5{87}>0.
\tag{H.5.15}
\]

At each fixed \(t\), \(A_s\) is affine in \(s\). Both terms in \(B_C\) are strictly convex increasing functions of \(C\) wherever their denominators are positive. Since \(C_s\) is affine, \(G_s(t)\) is concave in \(s\). It is therefore enough to bound its two order endpoints at \(t=3/2\).

There \(F=45/4\). For \(s=2\), \(C_2=39/2\), and

\[
A_2(3/2)=\frac{2116}{615}>3,
\qquad
B_{39/2}(3/2)=\frac{78}{59}+\frac{1014}{1765}
<\frac32+\frac35=\frac{21}{10}.
\]

Hence \(G_2(3/2)>9/10\). For \(s=3\), \(C_3=32\), and the exact comparison is

\[
G_3(3/2)=\frac{1971}{410}-\frac{64}{17}-\frac{208}{205}
=\frac{39}{1394}>0.
\tag{H.5.16}
\]

Concavity in \(s\) gives \(G_s(3/2)\ge39/1394\). Concavity in \(t\), (H.5.15), and \(5/87>39/1394\) consequently imply

\[
G_s(t)\ge\frac{39}{1394}>0
\qquad\left(2\le s\le3,\ \frac14\le t\le\frac32\right).
\tag{H.5.17}
\]

\textbf{6. The common-interface conclusion.}

For every \(t>0\), the energy \(E=C_s/F(t)\) lies strictly below the first separated Dirichlet eigenvalue on both sides. The total outward response at the same physical interface and the same physical energy is strictly positive by Sections 2 and 4--5. To identify the ODE responses with those in \ref{lem-interface}, normalize \(u\) to have interface value one. Its regular asymptotics give finite form energy and zero outer trace. Subtracting a smooth trace-one lifting supported near the interface leaves an element of the side zero-trace form domain. Integration by parts against compactly supported smooth zero-trace tests gives the weak equation, and Hardy's inequality and density extend it to every zero-trace test. Since \(E<\lambda_D(s)\), the shifted form is strictly coercive on that zero-trace space. Any other function with the same trace has the form \(u/u(1)+z\) and shifted energy equal to the energy of \(u/u(1)\) plus the strictly positive remainder energy of \(z\). Thus the normalized ODE solution is the unique minimizing extension. The free solution has the same property by the ordinary interval equation. The full form has compact embedding by \ref{def-high-operator}, and the potential has bounded one-sided limits at the interface. Hence every hypothesis of the variational criterion \ref{lem-interface} is satisfied.

To match its stated orientation explicitly, use \(x=(1+t-z)/t\). This sends the free interval to \((0,1)\), the singular interval to \((1,1+1/t)\), and the interface to \(x=1\). The energy becomes \(t^2E=q\) on both pieces, and each outward response is multiplied by \(t\). Thus the sign and the separated spectral hypotheses are preserved. The criterion gives \(\lambda_1(K_{s;1,t})>E\). The case \(t=0\) was established directly in Section 1.

This proves (H.5.4). For the geometric parameters in \ref{lem-high-cost-reduction}, \(t=b/\ell\) lies in \([0,3/2]\) and \(D\ge\ell^2F_0(t)\). The change of variables \(z=\ell x\) gives \[
\lambda_1(K_{s;\ell,b})
=\ell^{-2}\lambda_1(K_{s;1,t})
>\frac{C_s}{\ell^2F_0(t)}
\ge\frac{C_s}{D}.
\] Since \(C_2=39/2\) and \(C_3=32\), this is the claimed pair of crossing bounds.

\end{proof}

\begin{lemma}[H.6 Direct spectral bounds from positive comparison functions]\label{lem-no-crossing-sturm}

Use \ref{def-high-operator}. For real \(2\le\sigma\le3\), define \[
\Lambda_\sigma(\delta)=(1+\delta)^2\lambda_1(K_{\sigma;1,\delta}).
\] The function \(\Lambda_\sigma\) is nonincreasing for \(\delta\ge0\). Throughout \(0\le\delta\le1/4\), \[
\Lambda_2(\delta)>\frac{625}{32}>\frac{39}{2},\qquad
\Lambda_3(\delta)>\frac{1525}{48}>\frac{1587}{50}.
\tag{H.6.1}
\]

\end{lemma}

\begin{proof}[Proof of \ref{lem-no-crossing-sturm}]

\textbf{Scaling and the common form domain.} Use \ref{def-high-operator} and put \(L_\delta=1+\delta\). The unitary map \(u(x)\mapsto L_\delta^{1/2}u(L_\delta z)\) transforms \(L_\delta^2K_{\sigma;1,\delta}\) into the form \[
\int_0^1\left(|u'|^2+\sigma(\sigma-1)z^{-2}
\mathbf1_{\{z<1/(1+\delta)\}}|u|^2\right)dz.
\] The Hardy inequality proved in \ref{def-high-operator} gives the common domain \(H_0^1(0,1)\) and equivalence of every form norm to the Sobolev norm. The potential is pointwise nonincreasing in \(\delta\), so the Rayleigh principle proves that \(\Lambda_\sigma\) is nonincreasing for every real \(2\le\sigma\le3\) and all \(\delta\ge0\). It suffices to prove the two strict bounds at \(\delta=1/4\). On the complete interval \((0,5/4)\), use the pairs \[
(a,E)=\left(2,\frac{25}{2}\right),\qquad
(a,E)=\left(3,\frac{61}{3}\right).
\tag{H.6.2}
\]

\textbf{One positive comparison family on the singular interval.} For either pair define \[
\begin{gathered}
A=2a+1,\quad \alpha=\frac EA,\quad
\beta=\frac{\alpha^2}{A+2},\\
\gamma=\frac{2\alpha}{A+4},\quad
c=\frac1{23},\quad
\theta=\frac{E}{(A+2)(A+4)}.
\end{gathered}
\] Set \[
\begin{gathered}
H(t)=1-\gamma t-ct^2,\qquad
Y(t)=\alpha+\frac{\beta t}{H(t)},\\
w(z)=z^a\exp\left[-\frac12\int_0^{z^2}Y(t)\,dt\right].
\end{gathered}
\tag{H.6.3}
\] The two values of \(\gamma\) are \(5/9\) and \(122/231\). Thus \(H(t)\ge83/207>0\) on \([0,1]\). The function \(w\) is positive on \((0,1]\), with \(w(z)=z^a(1+O(z^2))\) and \(w'(z)=az^{a-1}+O(z^{a+1})\). It has finite Sobolev and inverse-square potential energy near zero.

Logarithmic differentiation gives \(w'/w=a/z-zY(z^2)\). With \(t=z^2\), its complete differential residual is \[
\begin{aligned}
\frac{(-\partial_z^2+a(a-1)z^{-2}-E)w}{w}
&=AY+2tY'-tY^2-E\\
&=\frac{\alpha^2t^3}{H(t)^2}
\left[\frac{c(A+6)}{A+2}-(\theta+ct)^2\right].
\end{aligned}
\tag{H.6.4}
\] To check the factorization directly, use \[
H-tH'=1+ct^2,\qquad
(A+2)\beta=\alpha^2,\qquad
(A+4)\gamma=2\alpha.
\] After multiplication by \(H^2\), these identities reduce the residual to \(t^3\{\beta c(A+6)-(d+\alpha ct)^2\}\), where \(d=\alpha\gamma-\beta=A\beta/(A+4)=\alpha\theta\). This is the displayed square in (H.6.4).

The two values of \(\theta\) are \(25/126\) and \(61/297\), so throughout \(0\le t\le1\), \[
\theta+ct\le\frac{61}{297}+\frac1{23}<\frac14,
\qquad
\frac{c(A+6)}{A+2}\ge\frac{13}{207}
=\frac1{16}+\frac1{3312}.
\] Here \(1/4-61/297-1/23=31/27324>0\). Hence \[
\frac{(-\partial_z^2+a(a-1)z^{-2}-E)w}{w}
>\frac{\alpha^2z^6}{3312H(z^2)^2}>0
\qquad(0<z\le1).
\tag{H.6.5}
\]

\textbf{The free polynomial and the complete comparison function.} Put \(q=E/16\), which is \(25/32\) or \(61/48\), and define \[
p_q(x)=x\left(1-\frac{qx^2}{12}\right)^2.
\] Since \(0<q<4\), this polynomial is positive on \((0,1]\) and satisfies \[
-p_q''-q p_q=\frac{q^2x^3}{144}(4-qx^2)>0.
\tag{H.6.6}
\] Use the reversed coordinate \(x=5-4z\) on the free interval and set \[
g(z)=
\begin{cases}
w(z)/w(1),&0<z\le1,\\[2mm]
p_q(5-4z)/p_q(1),&1\le z<5/4.
\end{cases}
\tag{H.6.7}
\] Then \(g\) is continuous and positive in the open interval, \(g(1)=1\), and \(g\in H_0^1(0,5/4)\) with finite potential energy. Its free-side residual is \[
(-\partial_z^2-E)g
=\frac{16}{p_q(1)}(-p_q''-q p_q)(5-4z)>0.
\] Thus both pieces use the same physical energy \(E\).

Set \(D_a=g'(1^-)=a-\alpha-\beta/H(1)\) and \(F_a=-g'(1^+)=4p_q'(1)/p_q(1)\). The free derivative is \(F_a=4[1-q/(3-q/4)]\). Direct substitution gives \[
\begin{aligned}
D_2+\frac{11}{4}&=\frac{27}{1162}>0,
&F_2-\frac{11}{4}&=\frac{195}{1436}>0,\\
D_3+\frac{21}{10}&=\frac{2071}{307260}>0,
&F_3-\frac{21}{10}&=\frac1{206}>0.
\end{aligned}
\tag{H.6.8}
\] Consequently \(J:=g'(1^-)-g'(1^+)=D_a+F_a>0\). This derivative jump is intentional: \(g\) is a form-domain comparison function, not an element of the operator domain with matching derivatives.

\textbf{The whole-interval form inequality.} Write \(V(z)=a(a-1)z^{-2}\mathbf1_{(0,1)}(z)\) and, on each open piece, put \(\rho=[-g''+(V-E)g]/g\). The two residual identities give \(\rho>0\) almost everywhere. This function is bounded: the singular residual is \(O(z^6)\) at zero, while the free residual divided by \(g\) is \[
\frac{q^2x^2(4-qx^2)}{9(1-qx^2/12)^2},\qquad x=5-4z,
\] which extends to zero at the other endpoint. All one-sided interface values are finite.

For \(u\in C_c^\infty(0,5/4)\), integration by parts separately on the two pieces gives \[
\begin{aligned}
\int_0^{5/4}\bigl(|u'|^2+(V-E)|u|^2\bigr)\,dz
&=\int_0^{5/4}g^2\left|\left(\frac ug\right)'\right|^2 dz\\
&\quad+\int_0^{5/4}\rho|u|^2\,dz+J|u(1)|^2.
\end{aligned}
\tag{H.6.9}
\] The outer boundary terms vanish; the two interface terms add to \(J|u(1)|^2\) because \(g(1)=1\). Discarding the first nonnegative integral, form-core density and the Hardy inequality from \ref{def-high-operator} extend the resulting inequality to every \(u\in H_0^1(0,5/4)\). Indeed, the shifted form is continuous in the form norm, \(\rho\) is bounded, and evaluation at \(1\) is continuous in \(H^1\).

The right side after this omission is strictly positive for every nonzero \(u\), since \(\rho>0\) almost everywhere. The compact embedding established in \ref{def-high-operator} ensures that the first Rayleigh quotient is attained. Applying the inequality to a first eigenfunction therefore gives \(\lambda_1(K_{2;1,1/4})>25/2\) and \(\lambda_1(K_{3;1,1/4})>61/3\). Multiplication by \((5/4)^2\) and scaling monotonicity prove (H.6.1) on the entire interval, including \(\delta=0\), with the exact margins \[
\frac{625}{32}-\frac{39}{2}=\frac1{32}>0,
\qquad
\frac{1525}{48}-\frac{1587}{50}=\frac{37}{1200}>0.
\]

\end{proof}

\begin{lemma}[Nonexistence of large-height candidates]\label{lem-high}

There are no data satisfying \ref{def-candidate} with \(h\ge4\).

\end{lemma}

\begin{proof}[Proof of \ref{lem-high}]

Suppose, to the contrary, that data in \ref{def-candidate} exist with \(h\ge4\). Use the rational bounds in \ref{lem-pi-rational} and set \[
\eta=\frac{ph}{S},\qquad r=\frac{\rho h}{p},\qquad
\delta=\frac1h\in(0,1/4],\qquad L=1+h.
\] Since \(S>ph\), the candidate lower bound \(\rho\ge S/(2L^2)\) and the strict upper bound \(\rho<\beta/[h(h+2)]\) yield \[
\frac8{25}\le\frac{h^2}{2(h+1)^2}<r<\frac{h+1}{h+2}
=\frac{1+\delta}{1+2\delta}<1.
\tag{H.1.1}
\]

\textbf{Choice of weight.} Define \[
\sigma=\begin{cases}2,&r\le2/3,\\3,&r>2/3,\end{cases}
\qquad A=\frac{\sigma(\sigma-1)p^2}{h^2}>0,
\qquad\tau=\frac{A}{\pi^2\rho^2}.
\tag{H.1.2}
\] Put \(\mu_\tau=\mu(A)\) and \(E_\tau=2(\pi^2\rho+A/\rho)/S\). The candidate condition requires \(\mu_\tau\le E_\tau\), whereas \[
E_\tau h^2=2\eta\left(\pi^2r+\frac{\sigma(\sigma-1)}r\right).
\tag{H.1.3}
\]

\textbf{Area and the operator with its complete interface.} If \(p<q\), \ref{lem-high-area-cost} gives \(h>k\), a crossing ratio \(0<\alpha<1/2\), and \(\eta<D^{-1}\), where \(D=1+4\alpha(1-\alpha)+3\delta+\delta^2\). In this case set \((\ell,b)=(1-\alpha,\alpha+\delta)\). If \(p\ge q\), the same lemma gives \(\eta<(1+\delta)^{-2}\), and we set \((\ell,b)=(1,\delta)\).

Reverse and rescale the original interval by \(z=(L-y)/h\) and \(v(z)=\sqrt h\,u(L-hz)\). This unitary map sends \(H_0^1(0,L)\) bijectively onto \(H_0^1(0,1+\delta)\) and satisfies \[
\int_0^{1+\delta}|v'|^2=h^2\int_0^L|u'|^2,
\qquad
\int_0^{1+\delta}\frac{h^2A}{f(L-hz)^2}|v|^2
=h^2\int_0^L\frac{A}{f(y)^2}|u|^2.
\] The long tail dominates for \(0<z<1-\alpha\) in the crossing case and for \(0<z<1\) in the noncrossing case. Thus \(f(L-hz)=pz\). The comparison potential from \ref{def-high-operator} therefore satisfies \[
V_K(z)=\frac{\sigma(\sigma-1)}{z^2}\mathbf1_{(0,\ell)}(z)
\le\frac{h^2A}{f(L-hz)^2}.
\tag{H.3.1}
\] On the remaining interval we only lower a nonnegative potential to zero. The functions retain their common \(H_0^1\) domain, with no additional zero trace at the interface. The image of every function in the original form domain lies in the comparison form domain. Comparing their Rayleigh quotients and then taking infima yields \[
h^2\mu_\tau\ge\lambda_1(K_{\sigma;\ell,b}).
\tag{H.3.2}
\]

\textbf{Crossing case.} The strictly convex function \(r\mapsto2(ar+c/r)\) attains its maximum on a closed interval at an endpoint. Since \(\pi^2<10\), for \(8/25<r\le2/3\) we have \[
2(\pi^2r+2/r)<\max\{189/10,58/3\}<39/2;
\] Similarly, for \(2/3<r<1\), \[
2(\pi^2r+6/r)<\max\{94/3,32\}=32.
\] With \(C_2=39/2,C_3=32\), equation (H.1.3), the inequality \(\eta<D^{-1}\), and \ref{lem-crossing-sturm} give \[
E_\tau h^2<\frac{C_\sigma}{D}
<\lambda_1(K_{\sigma;\ell,b})\le h^2\mu_\tau,
\] contradicting the candidate condition.

\textbf{Noncrossing case.} The bound \(2(\pi^2r+2/r)<39/2\) for \(\sigma=2\) was established above. For \(\sigma=3\), use \(2/3<r<1\) and \(\pi^2<987/100\). The same endpoint comparison for the strictly convex function gives the uniform bound \[
2(\pi^2r+6/r)
<\max\left\{\frac{779}{25},\frac{1587}{50}\right\}
=\frac{1587}{50}.
\] Set \(B_2=39/2\) and \(B_3=1587/50\). The area inequality \(\eta<(1+\delta)^{-2}\) from \ref{lem-high-area-cost}, equation (H.1.3), and the direct spectral bounds in \ref{lem-no-crossing-sturm} give, for either order, \[
E_\tau h^2
<\frac{B_\sigma}{(1+\delta)^2}
<\frac{\Lambda_\sigma(\delta)}{(1+\delta)^2}
=\lambda_1(K_{\sigma;1,\delta})
\le h^2\mu_\tau.
\] This contradicts the candidate condition throughout \(0\le\delta\le1/4\). All possible data have now been excluded.

To recover the complete Robin response, note that \(E_\tau<\mu_\tau\) also places \(E_\tau\) below both first eigenvalues obtained by imposing an additional zero trace at \(y=1\): a zero-trace function on either side extends by zero to an admissible test function for the complete Rayleigh quotient. Thus all spectral hypotheses of \ref{lem-interface} hold, and the complete response sum is strictly positive.

\end{proof}


\section{Coupled responses and middle-height exclusion}

\begin{definition}[T.5 The true tail and the full response]\label{def-coupled-model}

Use \ref{def-analytic}. For \((h,a,b)\in\mathcal D_*\) and \(0<k\le h\), set \(L=1+h\) and define on \((0,L)\) \[
f_k(y)=\begin{cases}1-dy,&0<y<1,\\
\max\{p(1-x/h),q(1-x/k)_+\},&y=1+x,\ 0<x<h.
\end{cases}
\] Let \(A_0=A=\Pi\rho_+^2/2\) and \(E=3\Pi\rho_+/S_0\). Denote by \(\mu_k(A_0)\) the first eigenvalue of the full Friedrichs--Dirichlet realization with nonnegative potential \(V_k=A_0/f_k^2\). The functions have zero trace at both outer endpoints and a common value at the interface \(y=1\). The side Dirichlet--Dirichlet problems impose an additional zero trace at that interface. If \(E\) lies strictly below both side first eigenvalues, use the unique trace-one minimizing extensions \(g_H,g_T\) in \ref{lem-interface} and define \[
\mathfrak d_H(E)=g_H'(1),\qquad
\mathfrak d_T^{\rm coupled}(E)=-g_T'(1).
\]

\end{definition}

\begin{proof}[Proof of \ref{def-coupled-model}]

By \ref{def-analytic}, \(p>0\), \(q\ge0\), \(0<w=p+q<1\), and \(h>0\). The head width is at least \(w>0\). The tail width is at least \(p(1-x/h)>0\), with a positive lower bound on compact interior subintervals. If \(k<h\), the second branch vanishes when \(x\) is sufficiently close to \(h\); if \(k=h\), the width equals \(m(1-x/h)\). Thus the potential is bounded on either side near the interface, and near the right endpoint it is a constant multiple of \((L-y)^{-2}\). The one-dimensional right-end Hardy inequality \(\int u^2/(L-y)^2\le4\int|u'|^2\) holds on \(H_0^1(0,L)\): prove it on the smooth core by integration by parts and Cauchy--Schwarz, then extend by density. Consequently the potential form norm is equivalent to the \(H_0^1\) norm, the form is closed, and the core is dense. The representation theorem for closed nonnegative forms gives a unique Friedrichs operator. Compact embedding and the interval Poincare inequality give a positive first eigenvalue that is attained. The interface-zero side forms have the same properties. Responses are defined only below the prescribed side Dirichlet eigenvalues; existence, uniqueness, and derivative traces of their minimizing extensions follow from \ref{lem-interface}.

\end{proof}

\begin{lemma}[R.6.3 and T.5.2 A lower bound for the true common-interface response]\label{lem-coupled-response-lower-bound}

For all parameters in \ref{def-coupled-model}, both side first Dirichlet eigenvalues are strictly greater than \(E\), and \[
\mathfrak d_H(E)+\mathfrak d_T^{\rm coupled}(E)\ge\mathcal G(h,a,b).
\] If \((h,a,b)\in\mathcal D\), the right-hand side is strictly positive and \(\mu_k(A_0)>E\).

\end{lemma}

\begin{proof}[Proof of \ref{lem-coupled-response-lower-bound}]

Use \ref{def-coupled-model} and \ref{def-analytic}. By \ref{lem-pole-free}, the first Dirichlet eigenvalues of the head and model tail exceed \(E\), \(0<A_0<1\), \(0<K=E-A_0<\Pi\), \(1<\sigma<5/2\), \(T=h^2E\) lies below the scaled model-tail first eigenvalue, and \(r,C,D>0\). The model-tail potential is \(V_m(1+x)=A_0/[m^2(1-x/h)^2]\).

Since \(k\le h\), we have \(0\le(1-x/k)_+\le1-x/h\) for \(0<x<h\). Hence \(f_k(1+x)\le m(1-x/h)\), and the true tail potential is at least the model potential. Both zero-trace form domains are \(H_0^1(1,L)\), by the right-end Hardy inequality and local boundedness of the potentials. Rayleigh comparison makes the true tail first Dirichlet eigenvalue at least the model value, and therefore strictly greater than \(E\). The head bound also follows from \ref{lem-pole-free}, so the two spectral hypotheses of \ref{lem-interface} hold.

To identify the head ODE solution with the variational minimizer, note that its potential is bounded and smooth on \([0,1]\). The solution belongs to \(H^2(0,1)\), with endpoint traces zero and one. Integration by parts against smooth zero-trace tests gives orthogonality for the shifted form, which extends by \(H_0^1\) density to the full side zero-trace domain. Since \(E\) lies below the first side Dirichlet eigenvalue, the shifted form is coercive on the zero-trace subspace. Any other function with the same traces differs from the solution by a zero-trace function \(z\), and its energy difference is \(Q_E(z)>0\) unless \(z=0\). Thus the solution is exactly the unique trace-one minimizing extension of \ref{lem-interface}.

The head potential is \(A_0/(1-dy)^2\), and its normalized minimizing extension is the solution of \ref{lem-head-response}. The bounds \(0<A_0<1<2\) and \(0<K<\Pi\) check all parameter hypotheses. That lemma and \ref{lem-cot-rational} give \[
\mathfrak d_H(E)\ge\mathsf H(K)+\mathcal B(A_0,w).
\]

For the model tail, set \(z=(h-x)/h\). Its equation becomes \(-u''+\sigma(\sigma-1)z^{-2}u=Tu\), because \(A_0h^2/m^2=\Pi R^2/2=\sigma(\sigma-1)\). By \ref{lem-tail-recurrence}, its regular solution gives model response \(D_\sigma(T)/h\ge\mathsf T(\sigma,T)/h\). The factor \(1/h\) comes from \(-\partial_x u((h-x)/h)=u'(z)/h\) at the interface \(z=1\), followed by division by \(u(1)>0\) for trace-one normalization.

The hypotheses of \ref{lem-regular-tail} hold with \(s=\sigma>1\) and \(T=h^2E\) below the scaled first Dirichlet eigenvalue. That lemma gives positivity, finite endpoint form energy, and the regular asymptotics \[
u_s=z^s(1+O(z^2)),\qquad
u_s'=s z^{s-1}+O(z^{s+1}).
\] It also identifies \(u_s/u_s(1)\) as the unique trace-one minimizing extension on \((0,1)\). Under \(z=(h-x)/h\), the physical shifted form is \(1/h\) times the normalized shifted form at energy \(T=h^2E\). Hence \[
g_m(1+x)=\frac{u_s((h-x)/h)}{u_s(1)}
\] is exactly the model-tail minimizing extension of \ref{lem-interface}, with interface trace one, zero outer trace, and outward response \(D_s(T)/h\). The normalization and both response comparisons use the same physical interface and energy.

The common affine trace-one form domain must also be checked. Since \(p>0\) and \(f_k(y)\ge p(L-y)/h\) for \(1<y<L\), we have \(0\le V_k(y)\le A_0h^2/[p^2(L-y)^2]\). The model potential \(A_0h^2/[m^2(L-y)^2]\) satisfies the same bound. For \(z\in H_0^1(1,L)\), right-end Hardy gives \(\int z^2/(L-y)^2\le4\int|z'|^2\). Both zero-trace form domains are therefore exactly \(H_0^1(1,L)\), rather than merely being related by inclusion.

Choose one common smooth lifting \(\psi\) with value one at \(y=1\), supported in a right-hand neighborhood of the interface where the potentials are bounded, and vanishing near \(L\). It has finite energy for both potentials. Any \(u\in H^1(1,L)\) with \(u(1)=1,u(L)=0\) satisfies \(u-\psi\in H_0^1(1,L)\); Hardy's bound and the finite energy of the lifting make both weighted integrals finite. Conversely, elements of either trace-one form domain have these Sobolev traces. Thus the affine domains are both \(\psi+H_0^1(1,L)\). The pointwise potential comparison and minimization below take place on this same set.

For a fixed interface trace, the minimum shifted energy is nondecreasing under pointwise increase of the nonnegative potential: compare the energies function by function on the common affine domain and take infima. By \ref{lem-interface}, each minimum is the corresponding outward response. Thus the true coupled-tail response is no smaller than the model response. Adding the head inequality gives the claimed full response bound.

For completeness, compactness for the full interval follows explicitly as well. The true full potential is bounded on the head and near the interface, and the tail satisfies \(V_k\le A_0h^2/[p^2(L-y)^2]\). Endpoint Hardy therefore makes the full form domain \(H_0^1(0,L)\), with equivalent form and \(H_0^1\) norms, closed form, and dense smooth core. The complete model potential has the same properties. Nonnegativity makes either form norm control the full \(H^1\) norm; Rellich's compact embedding \(H_0^1(0,L)\hookrightarrow L^2(0,L)\) gives compactness for both full form domains. The interface-zero side forms are compactly embedded for the same reason. Thus every compactness hypothesis of the strict interface criterion is satisfied.

The true potential is bounded on both sides near the interface, and its form is closed and compactly embedded, so \ref{lem-interface} applies. If the point belongs to \(\mathcal D\), then \ref{lem-terminal} gives \(\mathcal G>0\). The common-interface response is strictly positive, and the strict spectral criterion in \ref{lem-interface} gives \(\mu_k(A_0)>E\).

\end{proof}

\begin{lemma}[T.5.2--T.5.3 Closure of the middle-height spectral contradiction]\label{lem-middle-spectral-closure}

There are no data satisfying \ref{def-candidate} with \(3/8<h\le4\). More explicitly, if such data existed, setting \(a=\beta/h,b=Lq/h\) would give a point in \(\mathcal D\) of \ref{def-analytic}, and \(\rho_+=\beta/[h(h+2)]\), \(A_0=\Pi\rho_+^2/2\), and \(\overline E=3\Pi\rho_+/S_0\) would satisfy the contradictory spectral bounds \[
\mu(A_0)>\overline E,\qquad
\mu(A_0)\le\frac2S\left(\Pi\rho+\frac{A_0}{\rho}\right)
<\frac{3\Pi\rho_+}{S}\le\overline E.
\]

\end{lemma}

\begin{proof}[Proof of \ref{lem-middle-spectral-closure}]

Use \ref{def-candidate} and \ref{def-analytic}. The first assertion of \ref{lem-robin} places the parameter image in \(\mathcal D\). The original candidate has \(0<k\le h\), and its piecewise maximum width is exactly \(f_k\) from \ref{def-coupled-model}. Hence the first eigenvalue of the corresponding full Friedrichs--Dirichlet operator equals \(\mu_k(A_0)\). By \ref{lem-coupled-response-lower-bound}, \(\mu(A_0)>\overline E\).

On the other hand, the candidate's longitudinal spectral bound gives \(\rho\ge S/(2L^2)\). Its inequalities \(\beta\le h\), \(S=1+\beta+\gamma>1+\beta\), and \(\gamma>0\) imply \[
\frac\rho{\rho_+}\ge\frac{Sh(h+2)}{2\beta L^2}
\ge\frac{(1+\beta)h(h+2)}{2\beta L^2}
\ge\frac{h+2}{2L}>\frac12.
\] The outer-ear bound gives \(\rho<\rho_+\). For \(x=\rho/\rho_+\in(1/2,1)\), \[
\frac32-x-\frac1{2x}=\frac{(1-x)(2x-1)}{2x}>0.
\] The all-weight necessary condition in \ref{def-candidate}, applied to this same \(A_0>0\), gives \(\mu(A_0)\le2(\Pi\rho+A_0/\rho)/S<3\Pi\rho_+/S\). Also \(q=\gamma/(1+k)\) and \(\gamma\le k\) imply \(q/(1-q)=\gamma/(1+k-\gamma)\le\gamma\), so \(S_0=\beta+1/(1-q)\le\beta+1+\gamma=S\). The strict upper bound is therefore at most \(3\Pi\rho_+/S_0=\overline E\), contradicting the strict lower bound already obtained.

\end{proof}


\section{Emptiness of the candidate set}

\begin{lemma}[The candidate parameter set is empty]\label{lem-no-candidate}

The set defined in \ref{def-candidate} is empty.

\end{lemma}

\begin{proof}[Proof of \ref{lem-no-candidate}]

Use all definitions, notation, and hypotheses in the exact statement of \ref{def-candidate}. If such data existed, \ref{lem-small} would give \(h>3/8\). If \(h\ge4\), this would contradict \ref{lem-high}. If \(3/8<h<4\), \ref{lem-middle-spectral-closure} gives contradictory strict upper and lower spectral bounds at the same energy. Equivalently, \ref{lem-robin} maps the data into \(\mathcal D\) of \ref{def-analytic} and gives \(\mathcal G<0\), whereas \ref{lem-terminal} gives \(\mathcal G>0\). This contradiction exhausts all heights.

\end{proof}


\section{Complete proof of the main theorem}

\begin{proof}[Proof of \ref{thm-main-nonconvex-pentagon}]

Use the admissible class and spectral functional of \ref{def-domain}. If \(P\) has three or four active sides, the classical Polya--Szego (1951) area-spectral bounds for triangles and simple quadrilaterals give \(F(P)\ge2\pi^2\). By \ref{lem-square}, there is then a same-area genuine pentagonal competitor with strictly smaller eigenvalue. A simple triangle is itself necessarily convex, but the numerical threshold comparison still applies.

Now suppose that there are five active sides. No angle is flat, and the interior-angle sum of a simple pentagon is \(3\pi\). A nonconvex pentagon therefore has one or two strictly reflex angles. In the two-angle case, \ref{lem-double} gives a strict competitor. In the one-angle case, \ref{lem-single} either gives a strict competitor or places the polygon in the signed residual family.

Before using the corner variation formula, choose the first eigenfunction and verify its admissibility. The domain \(P\) is bounded and connected, with a simple polygonal Lipschitz boundary. Its nonnegative closed Dirichlet form on \(H_0^1(P)\) embeds compactly into \(L^2\), so the Rayleigh infimum is attained by an \(L^2\)-normalized function. Taking its absolute value makes it nonnegative. The weak eigenvalue equation, interior regularity, and the strong maximum principle then give a strictly positive first eigenfunction \(U\) on connected \(P\), with \(\int_PU^2=1\). Positivity and the standard simplicity of the first eigenvalue provide the required ground-state properties. These classical existence and positivity results can be found in Henrot (2006), Chapter 1, Sections 1.2.2 and 1.3.3. Boundedness, connectedness, and Lipschitz boundary verify the zero-trace form and maximum-principle hypotheses.

In the residual family, the differentiable vertex variation formula of \ref{lem-hadamard-corners} shows that, unless \(P\) is area-critical, there is an area-tangent vertex velocity \(V\) with \(D\lambda[V]<0\), reversing its sign if necessary. Realize these vertex velocities as follows. Place the closure of \(P\) inside a large rectangle. Choose a generic rotated direction in which all original vertices have distinct horizontal coordinates. Join the uniquely leftmost and rightmost vertices along outward horizontal rays to the corresponding rectangle sides. These two disjoint connecting segments lie outside \(P\) and split the exterior annular region into two simple polygons. Triangulate \(P\) and those two exterior polygons by finite nondegenerate triangulations, with matching subdivisions along common edges. Assign the chosen velocities \(V_i\) to the original vertices. At additional mesh points on an original boundary edge, interpolate its two endpoint velocities linearly; give all other mesh points and rectangle-boundary points velocity zero. Affine interpolation on each triangle, extended by zero outside the rectangle, produces a globally continuous Lipschitz field \(V\) that is affine along every entire original edge.

Let \(F_t=I+tV\). If \(|t|\operatorname{Lip}(V)<1\), then \(|F_t(x)-F_t(y)|\ge(1-|t|\operatorname{Lip}(V))|x-y|\). For any target point \(y\), the contraction \(x\mapsto y-tV(x)\) yields a unique preimage. Thus \(F_t\) is a global bi-Lipschitz homeomorphism. The image boundary is simple, and each original edge maps to one segment, creating no extra breakpoints. Every original edge has positive length and every adjacent-edge determinant is nonzero. These finitely many quantities depend continuously on \(t\), so they remain nonzero for sufficiently small \(|t|\). Hence the five active sides and nonflat angles persist. This construction both validates the vertex displacement and supplies the piecewise affine field required by the corner Hadamard formula. Write \(P_t=F_t(P)\). Its spectral and area derivatives are exactly those of the prescribed vertex velocities. Normalize to \(Q_t=(|P|/|P_t|)^{1/2}P_t\). Area is then preserved exactly. Since \(D|P|[V]=0\), we have \(\frac d{dt}\lambda_1(Q_t)|_{t=0}=D\lambda[V]<0\). A sufficiently small positive \(t\) therefore gives a strict competitor.

If \(P\) is area-critical, suppose for contradiction that \(F(P)\le2\pi^2\). By \ref{lem-geometry-input} it produces data satisfying \ref{def-candidate}, contradicting \ref{lem-no-candidate}. Thus \(F(P)>2\pi^2\), and \ref{lem-square} again supplies a strict competitor.

Finally, \ref{lem-five-edge-approximation} handles the requirement of exactly five active sides. If a strict competitor \(Q\) has fewer than five, introduce a nonflat breakpoint of amplitude \(\varepsilon\) at the midpoint of an open edge; if necessary, repeat on another open edge. Each sufficiently small perturbation preserves simplicity and adds exactly one active side. Piecewise affine interpolation on a fixed triangulation gives bi-Lipschitz maps converging to the identity, whose derivatives and inverse derivatives also converge to the identity. Two-sided Rayleigh comparison then gives continuity of area and of the first eigenvalue. This remains true after area normalization, so the original strict spectral gap persists. We may therefore always choose a competitor with exactly five active sides. Such a competitor contradicts global minimality of any nonconvex domain, proving the final assertion.

\end{proof}


\section{Degeneration, existence, and the conditional convex bridge}

Write \(\mathscr P_{\le n}(A)\) for the simple polygons of area \(A\) with at most \(n\) sides, and let \(\overline{\mathscr P}_{\le n}(A)\) be their polygonal closure in the complementary Hausdorff topology. This closure allows vertices touching other edges, zero-length edges, repeated edges, slits, and splitting. We use the standard external existence result that, for fixed \(n\ge3\) and \(A>0\), the first Dirichlet eigenvalue attains a minimum in this closure, at a bounded, possibly degenerate polygonal object with at most \(n\) segments. The result concerns general polygons: convexity is not assumed, and the minimizing object is not asserted to be a Jordan domain.

See Bogosel--Bucur \cite{nc:BB}, Section 1, pp.~20--21, and Section 7, p.~86 of the published version. Equation (1.2) formulates the problem for all simple polygons of fixed area \(\pi\), followed by the existence statement for a possibly degenerate closure minimizer. Other areas follow by scaling. The paper attributes existence to Henrot \cite{nc:henrot2006}, Chapter 3, particularly Section 3.3. We use this classical existence result rather than reprove its concentration-compactness or boundedness construction. Bogosel--Bucur's Theorem 7.1 bounds the diameter of an optimal object; it is not used as a uniform diameter bound for an arbitrary decreasing sequence with low eigenvalues.

\begin{definition}[A bounded five-segment limit chain and its domain]\label{def-polygonal-limit}

Consider a sequence of counterclockwise simple polygons of area \(A>0\), with at most five sides, lying in a common bounded container. Add repeated marked vertices so that each polygon has five marked vertices. Along a vertex-convergent subsequence, let the trace of the limiting oriented closed chain be \(\Gamma=[v_0v_1]\cup\cdots\cup[v_4v_0]\). For \(x\notin\Gamma\), let \(\operatorname{ind}_\Gamma(x)\) be its winding number. Define \(\Omega=\{x\notin\Gamma:\operatorname{ind}_\Gamma(x)=1\}\). This is a possibly disconnected open set of area \(A\). Denote the selected subsequence by \((P_j)\). The precise convergence assertion is that \(\mathbb R^2\setminus P_j\) converges in Hausdorff distance to \(\mathbb R^2\setminus\Omega\). Equivalently, for any closed rectangle \(B\) containing all polygonal closures in its interior, the compact sets \(B\setminus P_j\) converge to \(B\setminus\Omega\). For any nonempty bounded open set \(\Omega\), this appendix defines \(\lambda_1(\Omega)\) by the Dirichlet Rayleigh infimum over \(H_0^1(\Omega)\), without requiring connectedness.

\end{definition}

\begin{proof}[Proof of \ref{def-polygonal-limit}]

A common vertex-convergent subsequence exists by finite-dimensional compactness. After adding marked vertices, the tuples \((v_{0,j},\ldots,v_{4,j})\) lie in the fifth Cartesian power of one bounded set. As a bounded sequence in \(\mathbb R^{10}\), they have a subsequence along which all five vertices converge simultaneously, by Bolzano--Weierstrass. Relabel this one subsequence by \(j\); no incompatible vertexwise subsequences are used.

Parametrize the five edges linearly on five fixed parameter intervals, with repeated vertices giving zero segments. Let \(\gamma_j,\gamma\) be the approximating boundary curves and the limiting oriented chain, and set \(\varepsilon_j=\max_i|v_{i,j}-v_i|\to0\). Convex combinations on each edge give \(\sup_t|\gamma_j(t)-\gamma(t)|\le\varepsilon_j\). If \(x\notin\Gamma\) and \(\operatorname{dist}(x,\Gamma)>\varepsilon_j\), the straight homotopy between the two parametrized curves avoids \(x\), so their winding numbers agree. A counterclockwise simple approximating boundary has winding number one inside and zero outside. Thus the limiting winding number takes only the values zero and one. It is locally constant off \(\Gamma\), so \(\Omega\) is open, possibly disconnected, and remains in the common bounded container.

An explicit distance estimate proves complementary convergence. Choose a closed rectangle \(B\) whose interior contains the common container, and put \(F_j=B\setminus P_j\), \(F=B\setminus\Omega\). Both are nonempty closed sets, with \(\Gamma\subset F\) and \(\Gamma_j\subset F_j\). If \(x\in F_j\) and \(\operatorname{dist}(x,\Gamma)>\varepsilon_j\), the homotopy gives winding number zero, hence \(x\in F\). Otherwise \(\operatorname{dist}(x,F)\le\operatorname{dist}(x,\Gamma)\le\varepsilon_j\). Therefore \(\sup_{F_j}\operatorname{dist}(\cdot,F)\le\varepsilon_j\). Conversely, if \(x\in F\) has distance greater than \(\varepsilon_j\) from \(\Gamma\), its zero winding number gives \(x\in F_j\). Otherwise choose \(q\in\Gamma\) at distance at most \(\varepsilon_j\) from \(x\), and let \(q_j\in\Gamma_j\) have the same parameter. Its distance from \(q\) is at most \(\varepsilon_j\), so \(\operatorname{dist}(x,F_j)\le2\varepsilon_j\). Thus \(d_H(F_j,F)\le2\varepsilon_j\to0\). Outside the rectangle all points already belong to both complements, so this also proves Hausdorff convergence of the full complements, without an additional compactness assumption.

The finite union of segments \(\Gamma\) has area zero. Away from it, the indicator functions eventually equal the limiting winding number. The common bounded container supplies an integrable majorant, giving \(|\Omega|=\lim_j|P_j|=A>0\). To exhibit a nonzero test function, choose \(x\in\Omega\). Since \(\Omega\) is open, there is \(r>0\) with \(B_{2r}(x)\subset\Omega\). Define \(\varphi(y)=\exp[-1/(r^2-|y-x|^2)]\) for \(|y-x|<r\), and zero elsewhere. Every derivative inside the ball is the same exponential times a polynomial in \(y-x\) and \((r^2-|y-x|^2)^{-1}\). The limit \(t^{-N}e^{-1/t}\to0\) as \(t\downarrow0\), for each integer \(N\ge0\), shows that every derivative tends to zero at the sphere and matches the exterior zero function smoothly. Thus \(\varphi\) is a nonzero element of \(C_c^\infty(\Omega)\), supported in \(\overline B_r(x)\subset\Omega\), with positive mass and finite Dirichlet energy.

The Rayleigh test class is therefore nonempty, and the Dirichlet spectrum of this bounded open set is well defined through the standard closed form. No claim is made that slits have zero spectral effect, and connectedness of \(\Omega\) is not assumed.

\end{proof}

\begin{lemma}[A.1--A.2 Winding representation and area continuity]\label{lem-limit-winding-area}

With the construction of \ref{def-polygonal-limit}, \(\operatorname{ind}_\Gamma(x)\in\{0,1\}\), \(\Omega=\{x\notin\Gamma:\operatorname{ind}_\Gamma(x)=1\}\), and \(|\Omega|=\lim_j|P_j|\).

\end{lemma}

\begin{proof}[Proof of \ref{lem-limit-winding-area}]

Use \ref{def-polygonal-limit} and its construction proof. The five-segment parametrizations converge uniformly. For each \(x\notin\Gamma\), the straight homotopy between the approximating boundary and the limit chain therefore avoids \(x\) for sufficiently large \(j\). The Jordan theorem and homotopy invariance of winding number give \[
\operatorname{ind}_\Gamma(x)\in\{0,1\},\qquad
\Omega=\{x\notin\Gamma:\operatorname{ind}_\Gamma(x)=1\}.
\tag{A.1}
\] The complementary distance estimate \(d_H(B\setminus P_j,B\setminus\Omega)\le2\max_i|v_{i,j}-v_i|\) retains every one-dimensional slit: all points of \(\Gamma\) belong to the complement. Off \(\Gamma\), the indicator \(\mathbf1_{P_j}\) eventually equals \(\mathbf1_\Omega\). The finite-segment trace has area zero, and the common bounded container gives integrable domination. Hence \[
|\Omega|=\lim_j|P_j|=A.\tag{A.2}
\] This area identity makes no assertion about slit capacity; the spectral argument below separately uses domain inclusion and spectral continuity.

\end{proof}

\begin{lemma}[The five-segment degeneration lemma]\label{lem-five-segment-degeneration}

Let \(\Omega\) be the positive-area limit domain of \ref{def-polygonal-limit}. If, after deleting zero edges, the closed chain is not a Jordan polygonal boundary, then there are finitely many simple polygons \(Q_1,\ldots,Q_r\) with pairwise disjoint interiors, each with at most four active sides, such that \[
\Omega\subset U:=\bigcup_{j=1}^r Q_j,
\qquad |\Omega|=\sum_{j=1}^r|Q_j|=|U|.
\tag{A.3}
\] Each \(Q_j\) denotes its open interior. Passing from \(\Omega\) to \(U\) fills at most finitely many one-dimensional slits.

\end{lemma}

\begin{proof}[Proof of \ref{lem-five-segment-degeneration}]

Use \ref{def-polygonal-limit} and \ref{lem-limit-winding-area}.

First delete only zero edges. Temporarily retain redundant marked vertices on edges with the same direction; merge consecutive collinear edges only after the resulting loop is known to be Jordan. Deleting zero edges does not change the winding number or domain, and every remaining segment is still an original limiting edge. A transverse crossing between the relative interiors of two edges would persist under small vertex perturbations, contradicting simplicity of the approximating boundaries. Hence a remaining non-Jordan closed chain must have either a repeated vertex or a vertex lying on a nonincident edge. Collinear overlap also supplies such an endpoint-edge contact, and adjacent backtracking falls into the same class. Pure vertex coincidence may create extra zero edges in the splitting below, which only lowers the counts.

Consider first a five-segment chain. Relabel cyclically so that the contact is \(q=v_0\in[v_j,v_{j+1}]\), with \(j=1,2,3\). Insert \(q\) on the touched edge and split the oriented chain there into two closed subchains, with edge counts \[
j+1\quad\text{and}\quad5-j.
\] The only possibilities are \(2+4\), \(3+3\), and \(4+2\). If \(q\) is an endpoint of the touched edge, deleting the resulting zero segment decreases the count.

\textbf{Case 1: \(2+4\) or \(4+2\).} The two-segment loop traverses one segment in both directions and has winding number zero. Deleting it leaves the winding number unchanged off its trace and can only fill a one-dimensional slit. The remaining chain has at most four segments and no transverse crossings. Each open component of its complement contains a point avoiding the deleted backtracking segment. At that point its winding number equals that of the original chain and belongs to \(\{0,1\}\). Since winding number is constant on each component, the remaining chain has only winding numbers zero and one everywhere off its trace. If it is Jordan, it bounds a simple polygon with at most four sides. Otherwise split again at a contact. A four-segment chain splits into at most \(2+3\) or \(3+2\); a nondegenerate three-segment chain is a triangle, while a degenerate one has area zero. Continuing to remove zero-winding backtracking leaves a positive-area part that is a simple polygon with at most four sides. It contains the original domain, and the difference lies in finitely many segment traces.

\textbf{Case 2: \(3+3\).} Now \(q\in[v_2,v_3]\), and the two subchains are \[
(q,v_1,v_2,q),\qquad(q,v_3,v_4,q).
\] If either triangle loop is degenerate, it has winding number zero and leaves at most one triangle. Suppose both are nondegenerate and \(q\) is not an endpoint of the cut edge. Let their open triangles be \(T_1,T_2\), and let \(\ell\) be the line of that edge. The intersections of their closed triangles with \(\ell\) are respectively \[
[q,v_2],\qquad[q,v_3].
\] These segments lie on opposite rays from \(q\). Each triangle therefore has an interior point outside the other closed triangle: choose a relative-interior point of its base, which has positive distance from the other closed triangle, and move it slightly into its own interior. In particular, neither triangle contains the other.

Let the orientations of the two triangle loops be \(\sigma_1,\sigma_2\in\{-1,1\}\). Near the exclusive interior points just constructed, the winding numbers of the full chain are respectively \(\sigma_1,\sigma_2\). By (A.1), both must equal \(+1\). If \(T_1\cap T_2\) were nonempty, their intersection would be open and the winding number there would be two, again contradicting (A.1). Hence the triangle interiors are disjoint. The winding-one region of the full chain is their union, possibly with part of the original one-dimensional trace removed.

This rules out subtraction by negatively oriented loops, nested loops producing holes, and erroneous area counting along shared slits, without assuming that every contact is transverse.

If the initial chain already has fewer than five segments, the four- and three-segment arguments above apply. Take the interiors of the remaining positive-area simple loops as \(Q_j\). Formula (A.1) gives the inclusion, and the zero area of the finite-segment trace gives the area identity.

\end{proof}

\begin{lemma}[A.4 The spectral threshold for degenerate domains]\label{lem-degenerate-threshold}

In \ref{def-polygonal-limit}, if the chain after deletion of zero edges is not a Jordan boundary, or if \(\Omega\) is a simple polygon with at most four sides, then \(|\Omega|\lambda_1(\Omega)\ge2\pi^2\).

\end{lemma}

\begin{proof}[Proof of \ref{lem-degenerate-threshold}]

Use the positive-area limit open set of \ref{def-polygonal-limit}, the Rayleigh definition in \ref{def-domain}, and the inclusion and area identity in \ref{lem-five-segment-degeneration}.

Every simple polygon \(Q_j\) with at most four sides satisfies the classical triangle and quadrilateral inequality \[
|Q_j|\lambda_1(Q_j)\ge2\pi^2.
\] Triangles satisfy the stronger bound \(4\pi^2/\sqrt3\), which is unnecessary here.

The spectral identity for the finite union follows from the Rayleigh definition even when boundaries touch. The \(Q_j\) are open polygons with disjoint interiors. If a point of \(\partial Q_i\) were inside another \(Q_j\), nearby interior points of \(Q_i\) would also belong to \(Q_j\), a contradiction. Thus no boundary point of any \(Q_j\) lies in \(U\). For \(v\in C_c^\infty(U)\), its compact support has positive distance from the complement of \(U\), so \(v|_{Q_j}\in C_c^\infty(Q_j)\). Density and continuity of restriction extend this to \(H_0^1(U)\), with each restriction in \(H_0^1(Q_j)\). Conversely, approximation by smooth cores shows that zero extensions of elements of \(H_0^1(Q_j)\) belong to \(H_0^1(U)\); so do their finite sums. Hence \(H_0^1(U)\) is the orthogonal direct sum of these spaces, with mass and Dirichlet energy adding componentwise. Each component energy is at least \(\lambda_1(Q_j)\) times its mass, so every full Rayleigh quotient is at least \(\min_j\lambda_1(Q_j)\). For the reverse inequality, choose an index \(j_0\) attaining this finite minimum. For every \(\varepsilon>0\), choose a nonzero \(v_\varepsilon\in H_0^1(Q_{j_0})\) with quotient below \(\lambda_1(Q_{j_0})+\varepsilon\), and extend it by zero. The preceding core argument puts it in \(H_0^1(U)\) with unchanged mass and energy. Thus \(\lambda_1(U)\le\lambda_1(Q_{j_0})+\varepsilon\); let \(\varepsilon\downarrow0\). This argument does not require the Rayleigh infimum to be attained. Boundary contact neither adds an interface condition nor removes a zero trace.

The precise classical input is Freitas--Siudeja (2010), DOI 10.1051/cocv/2009018: equation (1.1) for general simple quadrilaterals and equation (2.3) for triangles. Convexity is not required.

Since \(\Omega\subset U\), zero extension gives \(H_0^1(\Omega)\subset H_0^1(U)\) and therefore \[
\lambda_1(\Omega)\ge\lambda_1(U)
=\min_j\lambda_1(Q_j).
\] Choosing a minimizing index \(j_0\) and using the total area identity yields \[
|\Omega|\lambda_1(\Omega)
\ge\left(\sum_j|Q_j|\right)\lambda_1(Q_{j_0})
\ge|Q_{j_0}|\lambda_1(Q_{j_0})
\ge2\pi^2.
\tag{A.4}
\] A slit may have positive capacity and substantially raise the eigenvalue. The argument only fills slits to enlarge the domain and lower its eigenvalue, in the correct direction. If the limiting chain is already a simple polygon with at most four sides, the same classical inequality gives (A.4) directly.

\end{proof}

\begin{lemma}[Spectral limits of bounded sequences of simple polygons]\label{lem-polygonal-spectral-continuity}

For the fixed-area counterclockwise sequence of simple polygons with at most five sides in a common bounded container, and its limit open set \(\Omega\) from \ref{def-polygonal-limit}, the chosen vertex subsequence satisfies \(\lambda_1(P_j)\to\lambda_1(\Omega)\). The limit may be disconnected or contain one-dimensional slits.

\end{lemma}

\begin{proof}[Proof of \ref{lem-polygonal-spectral-continuity}]

Use the bounded sequence and complementary Hausdorff convergence in \ref{def-polygonal-limit}. By \ref{lem-limit-winding-area}, \(|\Omega|=A>0\), so the limit open set is nonempty.

The external input is Sverak's planar domain spectral-continuity theorem: for open sets in a fixed compact set \(B\subset\mathbb R^2\), complementary Hausdorff convergence to an open set \(\Omega\), together with a common finite bound on the number of connected components of the full-plane complements, implies gamma-convergence and convergence of every fixed-index Dirichlet eigenvalue. The precise reference is Henrot (2006), Theorem 2.3.19, p.~31, with the complementary-distance convention of Definition 2.3.14, p.~30. The original result is V. Sverak, On optimal shape design, Journal de Mathematiques Pures et Appliquees 72 (1993), 537--551.

The sequence is planar, and its construction supplies a common container. Choose a fixed closed rectangle \(B\) containing all closures. Each \(P_j\) is the open interior of a Jordan polygon. The Jordan curve theorem makes \(\mathbb R^2\setminus P_j\) the connected closed exterior including its boundary, so the number of complement components is always one. The construction in \ref{def-polygonal-limit} gives \(B\setminus P_j\to B\setminus\Omega\), exactly the required distance, rather than ordinary Hausdorff convergence of polygonal closures. Neither uniform corner Lipschitz constants nor connectedness of the limit are required. All hypotheses hold, so the result follows by taking index one. Slit area zero is never used to infer spectral invariance.

\end{proof}

\begin{theorem}[Existence of a simple pentagonal minimizer]\label{thm-simple-minimizer-exists}

For every \(A>0\), the first Dirichlet eigenvalue attains its minimum in the fixed-area class of simple polygons with at most five sides defined in \ref{def-domain}. Every minimizer in this class is a genuine simple polygon with exactly five active sides.

\end{theorem}

\begin{proof}[Proof of \ref{thm-simple-minimizer-exists}]

Use \ref{def-domain}, \ref{def-polygonal-limit}, \ref{lem-degenerate-threshold}, and \ref{lem-square}. The external existence input is Bogosel--Bucur (2024), https://doi.org/10.5802/jep.250, Section 1, pp.~20--21, and Section 7, p.~86: a minimizer exists in the complementary Hausdorff closure of fixed-area simple polygons with \(n\) sides. Degeneracy is allowed and convexity is not assumed.

The truncated-square lemma gives a simple convex pentagon \(Q\) such that \[
|Q|\lambda_1(Q)<2\pi^2.
\] Scaling to area \(A\) shows that the closure minimum is strictly below \(2\pi^2/A\). To connect the external existence input with the bounded-limit convention of this appendix, use the boundedness construction in the proof of Henrot (2006), Theorem 3.3.1, pp.~47--48, rather than applying Bogosel--Bucur's optimal-diameter theorem to arbitrary decreasing sequences. That classical proof trims finitely many thin spikes from a fixed-area polygonal minimizing sequence, producing a minimizing sequence \(\widetilde P_j\) with at most \(N\) sides, uniformly bounded diameters, \(|\widetilde P_j|\le A\), and \(\lambda_1(\widetilde P_j)\le\lambda_1(P_j)+o(1)\). We use only this construction from the published existence proof; no diameter bound is asserted for every original decreasing sequence. The source is the \href{https://ndl.ethernet.edu.et/bitstreams/902d1f87-86bf-4004-bcc9-4091b345cb09/download}{public book text}, Section 3.3.1, pp.~47--48. Its theorem and proof allow degenerate limits and do not assume convexity.

Area normalization preserves boundedness. The eigenvalues above have a common upper bound \(C\). The classical Faber--Krahn inequality for nonempty bounded open sets, in Section 3.2 of the same book, gives \(|\widetilde P_j|\lambda_1(\widetilde P_j)\ge c_{\rm FK}>0\), where \(c_{\rm FK}\) is the first eigenvalue of the unit-area disk; its numerical value is not needed. Hence \(|\widetilde P_j|\ge c_{\rm FK}/C>0\). Set \(t_j=(A/|\widetilde P_j|)^{1/2}\); these factors are uniformly bounded. Change of variables in dimension two gives \(\lambda_1(t_j\widetilde P_j)=t_j^{-2}\lambda_1(\widetilde P_j)\), so the normalized domains remain a minimizing sequence with uniformly bounded diameters. Translate them into one closed ball.

Take a limit of this bounded fixed-area minimizing sequence. Add repeated marked vertices so that every polygon has five marked vertices, and choose a convergent subsequence in the fifth Cartesian power of the common closed ball. By \ref{def-polygonal-limit} and \ref{lem-limit-winding-area}, this gives a chain \(\Gamma\) and open set \(\Omega\) with \(|\Omega|=A>0\). All hypotheses of \ref{lem-polygonal-spectral-continuity} hold for this sequence, so \(\lambda_1(\Omega)=\lim_j\lambda_1(P_j)=\inf_{P\in\mathscr P_{\le5}(A)}\lambda_1(P)\). Putting \(P_*=\Omega\) explicitly proves that the selected vertex limit attains the closure minimum; uniqueness of a Hausdorff limit alone would not establish this. This limit satisfies every hypothesis of the degeneration threshold in this appendix. If its chain is not a Jordan boundary, apply that threshold. If the chain is Jordan but has fewer than five active sides after merging consecutive collinear edges of the same direction, apply the same threshold for simple triangles or quadrilaterals. No boundedness assumption is imposed on arbitrary approximating sequences.

The closure minimizer \(P_*\) just constructed can be neither a non-Jordan degenerate domain nor a simple triangle or quadrilateral, since all such objects satisfy (A.4). It is therefore a genuine simple pentagon: every edge has positive length, no angle is flat, and there is no self-contact, backtracking, or slit. Thus the minimum is attained in the original admissible class. The same threshold excludes fewer than five active sides for any minimizer in that class.

This step does not require a claim that finitely many Steiner deformations turn every nonconvex domain into a convex one.

\end{proof}

\begin{corollary}[A conditional bridge with the separate convex theorem as an explicit premise]\label{cor-conditional-convex-bridge}

Use \ref{def-domain}. Let \(R\) be a regular pentagon and \(\Lambda=|R|\lambda_1(R)\). If every convex polygon with at most five sides satisfies \(F(P)\ge\Lambda\), with equality only for Euclidean similarity images of a regular pentagon, then the same bound and equality classification hold for every simple polygon with at most five sides. The class with exactly five active sides has the same minimum, and every nonconvex domain satisfies a strict inequality.

\end{corollary}

\begin{proof}[Proof of \ref{cor-conditional-convex-bridge}]

Use \ref{def-domain}. The convex theorem is only the explicit hypothesis of this corollary. By \ref{thm-simple-minimizer-exists}, the full class of area-\(A\) simple polygons with at most five sides has a minimizer \(P_*\). If it were nonconvex, \ref{thm-main-nonconvex-pentagon} would supply a same-area competitor with a strictly smaller eigenvalue, contradicting minimality. Hence \(P_*\) is convex, and the assumed convex inequality gives \(A\lambda_1(P_*)\ge\Lambda\).

Take \(t=(A/|R|)^{1/2}\). A two-dimensional change of variables leaves Dirichlet energy unchanged, multiplies mass by \(t^2\), and gives a bijection between the corresponding \(H_0^1\) spaces. Thus \[
|tR|=A,\qquad A\lambda_1(tR)=|R|\lambda_1(R)=\Lambda.
\] This regular pentagon belongs to the admissible class, so conversely \(A\lambda_1(P_*)\le\Lambda\). Minimality now gives, for every simple polygon \(P\) with at most five sides, \[
F(P)\ge\Lambda.
\] If equality holds, \(P\) is itself a minimizer in the full class. Nonconvex exclusion makes it convex, and the assumed convex equality classification makes it a Euclidean similarity image of a regular pentagon. Conversely, rigid motions and the displayed scaling law ensure equality for those images. Every nonconvex domain therefore satisfies \(F(P)>\Lambda\).

The exactly-five-side class contains the regular pentagon and is contained in the at-most-five-side class. Its minimum and equality set are consequently the same.

\end{proof}


\clearpage
\begin{thebibliography}{10}

\bibitem{nc:BB}
Beniamin Bogosel and Dorin Bucur.
\newblock On the polygonal faber--krahn inequality.
\newblock {\em Journal de l'École polytechnique — Mathématiques},
  11:19--105, 2024.
\newblock Section 1, pages 20--21; Section 7, page 86.

\bibitem{nc:brock}
Friedemann Brock.
\newblock Continuous steiner-symmetrization.
\newblock {\em Mathematische Nachrichten}, 172:25--48, 1995.

\bibitem{nc:brock2000}
Friedemann Brock.
\newblock Continuous rearrangement and symmetry of solutions of elliptic
  problems.
\newblock {\em Proceedings of the Indian Academy of Sciences: Mathematical
  Sciences}, 110:157--204, 2000.
\newblock Theorem 3.2 and Corollary 3.2.

\bibitem{nc:Capriani}
Giuseppe~Maria Capriani.
\newblock The steiner rearrangement in any codimension.
\newblock {\em Calculus of Variations and Partial Differential Equations},
  49:517--548, 2014.
\newblock Theorem 2.3 and conditions (2.12)--(2.14).

\bibitem{nc:CF}
Andrea Cianchi and Nicola Fusco.
\newblock Steiner symmetric extremals in pólya--szegö-type inequalities.
\newblock {\em Advances in Mathematics}, 203:673--728, 2006.
\newblock Theorem 2.2.

\bibitem{FS}
Pedro Freitas and Bartłomiej Siudeja.
\newblock Bounds for the first dirichlet eigenvalue of triangles and
  quadrilaterals.
\newblock {\em ESAIM: Control, Optimisation and Calculus of Variations},
  16:648--676, 2010.

\bibitem{nc:grisvard}
Pierre Grisvard.
\newblock {\em Elliptic Problems in Nonsmooth Domains}.
\newblock Pitman, 1985.

\bibitem{nc:henrot2006}
Antoine Henrot.
\newblock {\em Extremum Problems for Eigenvalues of Elliptic Operators}.
\newblock Birkhäuser, 2006.
\newblock Chapter 3, especially Section 3.3.

\bibitem{nc:henrot}
Antoine Henrot and Michel Pierre.
\newblock {\em Variation et optimisation de formes}.
\newblock Springer, 2005.

\bibitem{nc:kato}
Tosio Kato.
\newblock {\em Perturbation Theory for Linear Operators}.
\newblock Springer, corrected second edition, 1980.
\newblock Chapter VII, Section 4; Chapter VI, Section 2.

\bibitem{PS}
George Pólya and Gábor Szegő.
\newblock {\em Isoperimetric Inequalities in Mathematical Physics}, volume~27
  of {\em Annals of Mathematics Studies}.
\newblock Princeton University Press, 1951.

\bibitem{nc:sverak}
Vladimír Šverák.
\newblock On optimal shape design.
\newblock {\em Journal de Mathématiques Pures et Appliquées}, 72:537--551,
  1993.

\end{thebibliography}

\end{document}
